\documentclass[../main.tex]{subfiles}
\begin{document}
%==========================================TIÊU ĐỀ TẬP========================================
\begin{table}[h]
  \begin{adjustwidth}{-1cm}{}
    \begin{tabular}{>{\centering\arraybackslash}p{6cm}|>{\raggedright\arraybackslash}p{12.5cm}}
      \multirow{5}{6cm}
      % Cột bên trái: ÉPISODE và số tập
      {\\[-40pt]
        \flaregothic\fontsize{20pt}{20pt}\selectfont \centering
        {\contour{errorfive}{\textcolor{black}{ÉPISODE}}}\color{black}\\
        \burbank\fontsize{72pt}{72pt}\selectfont
        \includegraphics[height=3.12359cm]{../../../../Image/Book?/number/num32.png}
        \\[18pt]
        \flaregothic\fontsize{18pt}{18pt}\selectfont
        {\contour{errorfive}{\textcolor{white}{CHOIX}}}
        \color{black}
      }
       &
      % Dòng thứ nhất: tên tập tiếng Nhật
      {\fotskip\fontsize{18pt}{18pt}\selectfont \ruby{愛}{あい}の\ruby{治}{ち}\ruby{療}{りょう}}
      \\[6pt]
      % Dòng thứ hai: tên tập tiếng Anh
       & {\cafeteria\fontsize{18pt}{18pt}\selectfont Remedy for Love}                                                       \\[4pt]
      % Dòng thứ ba: tên tập tiếng Việt
       & {\vnmsans\fontsize{18pt}{18pt}\selectfont Liều thuốc tình yêu}                                                     \\[8pt]
      % Dòng thứ tư: nhân vật xuất hiện
       & {\caviar{\fontsize{14pt}{14pt}\selectfont Nhân vật: Tất cả mọi người xuất hiện từ trước tới nay, trừ \emp{Choco}}}
      % \\[6pt]
      % Dòng cuối cùng: Thời gian diễn ra của tập (thường sẽ là ngày phát hành)
      % & {\continuum\fontsize{16pt}{16pt}\selectfont Ngày phát hành: 11/06/2022}\\[8pt]
    \end{tabular}
  \end{adjustwidth}
\end{table}
%===========================================NỘI DUNG==========================================
\color{black}

\pov{Hikawa Kuon}

Nếu phải tóm gọn lại những gì đã xảy ra tối hôm qua, thì quả thật, đó là một câu chuyện mà chúng tôi vẫn chưa hoàn toàn có thể hiểu được. Những gì mà ba anh em chúng tôi, Shino, Sakura, và những cộng sự đi cùng tới lúc này tựa không khác gì một thước phim tua chậm, chỉ để rồi mọi thứ dần sáng tỏ, nhưng những sự mơ hồ và những lời giải đáp bỏ ngỏ vẫn cứ thế vất vưởng trên đầu chúng tôi.

Lúc đó, cả nhóm lại thắp đèn rọi vào khu rừng, một lần nữa. Gió lạnh ban đêm len lỏi qua tán lá, mang theo mùi hơi ẩm của đất và điều gì đó nặng trĩu hơn. mùi của quá khứ chưa kịp phai. Chúng tôi bước đi, tiếng cành khô gãy dưới chân như những mảnh ký ức vỡ vụn. Trước mặt, tấm màn sương mỏng dần dần hé mở, để lộ ra câu chuyện mà dân làng từng cố quên lãng: những đứa trẻ bị dẫn vào rừng, không vì tội lỗi, không vì được chọn lựa, chỉ vì người lớn cần một lời bào chữa cho mùa màng thất bát, cho dịch bệnh, cho nỗi sợ hãi không tên. Một lời giải thích đơn giản: ``Trời trách, nên hiến con.'' Cái giá quá rẻ, quá quen, đến mức máu của trẻ thơ đã thấm sâu vào rễ cây, thành dòng chảy ngầm nuôi nấng linh hồn oán hận suốt ngàn năm.

Tôi vẫn thấy cô bé ấy đứng giữa bạch đàn cổ thụ. Đôi mắt xanh biếc không ánh lệ, chỉ phản chiếu lại khoảnh khắc cuối, khi bóng tối sà xuống, khi bàn tay người lớn buông lỏng, khi tiếng khóc nghẹn giữa cổ họng. Cô không van xin, chỉ nhìn, như thể muốn khắc vào trời đất hình hài của mình, để đời sau không bao giờ tưởng rằng nỗi đau này chỉ là lời nguyền xa xôi. Chính ánh nhìn ấy đã níu giữ linh hồn cô, biến khu rừng thành nấm mồ chung của cả làng, và biến chúng tôi - những kẻ đến sau - thành nhân chứng bất đắc dĩ.

Hiểu thì đã hiểu, nhưng biết phải làm gì tiếp theo vẫn là ẩn số. Tôi thì thầm, chủ yếu để trấn an chính mình: ``Lần này, ta sẽ không lặp lại sai lầm.'' Giọng tôi nhỏ đến mức gần như bị gió cuốn đi. ``Không lặp lại,'' nghe thật dễ, nhưng ``đúng'' ở đây là gì? Đưa linh hồn các em về với ánh sáng? Hay để khu rừng tiếp tục ôm mối hận, nhắc nhở đời sau đừng bao giờ dùng máu trẻ thơ để đánh đổi yên bình?

Tôi chưa có lời đáp, chỉ biết bước tiếp, để tiếng nức nở của rừng già theo chân, như hồi chuông ngân dài trong đêm, kêu gọi một phương án mà chúng tôi chưa dám thốt ra.

Có lẽ lúc này\dots\ chúng tôi vẫn nên tiếp tục sống với cuộc sống thường nhật trước mắt. Tiếp tục với bài vở, tiếp tục với ca làm việc chiều, lẳng lặng theo dõi mọi thứ xung quanh. Biết đâu đấy, sẽ có một câu trả lời nào đó mà chúng tôi chưa tìm thấy, ở một thời điểm bất thình lình nào đó.

\il

\begin{center}
  \uline{Sáng hôm sau.}
\end{center}

Buổi sáng ngay sau sự kiện hôm đó. Một ngày khoảng giữa tháng Năm.

Thời tiết bắt đầu ấm dần lên, cái lạnh của đầu xuân gần như đã biến mất, nhường chỗ cho một thứ không khí nhẹ hơn, nhưng cũng báo hiệu mùa hè đang đến gần. Các học sinh của trường Aogami cũng đã bắt chuyển dần sang đồng phục hè để đối phó tiết trời đang dần nóng lên.

Tôi đi cạnh Kaigou-chan, Suzuki thì thấp thoáng phía trước, cả ba cùng hòa vào dòng học sinh đang tíu tít áo sơ-mi ngắn tay. Chúng tôi vừa đẩy cánh cửa lớp 1-C, vừa nghĩ tới bài kiểm tra Văn chiều nay thì tất cả đồng loạt dừng bước.

Không khí như ngưng đọng. Mùi bột mì nướng ấm ấm, thoang thoảng mùi trứng tanh nhẹ, lẫn chút hương vani hóa học từ máy pha nước nóng góc hành lang len lỏi vào mũi.

Tôi chớp mắt. Hôm qua đây vẫn còn bảng đen, bàn ghế thẳng tắp, còn bây giờ, trước mắt chúng tôi\dots

``Ờm, thế quái nào\dots?'' Kaigou-chan lẩm bẩm, tay vẫn nắm chặt cặp.

Tôi không đáp, chỉ đưa mắt quét một vòng. Khung cảnh lớp học vẫn như mọi ngày, nhưng toàn bộ bàn học đã xếp chồng lên nhau, kê sát vách và để ở cuối lớp. Ở trung tâm, bốn cái bàn dài được ghép thành hình chữ nhật, phủ khăn kẻ carô đỏ trắng. Trên mặt khăn, thớt gỗ, âu inox, rổ trứng, túi bột mì vẫn còn nguyên tem, chai dầu ăn, lọ muối hồng, và chiếc máy đánh trứng cầm tay màu hồng nhạt đang nằm chễm chệ như chờ việc.

Kaigou-chan khẽ kéo tay tôi, chỉ lên bảng lớn. Dưới dòng chữ ``Ngày lễ làm sô-cô-la nhân ngày Lễ tình nhânnhân'' viết phấn trắng, còn treo sẵn mấy chiếc tạp dề xanh lá có thêu tên lớp. Một tờ giấy A4 hồng dán ngay cửa ra vào: ``Học sinh vui lòng rửa tay trước khi vào bếp. Đeo tạp dề vào để tránh bị bẩn đồng phục.''

Tôi nuốt khan. Cảm giác như vừa bước nhầm cửa vào lớp học nấu ăn của khách sạn, chứ không phải nơi thường ngày chúng tôi ngáp dài chờ hết giờ Vật lý.

``Chào buổi sáng, ba người nhà Hikawa\~{}!''

Tiếng reo vang lên, kéo chúng tôi khỏi cơn ngỡ ngàng. Ayame-san đứng giữa lớp, tay chống hông, nụ cười toe toét như thể cô vừa thắng một ván game khó nhằn. Bên cạnh, Wata-san loay hoay thắt dây tạp dề, mặt ửng hồng vì gấp.

Kaigou là người đầu tiên lấy lại giọng: ``Hai cậu đang biến lớp học thành gì thế này?''

``Lễ tình nhân chứ còn gì!'' Ayame nhún vai, đôi mắt cô sáng lên. ``Hôm nay cả trường sẽ làm sô-cô-la đấy.''

Tôi và Suzuki đồng thanh, nét mặt không giấu nổi sự hoang mang: ``Tháng Năm mà cũng Valentine sao?''

Nàng tóc vàng gật đầu, giọng nhỏ nhẹ như đang dỗ trẻ con: ``Ừ, lễ tình nhân của chúng mình.''

Không khí như bị bơm thêm hơi lạ, lạnh buốt đến mức mùi bột mì nướng cũng phải đóng băng giữa không trung. Kaigou-chan khẽ cau mày, giọng khẽ run: ``Tôi tưởng Valentine thì phải 14-2, còn White Day 14-3 chứ nhỉ?'' Câu nói vừa dứt, cả lớp bỗng chốc im lặng như thời gian ngưng trôi, mọi ánh mắt đồng loạt ngoái lại nhìn chúng tôi, không chỉ ngạc nhiên, mà còn như đang quan sát những sinh vật ngoài hành tinh vừa rơi xuống từ lỗ đen vũ trụ.

Shino-san chớp mắt, đôi con ngươi xanh không cảm xúc như tấm gương bể: ``14-2 là ngày gì? Cậu đang nói gì vậy, Kaigou-san?'' Giọng cô trầm đến mức tiếng phấn rơi trên bảng cũng có thể nghe thấy.

Uzuki-san nghiêng đầu, dấu chấm hỏi hiện rõ: ``White Day\dots\ nghe giống như một ngày tôn vinh màu sắc nhỉ? Có lẽ là lễ hội bột giấy?'' Cô mỉm cười như thể vừa nghĩ ra một câu đố hóc búa.

Ayame-san cũng nhướng mày, cặp mắt đỏ lóe sáng: ``Lạ nhỉ, bọn này còn chưa từng nghe. Nó là gì vậy? Có ăn được không?'' Cô khẽ liếm môi, như thể vừa ngửi thấy mùi sô-cô-la lạ lẫm từ chiều không gian khác.

Tôi quay sang Kaigou-chan, rồi sang Suzuki. Ba đứa nhìn nhau, rõ ràng là ngỡ ngàng với những câu hỏi khác lạ như thế này, giống như vừa phát hiện ra mình đang đứng giữa một bản giao hưởng mà bản nhạc bị xáo trộn hoàn toàn.

``Khoan đã\dots'' rôi xoa trán, cảm giác như có làn sương mù mỏng đang bò qua ký ức, xóa nhòa những ngày lễ quen thuộc, ``\dots Đừng nói là ở đây không có Valentine tháng Hai thật đấy nhé.''

``Không có đâu.'' Một giọng nói chen vào, nhẹ nhàng nhưng đủ sức cắt đứt mọi âm thanh.

Touka-san bước ra từ phía sau, tay cầm một túi nguyên liệu, tà áo trắng phất phơ như cánh cổ tích. ``Ít nhất là không phải theo cách mà cậu đang nghĩ,'' cô nàng Thần đồng mỉm cười, nụ cười ẩn chứa cả dải ngân hà, một nụ cười của kẻ biết rõ chúng tôi vừa lạc vào cánh cửa không có tên.

\dots À phải rồi. Đây là một vũ trụ khác, nơi mà những ngày lễ không tính bằng lẽ thông thường, nơi mà trái tim không cần phải đợi đến 14-2 mới được phép đập nhanh hơn một nhịp. Cũng không ngạc nhiên nếu nơi này có những sự khác biệt so với nơi tôi sống, giống như việc bánh mì có thể mọc cánh, hay như việc tình yêu ở đây được nướng bằng lò than thay vì nồi áp suất thời gian.

Cô bạn tóc xanh lam đặt túi nguyên liệu xuống bàn, tiếng nhựa kêu khẽ.  Cô nhìn chúng tôi, mắt xanh như mặt hồ sương sớm, rồi từ tốn giải thích: ``Ở Aogami, Valentine không đến vào tháng Hai. Nó đến vào tháng Năm, khi hoa phượng chưa nở nhưng gió đã nghe mùi hè.''

Kaigou-chan chau mày, giọng hạ thấp: ``Sao lại dời lễ hội vậy? Hẳn phải có lý do sâu xa nào đó đằng sau chứ, đúng không?''

Touka-san khẽ nhún vai, mái tóc xanh rung như nhánh liễu. ``Sau đợt `tái sinh', hội hè ở đây đều được sắp lại như bộ bài mới. Người ta nói rằng, nếu đã chết một lần, thì sinh nhật cũng nên đổi.  Ngày Lễ tình thân vì thế được chọn để mang một nghĩa khác.''

Suzuki khoanh tay, lắng nghe. ``Ý cậu là sao?''

Cô bạn Thần đồng quét mắt qua lớp học, nơi những nữ sinh đang diện những tạp dề còn thơm mùi vải mới, những khuôn chocolate đã xếp thành hàng như quân cờ. Cô nói chậm rãi, như thể từng chữ đều có trọng lượng: ``Tình cảm nảy sinh sau tái sinh.''

Tôi nghe tiếng mình thở. ``Cậu nghe điều đó ở đâu vậy?''

``Saki-san nói với mình như vậy. Cô ấy bảo: khi thị trấn còn chưa kịp khóc, thì hãy để những trái tim mới được phép đập sai nhịp. Hoặc ít nhất, đó là những gì cậu ấy nói với mọi người như thế. Mình cũng không thực sự tin lắm đâu.''

Im lặng. Ngoài hành lang, tiếng chuông báo giờ học vang lên, nhưng trong lớp không ai nhúc nhích.

Azuki-san cười nhẹ, mắt nhìn thẳng vào không trung. ``Nghe cũng\dots\ hợp lý. Sau cơn mưa máu, người ta vẫn trồng hoa. Sau đêm con nít khóc, người lớn vẫn dám yêu.''

Shino-san gật, giọng khô như lá rụng: ``Bắt đầu lại\dots\ không phải là câu cửa miệng. Đps là cách họ chọn để không chết lần hai.''

``Mọi người văn vẻ quá nhể?'' Kaigou-chan nhìn hai cô bạn, đổ mồ hôi hột.

Tôi đứng giữa họ, cảm thấy mình như viên sỏi trôi vào bể kính. Mùi bột cacao phả vào mặt, ngọt đắng. Tiếng cười của Ayame, tiếng lách cách đánh trứng, tất cả đều thật, đều sống, đều\dots\ không thuộc về những điều mà tôi đã từng biế.biết.

Ngày Lễ tình nhân ở nơi tôi từng sống là ngày 14-2, có hồng, có socola hộp vuông, có những màn hẹn gặp viết sẵn trong thiệp. Đó là ngày mà cả thế giới thở cùng nhịp, trao cho nhau những món quà mà người ta thầm yêu mến.

Còn ở đây, nó là một cái tên được gọi lại, một mùa được gieo trễ, một lời thề thầm sau lưng người vừa đi qua cửa tử. Nó giống như một ngày tượng trưng cho sự tái sinh, một thứ tình yêu nảy nở sau sự đổ sập tối tăm hơn là một ngày lễ thông thường.

Tôi khẽ cười, không vui, không buồn. Chỉ thấy lòng nặng trĩu như khối bột đang nhào. ``Tình cảm nảy sinh sau tái sinh\dots\ à.''

Nếu thật vậy, thì đêm qua chúng tôi đã thấy máu của trẻ con thấm vào rễ cây, nay lại thấy đường kính của trái tim mới đo bằng khuôn thanh sô-cô-la. Và nếu có thứ gì đang mọc lên ở thị trấn này, thì không chỉ là nụ hôn đầu, mà còn là cơ hội\dots\ để sửa lỗi ngàn năm qua.

\ct

Không khí trong lớp 1-C nhanh chóng trở nên nhộn nhịp một cách khó tin.
Tiếng kéo ghế, tiếng bát đĩa chạm nhau và cả những tiếng cười nói đan xen khiến căn phòng chẳng khác gì một xưởng làm bánh thu nhỏ. Mùi sô-cô-la nóng chảy lan dần trong không khí, ngọt nhưng không gắt, đủ để khiến người ta vô thức thấy dễ chịu hơn.

Một vài người đã bắt đầu đeo tạp dề từ sớm, trong khi số khác vẫn còn loay hoay với việc phân biệt đâu là nguyên liệu chính. Có những người đã từng có kinh nghiệm làm sô-cô-la từ trước đó và biết được những bước cơ bản, nhưng cũng có những người chưa từng động dao bao giờ, chỉ đứng đó chết trân mà không biết làm gì đầu tiên. Không ai ngồi ngoài cuộc, chỉ khác nhau ở chỗ: người đã nhập cuộc, kẻ còn đang\dots\ nhập thân phận.

Với chúng tôi, tôi đã từng dấn thân vào bếp núc không ít lần khi bố mẹ hoặc khi Ryuuhou bận việc. Với chị em nhà Hikawa, họ còn quen thuộc hơn vào nơi đầy rẫy công cụ làm bếp và đồ ăn mua ở chợ. Nhưng, nếu nói về sô-cô-la hay làm bánh kẹo nói chung, thì ba chúng tôi chẳng khác nào những con gà mờ.

Wata-san đập tay hai cái, vừa đủ át tiếng ồn. ``Mọi người ơi, nghe này! Làm sô-cô-la không khó, nhưng mà làm ẩu thì ăn vào ôm bụng đấy nha\~{}!''

Cô nàng quét mái tóc vàng lềnh bềnh, nụ cười vừa tỏa sáng vừa có gì đó\dots\ đe dọa. Tôi biết cô ấy là người có kinh nghiệm làm bánh kẹo hơn tất cả mọi người trong lớp, nên có lẽ cô ấy có thể tin tưởng được.

Touka-san nghiêng nghiêng muôi gỗ, tiếp lời phó hội với tông giọng bình tĩnh hơn hẳn: ``Những bước cũng cơ bản thôi: cắt nhỏ, đun nhỏ cách thủy, khuấy đều tay. Nếu không muốn sô-cô-la thành\dots\ sô-cô-cục.''

Cô cầm muỗng khuấy thử một lượt, động tác gọn gàng, làm mẫu cho các bạn cùng lớp, miệng nhẹ nhàng đếm nhịp: ``Một, hai\dots\ một, hai\dots'' Cả lớp im phăng phắc, mắt dán vào cánh tay uyển chuyển của cô như thể đang xem một màn ảo thuật. Một vài người bắt chước ngay, giơ muỗng lên không trung rồi hạ xuống chậm rãi; những người khác thì cúi sát vào âu, môi mím chặt, cố gắng sao cho đường khuấy tròn trịa không văn một giọt. Tiếng kim loại va chạm nhẹ, tiếng hít thở khẽ và mùi cacao nóng hòa vào nhau tạo thành một bản giao hưởng nhỏ ấm áp, khiến không khí căng thẳng ban nãy bỗng tan biến, nhường chỗ cho niềm háo hức khó giấu.

Khi cô nàng Thần đồng dừng tay, khóe môi cô khẽ nhếch, đôi mắt xanh lướt qua từng khuôn mặt như thể thầm hỏi: ``Các cậu đã sẵn sàng chưa?'' Cả lớp đồng loạt gật đầu, tiếng ``rồi!'' vang lên rộn rã, lập tức hòa vào tiếng ghế kéo, tiếng mở túi bột và tiếng cười khúc khích - một âm thanh duy nhất của sự khởi đầu.

``Sau đó đổ vào khuôn, để nguội, và--'' Touka-san tiếp tục, để rồi Inari-san nói nốt, giọng khẽ tới mức nghe như gió lùa: ``\dots g-gói lại thật đẹp để tặng cho người quan trọng của mình. C-Có thể là một người bạn, một người mà các cậu yêu thầm, h-hoặc là\dots\ H-Híc!''

Cô nàng khẽ lùi về phía sau một chút, mặt đỏ bừng lên, luống cuống: ``M-Mọi người cứ làm từ từ thôi. N-Nếu có gì không hiểu thì cứ đừng ngần ngại hỏi mình!''

Nhìn cô bạn bờm Cáo trắng đỏ mặt vì ngượng ngùng, tôi khẽ huých Kaigou-chan, ghé sát tai nói thầm: ``\vie{\hl{Cậu ấy trông căng thật. Dù là lớp trưởng lớp mình thế kia\dots}}''

``\vie{\hl{Trách làm sao được. Cậu ấy có vẻ không phải người thích đối diện với đám đông như Hanabi đâu,}}'' cô nàng nhỏ giọng đáp lại, nhìn cô bạn Cáo trắng vẫn đang run rẩy kia, trước khi tiếng của cô bạn thuở nhỏ của mình vang lên những tiếng động viên.

Suzuki gãi cằm, cười khẽ: ``\vie{\hl{Vẫn giống như Inari-san mà ta biết nhỉ. Cậu ấy cũng cố gắng lắm đấy chứ.}}''

Ayame vung tay vòng qua đầu, ra hiệu như chỉ huy trưởng: ``Mọi người đã nắm bắt được chưa? Nếu rồi thì\dots\ xắn tay áo lên nào! \textbf{Bắt đầu vào việc!}''

Ngoài cửa sổ, những cánh hoa anh đào mùa xuân đang dần phai tàn chỉ còn cành trơ, như đang chờ một mùa hè mang mùi cacao thay cho mùi hoa. Giữa tiếng kim loại va chạm, tiếng cười khúc khích, tiếng nói cười giúp nhau, lớp 1-C bỗng trở thành một bếp lửa khổng lổ, nơi mỗi người đều đang nung nấu một điều chưa dám thốt thành lời.

Lớp nhanh chóng chia thành từng nhóm. Những ai đã từng học nấu ăn tự động tụ lại thành một nhóm quanh Touka-san, mắt vẫn dán vào muôi gỗ như thể đó là cây đũa thần. Nhóm bên cạnh - toàn những người vừa mới lần đầu cầm khuôn - thì tíu tít hỏi Wata-san cách phân biệt cacao đắng và đường nâu, giọng hồn nhiên đến mức khiến cả góc lớp bật cười. Một vài bạn lẻ tẻ đứng xa, mắt lấp lánh nhưng tay vẫn lúng túng, chỉ dám chạm vào túi bột mì khi thấy Ayame-san gật đầu khích lệ. Có đứa còn len lén nhét chiếc khuôn hình trái tim vào ngăn bàn, như giấu một lá thư chưa kịp ký tên. Trên bảng, phấn trắng vương vãi thành những vệt mờ, ghi dấu những công thức vội vã xếp thành hàng như dãy mật mã của tuổi mới lớn.

\ct

Góc cửa sổ bên phải lớp 1-C trở thành đấu trường của những người nhiệt huyết nhất. Hanabi vung muỗng như đang điều khiển đội quân, mắt sáng quắc: ``Ai khuấy nhanh nhất sẽ được ăn miếng đầu tiên!''

Akane không chịu thua, tay trái chống hông, tay phải quay như máy, bột cacao bắn tung tóe lên tóc. ``Tôi sẽ không thua đâu!''

Kaoru đứng giữa, vừa đếm nhịp vừa lùi dần, môi lẩm bẩm: ``Từ từ, chậm thôi, chậm thôi! Nếu khuấy nhanh quá thì--'' Một giọt bột đậu lên mũi cậu, cả bọn phá lên cười, thành ra tiếng ``bùm'' nhỏ vang khắp góc lớp như pháo thăng hoa.

``\dots sẽ bị như thế đấy,'' cô bạn tóc ngắn lẩm bẩm, mặt nhăn lại như khỉ vì bị châm chọc.

Nhóm ba người họ tiếp tục với những tiếng nói cười râm ran, nhưng dường như họ đang bắt đầu nghịch với đồ ăn hơn là nghiêm túc làm việc.

``Khoan đã, tôi nên cho đường nâu hay đường trắng đây?'' Hanabi mở nắp hộp, mắt tròn xoe.

``Hình như Touka-chan nói là\dots\ đường nâu cho mùi khói,'' Akane cầm thìa quấy điên cuồng, bột cacao bắn tung tóe lên tóc.

Kaoru ngồi trên thành ghế, một tay cầm hộp sữa tươi, một tay chỉ trỏ: ``Đổ thêm sữa vào nữa. Đổ đi, đổ đi!''

Cô nàng tóc nâu nghe theo, nhưng thay vì đổ sữa vào hỗn hợp bột, cô ấy lại\dots\ rót vào miệng mình. ``Ủa, ngon quá, có vani!'' Cô giấu vội hộp sau lưng khi cô nàng bờm Chó quay sang.

``Á à! Hanabi-chan ăn vụng rồi kìa!'' Akane rốt cuộc cũng phát hiện, nhưng tay vẫn không ngừng quay. Nồi cách thủy bốc hơi nghi ngút, sô-cô-la trong âu sệt lại thành khối đặc quánh.

``Chắc phải nóng hơn nữa,'' cô nàng tóc ngắt nhấp thêm ngụm sữa, môi ướt át, ``lửa to lên chút để hỗn hợp nở ra, rồi ta từ từ hạ lửa xuống.''

Akane gạt nút bếp sang phải cùng. Ngọn lửa ``rốp'' một cái, bốc phừng lên dữ dội, khói xám xịt bốc lên. Khối sô-cô-la vốn đã cạn nước giờ kết tinh thành một lớp đen bóng, cứng ngắc, như mảng đá bazan. Cô thử cạy bằng thìa, và  *CỐP*, thìa gãy.

``Nó cứng quá, không đổ khuôn được!''

``Thì\dots\ đập ra làm kẹo bỏ túi,'' Hanabi gợi ý, vẫn núp sau ghế tránh đòn.
Kaoru gãi đầu: ``Không sao, thêm sữa vào lại\dots\ à quên, sữa tôi uống hết rồi.'' Cô lắc vỏ hộp, nghe tiếng ``bình bịch'' rỗng tuếch.

Ba người nhìn nhau, cùng cười hềnh hệch. Hanabi lấy khuôn hình ngôi sao, đặt lên bàn, gõ gõ khối đá sô-cô-la cứng ngắc: ``Đập ra, mỗi người một mảnh, tặng bạn gọi là `vụn tình yêu'.''

``Tôi xí khối to để tặng cho Yuka nha\~{}'' nàng Khuyển nói. Kaoru đứng dậy, phất tay: ``Nói vậy thôi, chứ ta phải làm mẻ khác rồi. Để tôi chạy hỏi xem các thành viên câu lạc bộ Nấu ăn còn nguyên liệu không. Chứ Valentine mà không có sô-cô-la là hỏng đấy.''

``Đúng vậy nhỉ|\~{}''

Tôi đứng xa, chứng kiến cảnh tượng ấy, bất giác nhớ đến ``đám hề hololive'' hồi đó - những ý tưởng điên rồ mà các em ấy đưa ra, hay những lần phá tung cả văn phòng khiến tôi và A-san dọn dẹp mệt bở hơi tai. Họ cũng như thế: loay hoay, đổ vỡ, rồi lại cười tươi rói giữa đống hỗn độn.

``\textit{\dots Đúng là mình không thể nào thoát khỏi họ được mà,}'' tôi thì thầm. Không phải vì công thức, mà vì họ đang tự viết lấy công thức của riêng mình, một thứ sô-cô-la không cần hoàn hảo, chỉ cần đủ điên rồ, đủ chân thật.

\ct

Ở phía đối diện của lớp, nhóm ``nổi loạn'' cũng không hề kém cạnh.

Shion trước tiên dựng nồi lên bếp ga mini, không vội đổ sô-cô-la mà vẽ một vòng tròn mờ bằng ngón tay quanh thành nồi, miệng lẩm nhẩm: ``Ta sẽ truyền năng lượng bóng tối vào đây! \textbf{Hãy nổi lên nào, sinh vật huyền bí của ta!}''

Aku đang mở túi bột cacao nghe vậy giật mình, khuôn inox lập tức rơi ``choang'' xuống sàn. Cô lúi húi nhặt, lẩm bẩm về phía cô bạn ảo tưởng của mình: ``S-Shion-chan, đừng làm mình sợ chứ\dots''

``Không cần phải để ý con điên kia đâu,'' Koishin đỡ lời, giọng mềm. Nàng tsundere vẫn chừng mực, nhưng đôi tay lại cực nhanh như một tay nghề lão luyện: đong đo cẩn thận, trộn bột thành nhuyễn, đổ vào từng khuôn, đặt lên nồi hấp cách thủy. Động tác vừa dứt, cô lập tức lùi một bước, như sợ ai hiểu lầm mình đang quan tâm.

Shion bỗng reo lên: ``Xong phép rồi!'' rồi hạ lửa. Nồi sô-cô-la vốn nâu bóng bỗng sôi ùng ục, bề mặt nổi những bong bóng đen như mực. Mùi vẫn thơm, nhưng nhìn có vẻ\dots\ rất ma quái.

Nàng Hầu gái trố mắt: ``\textbf{N-N-Nó sống dậy kìa!!}'' Cô lùi lại đụng phải khay khuôn, ``choang'' lần hai. Khay vừa chạm đất, cả khối sô-cô-la đang sủi bọt trào ra ngoài, chảy lênh láng trên mặt bàn. Aku cuống quýt dùng thìa gạt, vô tình bê cả khăn lau chùi, và một lần nữa, cô trượt chân, ngồi bệt xuống sàn.

Tôi đứng xa, không kìm được lẩm bẩm: ``\textit{Aku-san vẫn cứ là Aku-san thôi nhỉ? Cô ấy thậm chí còn vụng hơn cả Hành Tím nhà mình nữa\dots}''

Koishin thở dài, cầm khăn ướt lau vội cho Aku. Xong, cô quay sang Shion, mặt vẫn điềm tĩnh: ``Shion-san à, đừng có bình phẩm mấy câu vô dụng nữa được không? Cậu làm bọn tôi khó chịu đấy.''

``Nhưng mà--''

``Không nhưng nhị gì nhé.'' Cô nàng tsundere lặng lẽ đặt lên bếp một nồi khác. Lửa vừa, sô-cô-la tan chậm, mùi thơm nguyên chất lan ra. Cô rót vào khuôn trái tim nhỏ, mặt không cảm xúc, không ngừng lầm bầm như sợ ai đó đang nhìn: ``T-tôi không làm cho ai cả đâu, chỉ là\dots\ tôi không muốn lãng phí nguyên liệu thôi.''

Shion nháy mắt với Aku, Aku cười trừ. Koishin làm xong, lập tức giấu riêng một hộp vào ngăn bàn, khóa kỹ. Cả hai kia đồng thanh thở phào: ``Ồ\~{}''

Koishin đỏ tai, quay đi: ``Nhìn cái gì? Làm đi, còn nhiều khuôn trống đấy!''

Tôi khẽ cười, xoay người về phía bàn mình. Mùi cacao nồng đã bắt đầu phủ kín lớp, nhưng giữa những âm thanh leng keng, tôi vẫn nghe thấy tiếng Shion lẩm nhẩm đọc thần chú. Có lẽ cô nàng lại sắp ``nâng cấp'' lửa lần nữa.

Bất chợt, khói bốc lên nghi ngút. Cả nhóm vừa hét vừa chạy tán loạn, từ những công đoạn cuối. Nàng Phù thủy giật mình thốt lên, giọng the thé: ``Hể? Sao nó lại cháy rồi!?''

Bạn thân của cô cau mày, lập tức đổ lỗi: ``Cậu nấu kiểu gì mà như đốt lò vậy!?''

``Đâu, ta mới chỉ điều chỉnh hỏa phép của nó lên thôi mà!'' Shion vừa nói vừa chỉ vào cái nồi đang nhảy khói.

Koishin cúi xuống nhìn, mặt bình thản nhưng giọng đầy nghi ngờ: ``\dots Lên là bao nhiêu độ?''

Shion gãi má, cười trừ: ``\dots Chắc là tầm lửa to nhất có thể?''

Cô nàng tsundere thở dài, lắc đầu: ``Thế này thì toang rồi.''

Nàng tóc tím-hồng cũng ngán ngẩm chẳng kém: ``Hết cứu luôn rồi đó.''

Nhìn vẻ mặt bất mãn của hai người, nàng tóc trắng-xám trố mắt: ``Ơ? T-Ta đã làm sai ở công đoạn nào vậy!?''

\ct

Còn tôi đứng ở một góc bàn cùng Kaigou, Suzuki và Sakura. Nhóm quan sát.

``Chúng ta nên bắt đầu từ đâu đây?'' Suzuki hỏi, mắt vẫn dán vào đống nguyên liệu chưa mở hết tem.

``Làm đúng quy trình trước đã,'' Sakura đáp, tay đã lật tờ công thức in nhỏ như một bản nhạc giao hưởng. ``Cắt nhỏ bơ cacao, đường nâu, bột cacao đắng, thêm chút vani và muối hồng.''

Cô cầm dao, cắt từng viên bơ cacao thành hạt lựu đều tăm tắp, rồi xếp thành hàng như quân cờ trên thớt. Bột cacao được rây mịn qua rây inox, lớp bụi nâu li ti bốc lên, phủ một mùi đắng ngọt lên đầu ngón tay.

``\dots Chuẩn đến từng gam luôn kìa,'' tôi lẩm bẩm, nhìn cô bạn thao tác như đang chơi một trò chơi xếp hình.

``Tôi chỉ làm theo lý thuyết thôi,'' Sakura đáp, khẽ đẩy gọng kính. ``Còn thực hành\dots\ cậu biết đấy, lý thuyết suông khô khan lắm. Phải tự tay làm mới vỡ ra được.'' Quả không hổ danh người đã từng làm trong ngành sinh vật học, nơi mà thời gian thực nghiệm nhiều tới mức còn lấn át cả thời gian ngủ.

Kaigou cười khẽ, gãi má: ``Vậy thì để tôi căn chỉnh bếp!'' Cô bật bếp mini, đặt nồi cách thủy lên, nước trong nồi bắt đầu rít lên như tiếng mèo con. Suzuki lùi lại một bước, tay giơ cao: ``Mình sẽ phụ khuấy!''

Tôi nhìn ba người họ, bỗng thấy mình như khán giả đứng ngoài rạp. Tôi cầm chiếc khuôn hình chiếc lá, nhẹ nhàng đặt lên khăn carô. ``Tôi sẽ làm một viên sô-cô-la hình lá, để khi tan ra, người ăn sẽ thấy mùa thu ở lưỡi.''

``Sến quá đấy,'' Kaigou nhăn mũi, nhưng khóe môi lại cong lên. ``Nhưng mà\dots\ cũng được, tôi sẽ làm hình trăng khuyết.''

Suzuki gãi đầu: ``Vậy\dots\ em làm hình quả trứng được không? Dễ đổ khuôn.''

``Quả trứng thì nghe như đang làm bánh,'' Sakura cau mày, nhưng rồi lại gật. ``Ừ, cũng được. Miễn là cậu không đập vỡ.''

Chúng tôi bắt đầu. Bơ cacao tan chậm trong nồi, mùi béo ngậy bốc lên, quyện với mùi vani và đường nâu. Kaigou khuấy đều tay, mắt dán vào nồi canh chừng từng phút một. Suzuki thì loay hoay lau khuôn, miệng lẩm nhẩm: ``Không bong, không vỡ, không dính\dots''

Tôi đổ hỗn hợp vào khuôn lá, mặt nâu bóng in rõ từng gân. Cô nàng Vu nữ đứng sau, tay cầm đũa tre, khẽ gõ vào thành khuôn cho bọt khí thoát ra. ``Như thế này, khi nguội sẽ không có lỗ rỗ.''

Khi mọi thứ đã đầy, cô em út gom khay khuôn vào một cái khay lớn, mặt căng như người đang bế trứng. ``Để mình mang ra phòng CLB nấu ăn!'' Cô bước đi, chân hơi chòng chành, như sợ một cái sẩy chân sẽ làm vỡ cả thế giới.

Chúng tôi nhìn theo, rồi quay lại nhìn nhau. Trên bàn, chỉ còn lại vài giọt sô-cô-la đã đông lại, như những vũng nước mùa thu bị quên lại. Kaigou khẽ cười: ``Giờ thì\dots\ ta chờ thôi.''

``Ừ,'' tôi đáp, tay khẽ vuốt khuôn lá còn nóng.

Tôi quay lại, mắt lướt qua từng góc lớp, nơi những thành viên còn lại của lớp vẫn đang tất bật làm những món sô-cô-la chất lượng nhất.

Shino lặng lẽ đứng cạnh nhóm Shion, thỉnh thoảng cúi xuống chỉnh lại độ nghiêng của khuôn, tay cô đặt lên vai Aku như muốn trấn an.

Kana và Kanade túm đầu vào một tập khuôn hình thù lạ mắt, từ ngôi sao, cánh sen, cả con cá heo, hai người tranh luận ríu rít xem hình nào dễ gây ấn tượng hơn.

Yae và Shiina thì im phăng phắc, chỉ có tiếng thìa gõ nhẹ vào thành âu, đều đặn như máy. Không cần nói, họ vẫn kịp gấp rưỡi tiến độ so với những nhóm khác.

Mari đứng thẳng người trước khay khuôn trái tim, hai tay buông thỏng. Khối sô-cô-la nâu óng ả vẫn còn nguyên trong nồi, chưa giọt nào được rót ra. Cô chỉ nhìn, như thể đang chờ một tín hiệu từ chính trái tim mình.

``Mari-san?'' tôi khẽ gọi, sợ làm vỡ mảng im lặng mỏng manh quanh cô. Cô giật mình, vai rung nhẹ, rồi quay lại với nụ cười gượng: ``À\dots\ chẳng là, mình đang nghĩ\dots\ nên đổ hình gì thì hợp. Cậu nghĩ sao, Kuon-san?''

``Đừng nghĩ quá lên làm gì,'' tôi nói. ``Cứ làm cái hình mà cậu thấy hợp với tính cách của mình là được. Một cái gì đó\dots\ thuộc về bản thân ấy.''

Cô gật chậm, mắt vẫn dán vào khuôn: ``Ừ, thuộc về mình\dots''

Suzune thì lon ton giữa hai dãy bàn, đầu ngoáy như trẻ con mẫu giáo lần đầu được chơi đồ hàng: ``Oa\~{} cái nồi này nhỏ xíu mà có vani thật kìa!''

``\textbf{Đừng chạm vào!}'' Azuki hét vọng từ đầu lớp, giọng vừa dứt thì Suzune đã rụt tay về, mặt hớt hải.

``Ể? Không được chạm à?'' cô khẽ giật bắn, như thể vừa phát hiện ra luật lệ mới.

Tôi thở ra một hơi dài, đủ làm mùi cacao nóng bốc lên mặt. Ồn ào. Lộn xộn. Nhưng mỗi tiếng cười, mỗi vệt bột bắn tung tóe đều nhấp nháy như đèn lồng đêm hội. Khác hẳn tối qua, khi chỉ còn mùi máu lẫn lá cây ẩm thấp.

``\textit{Tình cảm nảy sinh sau tái sinh à\dots}''

Câu ấy tự nhiên chạy vòng trong đầu tôi, nhịp theo tiếng thìa khuấy đều đặn của Yae bên kia. Đó không đơn thuần là một câu nói, mà là lời hứa thầm của cả thị trấn này: rằng sau khi đã chết đi một lần, họ vẫn dám yêu, vẫn dám tin, vẫn dám mở lòng đón những vị khách lạ như tôi.

\dots Chỉ có điều, liệu tôi có xứng đáng với tình cảm mà mọi người dành cho tôi? Sớm muộn gì, tôi cũng sẽ không còn sống ở nơi này nữa; chúng tôi sẽ quay trở về thế giới gốc của mình. Trong hai tháng ngắn ngủi ấy, tôi đã chứng kiến Shino và Sakura chào buổi sáng mỗi lần chúng tôi đi tới trường, Touka mỉm cười chỉnh lại cà vạt lệch của tôi trước khi lên lớp với những lời nhắc nhở nghiêm túc, Ayame nghịch ngợ giấu vội bức thư không ký tên trong ngăn bàn tôi, hay Mari đỏ mặt khi thấy tôi bất ngờ quay lại lớp sau giờ ra chơi. Những tình cảm ấy, dù nhỏ bé, lặng lẽ, nhưng lại sâu sắc đến mức khiến tôi lo sợ: nếu một ngày tôi biến mất, liệu họ có còn vững bước để tiếp tục không?

\ct

Khắp lớp 1-C, tiếng cười, tiếng nồi chảo va đập, tiếng khuấy đều tay hòa vào nhau thành một bản nhạc rộn rã; mùi cacao nóng quyện vani lan tỏa, bám vào từng mái tóc, từng tà áo. Kẻ còn loay hoay cân đường, người đã hớn hở rót sô-cô-la vào khuôn; vài bạn chạy qua chạy lại mượn thìa, mượn khuôn, khiến dãy bàn như con tàu đang chòng chành giữa biển người. Dưới ánh đèn huỳnh quang, những dòng sô-cô-la óng ả chảy ròng, phản chiếu bóng học sinh đang cúi sát, mắt sáng như thể mỗi giọt đều là một lời thủ thỉ muộn màng của mùa hè.

Ở góc gần cửa sổ, ánh nắng cuối giờ chiếu xiên qua khung kính, đậu lên hai bóng dáng cúi sát vào nhau. Trên mặt bàn, khuôn sô-cô-la hình ngôi sao nằm lọt thỏm giữa đối nguyên liệu, viền cánh sao đã bắt đầu tràn ra ngoài, loang thành một dải thiên hà thu nhỏ đen óng ả.

Kana nhíu mày khi thấy cô bạn Thiên sứ bắt đầu hơi quá tay: ``Bà đổ hơi nhiều rồi đó.''

Kanade không ngẩng lên, tay vẫn giữ thìa, ánh mắt hiện rõ sự tự tin: ``Không sao đâu, tôi tính toán rồi. Cứ để nó tràn một chút sẽ đầy đặn hơn.''

``\dots Nó đang tràn ra ngoài kìa,'' nàng Ác ma chỉ vào mép khuôn, nơi dòng sô-cô-la nóng hổi đang len lách xuống thớt gỗ, tạo thành những vệt đen nhánh.

Kanade giật mình, mắt tròn xoe: ``Thôi chết!''

Một giây im lặng. Cả hai cùng nhìn chằm chằm vào vũng sô-cô-la đang lan rộng, như thể đang chờ nó tự thu mình lại.

``\dots Thôi kệ,'' Thiên sứ cười trừ, tay cầm thìa quay vòng trong không trung như vẽ một chòm sao mới. ``Thành hình thiên hà luôn đi. Một ngôi sao không đủ chứa hết hai chúng ta, thì biến thành cả dải ngân hà vậy.''

Kana thở dài, vội lấy khăn giấy thấm vệt sô-cô-la sắp chạm tới cuốn sách bên cạnh. ``Đừng có sáng tạo lúc này! Làm vậy cả bàn sẽ dính bết đấy.''

Ở phía dưới họ, một cuộc chiến nhỏ đang diễn ra. Azuki đứng chắn trước bàn học, hai tay dang ra như thể đang bảo vệ một bảo tàng quý giá khỏi một cơn bão lỡm. Đôi mắt cô dán chặt vào Suzune, người đang háo hức với những dụng cụ nhà bếp mà có lẽ cô chưa được tận mắt thấy bao giờ.

``Oa! Mình có thể thêm cái này vào không?'' Suzune reo lên, giơ cao một lọ thủy tinh nhỏ, bên trong là chất lỏng màu xanh lét, lấp lánh như nước mắt cá rồng. Không ai biết nó từ đâu ra. Có lẽ cô vừa ``vô tình'' lượm được từ ngăn kéo nào đó trong phòng thí nghiệm hóa học mà chẳng ai để ý.

Cô nàng Mạo hiểm - dù có phiêu lưu tới cỡ nào, nhưng cũng biết rõ giới hạn ở đâu - không để cô bạn kịp nói tiếp. ``Không!'' một tiếng cấm như dao cắt nghẹt không khí. Cô lao tới, giật lấy lọ như thể đó là bom hẹn giờ: ``Cậu định biến lớp mình thành phòng thí nghiệm của bác Frankenstein à?''

Cô nàng Linh vật chu môi, mắt to tròn như thể vừa bị tước mất quà sinh nhật: ``Nhưng nhìn vui mà! Màu này lên sô-cô-la chắc chắn sẽ nổi bật lắm!''

``Chúng ta đang làm sô-cô-la,'' Azuki nhắc lại, từng tiếng như đóng đinh vào đầu, ``không phải thí nghiệm hóa học. Và cũng không phải phim kinh dị đâu!''

Suzune rụt vai, ``Ể\dots'' một tiếng thất vọng nhỏ xíu, rồi bĩu môi, lẩm bẩm điều gì đó nghe như ``nhưng mà nó có ánh kim tuyến kìa\dots'' trước khi lủi thủi quay đi, để lại Azuki đứng đó, ôm lọ chất lỏng màu xanh như đang giữ một quả trứng của rồng, không biết nên cất nó đi đâu cho an toàn.

Ngay giữa lớp, nơi bàn trung tâm được kê thành một quầy thao tác nhỏ, Wata xoay nồi cách thủy cho lớp nước ngoài cạn bớt, giọng vẫn tỉnh queo như phát loa phường: ``Nhớ đậy nắp thật khít, nước bắn vào sô-cô-la là coi như đổ sông đổ biển đấy!''

Touka đứng bên cạnh, tay cầm nhiệt kế đỏ au, khẽ gật, giọng đều đều như máy đếm: ``Nhiệt độ nên để ổn định 45-50 độ, quá 55 độ là đứt gãy tinh thể, ăn sẽ xước cả răng.''

Inari lùi một bước, hai tay níu tạp dề, mắt lia xuống đất, nhưng vẫn cố nhích lại gần nhóm, thì thào:
``\dots Cứ\dots\ từ từ thôi\dots\ Đừng có dọa họ như vậy chứ.''

Ayame dựa cửa sổ, cánh tay gấp trước ngực, mắt lim dim theo dõi. Hương cacao bốc lên, phủ một lớp mờ ảo lên khuôn mặt cô. Cô mỉm cười, khẽ nói như tự giãi bày: ``Mỗi người một nồi, một khuôn một vẻ, nhưng cùng đổ cả trái tim vào. Valentine năm nay, thấm thía hơn mọi năm rồi.''

Mari vẫn đứng đó, bên chiếc bàn gần cửa sổ, nơi ánh nắng cuối buổi chiều đang dần tắt. Trước mặt cô, khuôn tròn nhỏ nằm lặng lẽ, như đang chờ một lời giải thích cuối cùng. Cô không còn do dự nữa. Tay phải cầm nồi, tay trái đỡ nhẹ, cô nghiêng chậm, để dòng sô-cô-la nâu óng chảy vào khuôn, từng giọt, từng giọt, như thể đang đếm thời gian.

Tôi bước lại gần, không dám làm ồn. ``Làm khuôn tròn à?'' Cô không quay lại, chỉ khẽ ``ừ'' một tiếng, mắt vẫn dán vào khuôn.

``\dots Đơn giản thôi,'' cô nói tiếp, giọng nhỏ như thể đang nói với chính mình, ``nhưng chắc là đủ rồi.''

Ở góc bên kia, dưới ánh đèn huỳnh quang phủ một lớp vàng nhạt, Nanase và Honoka đang cúi sát vào nồi sô-cô-la vừa mới tan chảy. Không khí quanh họ như nặng hơn một chút, bởi cả hai đều có vẻ không chịu nhường nhau.

Cô cháu gái của hiệu trưởng trường khẽ nhấc thìa, khuấy nhẹ một vòng một cách cẩn thận, mũi nhíu lại: ``Tôi nghĩ nên thêm chút cacao đắng nữa. Vị này còn ngọt quá, không có chiều sâu.''

Cô bạn bờm Cáo nghe xong liền bật dậy, hai tay chống hông, mắt tròn xoe: ``Gì chứ? Đắng ư? Không được! Valentine mà, ai lại đi tặng sô-cô-la đắng như thuốc bao giờ!''

``Nhưng nếu chỉ toàn đường thì nó sẽ như kẹo dẻo, không còn là sô-cô-la nữa,'' Nanase đáp lại, giọng bình thản nhưng đầy kiên quyết. ``Tôi muốn người ăn cảm nhận được vị đắng nhẹ, rồi ngọt dần. Như thế mới tinh tế.''

Honoka hừ một tiếng, tay vẫy vẩy thìa trong không trung: ``Tinh tế gì chứ! Người ta tặng sô-cô-la là để thấy ngọt ngào, không phải để ăn xong rồi nhăn mặt! Cậu muốn người ta ăn xong phải uống cả cốc nước à?''

Nàng tóc xanh khẽ cười, mắt nhìn xoáy vào nồi: ``Ít ra mình có gu của người trưởng thành. Không như cậu, còn mê đồ ngọt kiểu trẻ con lắm.''

``Cậu\dots\ vừa nói ai trẻ con hả!?'' nàng Hề nghiến răng, má bánh bao phồng lên, như thể chỉ cần thêm một giây nữa là cô sẽ nhảy vào nồi sô-cô-la luôn.

Cả hai đứng đó, chẳng ai chịu nhường, giữa khói trắng bốc lên nghi ngút, như thể đang không chỉ nấu sô-cô-la, mà còn đang nấu chín cả tình cảm của mình.

Ở góc cuối lớp, Shiina và Yae dựng một chiếc bàn nhỏ riêng biệt, như thể muốn tách khỏi cái ồn ào của thế giới. Không ai chỉ huy, cũng không ai reo hò, chỉ có tiếng thìa inox khẽ chạm vào thành âu, đều đặn như nhịp kim đồng hồ.

Giữa lúc cả lớp còn đang loay hoay tìm cách cứu vãn những mẻ sô-cô-la cháy khét hay đóng cục, hai cô gái ấy chỉ lặng lẽ làm việc. Yae cúi xuống, mắt dán vào nồi cách thủy, tay khuấy chậm rãi theo một vòng tròn hoàn hảo. Shiina đứng đối diện, lặng lẽ lọc bột qua rây, mặt không một chút cảm xúc.

Thỉnh thoảng, tiếng nói vang lên, nhưng chỉ đủ để hai người nghe thấy.

``Được chưa?''  Yae hỏi, không ngẩng lên.

``Ừ.''  Shiina đáp, tay vẫn rây bột.

Một lúc sau, khi mùi cacao đã ngậm đủ hơi nóng, nàng bờm Sư tử khẽ gật: ``Đổ đi.''

Cô bạn Thể lực không nói gì, chỉ lặng lẽ nghiêng nồi. Dòng sô-cô-la óng ả chảy vào khuôn, không một bọt khí, không một vệt loang. Shiina đặt khuôn xuống khay, tay khẽ gõ nhịp, như thể đang đọc một bản nhạc im lặng.

Cứ thể, hai người họ chỉ tập trung làm việc, không bình phẩm lấy một câu. Nhưng khi khuôn được mở ra, những viên sô-cô-la nâu bóng nằm im lặng trong lòng khuôn, tròn trịa, mịn màng, như thể chúng chưa từng biết đến sự vội vã hay lỗi lầm.

Và trong im lặng ấy, họ chỉ nhìn nhau một lúc, rồi tiếp tục. Quả thật, người ít nói thường hay có năng suất làm việc tốt hơn, như những gì mà bố tôi từng dạy.

Suzuki không chịu đứng yên. Cứ mỗi khi một nhóm lên tiếng reo hò hay tiếng ai đó thất thanh khi gặp chuyện, em ấy lại lách qua khe bàn, váy xoè nhẹ như lá bàng rơi sán tới. Tay em thoăn thoắt: nâng nhẹ cổ tay người này để họ khỏi quậy điên cuồng, xoay vặn lửa cho người kia, hoặc chỉ đơn giản là gạt bớt bọt khí đang trào ra mép âu.

``Chỗ này khuấy nhẹ tay thôi, kẻo bột nổ tung đấy,'' Suzuki cười trước vẻ mặt hoảng hốt của Inari, người vừa bị bột bắn tung lên mặt.

``Để nguội thêm xíu nữa, đổ ngay bây giờ là sẽ loang đó,'' em ấy gỡ chiếc khuôn khỏi tay Suzune đang hấp tấp, giọng dịu như ru.

``Không sao, còn cứu được mà, tí nữa cho thêm xíu sữa rồi đánh lại,'' Cô vỗ vai Mari đang nhăn nhó vì khối sô-cô-la vón cục, đôi mắt nửa trách nửa an ủi.

Không dạy bảo, không ra lệnh, mà chỉ bảo một cách nhỏ nhẹ, cẩn thận chỉnh lại, rồi bỏ đi sau khi nhận lời cảm ơn. Vậy mà cả lớp đều quay lại nghe theo, như thể có một dây chuyền bí mật truyền âm thanh từ giọng nói cô.

Tôi khẽ huých khuỷu Kaigou. ``\vie{Không ngờ là em gái cô lại tháo vát như thế đấy. Tới mức còn hướng dẫn cả các bạn nữa kìa, mặc dù đây là lần đầu tiên em ấy là sô-cô-la.}''

Cô chị không đáp ngay. Cô đang nhìn theo bóng Suzuki cúi xuống lau vũng sô-cô-la Suzune vừa đánh đổ; tóc cô rủ xuống che khuôn mặt, chỉ còn lộ đôi tai hơi đỏ vì hơi nóng. Một lúc sau, Kaigou mới ``ừ'' khẽ, giọng trầm hơn mọi khi. ``\vie{Nó lúc nào cũng vậy. Nhưng giờ\dots\ có lẽ cũng đã trưởng thành phần nào rồi. Như là một cách để gắn kết với các bạn trong lớp hơn.}''

``\vie{Công nhận.}''

Khi nhịp ban đầu hạ xuống, cả lớp bước vào giai đoạn ``muôn hình vạn trạng'', mỗi người bắt đầu chút chắt cho sản phẩm củ của mình. Người cúi rây bột, kẻ khều lửa nhỏ, tiếng thìa inox gõ nhẹ vào thành âu nhịp nhàng như mưa rơi trên mái tôn. Những câu chuyện râm ran không dứt hẳn mà thu mình lại, chỉ còn đủ để hai người ngồi cạnh nhau nghe thấy.

Nhóm chúng tôi đã xong từ lâu. Bơ cacao đã tan, đường đã tan, khuôn đã xếp thành hàng lá mùa thu trên khay, tất cả đều đã được Suzuki đem hết vào tủ lạnh và chỉ chờ tới lúc hoàn thành công đoạn trang trí phần cuối.

Kaigou dựng cằm lên thành ghế, mắt liếc qua đám khói bốc lên từ nồi của nhóm Shion, mỉm cười khẽ: ``\vie{Mọi thứ cũng khá ổn,}'' cô thủng thẳng, ``\vie{may là chưa có vụ nổ nào sập lớp.}''

Tôi gật, vừa kịp né đứa bạn chạy qua tay bưng khay nước: ``Ừ, ít nhất thì lần này chúng ta sẽ không phải gọi cứu hỏa nữa.''

\ct

Mười một giờ rưỡi trưa, ánh nắng đứng bóng ngoài hành lang lọt vào lớp học, loang vệt vàng lên những khay sô-cô-la vừa được đổ đầy. Không khí nóng hổi của bếp ga mini hãy còn phảng phất, nhưng đã lẫn mùi vani ngọt lịm, báo hiệu mọi thứ đang bước sang giai đoạn ``nghỉ ngơi'' để đông cứng.

Ayame vỗ tay hai cái thật to, giọng vang lên trên cả lớp: ``\textbf{Tạm dừng nào, hết giờ rồi!}''

Wata đứng bên cạnh, nụ cười vẫn tươi như sáng sớm: ``Mọi người nhẹ tay xếp thành phẩm vào tủ lạnh nhé. Nhớ dán nhãn tên cho cẩn thận, kẻo chiều lấy ra không biết của ai với của ai đâu!''

Tiếng ghế kéo lê rào rào, tiếng bước chân lộp cộp. Các khay nhôm, khuôn silicon, hộp nhựa được bưng lên, xếp hàng dọc theo lối đi. Những viên sô-cô-la hình lá, hình trăng khuyết, hình ngôi sao hay đơn giản chỉ là khối tròn bóng mỡ còn rung rinh theo nhịp bước, như thể cũng háo hức được vào lạnh. Ở đâu đó, Suzune reo lên: ``Nhìn viên hình trứng của em nè, tròn vo như trăng rằm!'' cùng với tiếng cười đáp của Azuki: ``Khéo đóng gói kẻo nó bung ra nhé.'' Hanabi xoa khay ngôi sao, híp mắt: ``Vụn tình yêu đây, ai dám nhận?'' còn Koishin lẩm bẩm: ``Chỉ cần không bị trả lại là ổn rồi\dots''

Hội trưởng nhìn theo các bạn, gật gù hài lòng, mắt khẽ liếc sang Ayame như chờ lời bình. Ayame mỉm cười, giọng nhẹ nhàng nhưng đủ để cả lớp nghe thấy: ``Ở đây, mỗi người đều đang nấu một món tình cảm riêng. Không cần hoàn hảo, chỉ cần thật lòng là đủ rồi.''

Touka đứng cạnh bàn, tay vẫn cầm nhiệt kế, mắt lướt qua từng nhóm. Cô khẽ nói, giọng đều đều nhưng không lạnh lùng: ``Mình cứ nghĩ Valentine chỉ là một ngày bình thường. Nhưng nhìn mọi người hăng say tới mức đó\dots\ mình bắt đầu hiểu vì sao người ta lại cần một ngày để nói ra điều thường ngày không dám nói.''

Inari lùi về phía cuối lớp, tay vẫn níu tạp dề, mắt nhìn xuống đất. Cô thì thào, như nói với chính mình: ``M-Mình không dám nghĩ ai sẽ nhận sô-cô-la của mình đâu. Nhưng mà\dots\ chỉ cần làm thôi, cũng đã là một bước tiến rồi phải không?''

Wata đứng trên bục giảng nhỏ, tay chống hông, nhìn cả lớp với ánh mắt hiếm khi thấy được: dịu dàng. Cô cười khẽ: ``Mọi người cũng đã cố gắng hết sức rồi đấy chứ. Từng viên sô-cô-la mà họ đã làm vì người dành tình cảm nhất\dots\ đó cũng là một kỷ niệm tuổi học trò đáng nhớ đấy chứ.''

Trong dòng người đang mang sản phẩm của mình ra với những sự mong chờ, tôi khẽ dựa lưng vào vách tủ, nhìn theo bóng Mari còn chưa kịp khuất hẳn ngoài hành lang. Mùi cacao nồng vẫn bám trên tà áo, nhưng điều khiến tôi bâng khuâng lại là những ánh mắt lén lút đang lướt qua mình: một sự vội vã, e thẹn, rồi lại chạy đi nơi khác. ``\vie{Tôi có cảm tưởng rằng mọi người đang nhắm tới mình\dots}'' tôi lẩm bẩm tay vô thức siết nhẹ chiếc khuôn lá còn nóng, ``\vie{nhưng chúng ta chỉ là kẻ qua đường, khi lời nguyền ở nơi đây được giải quyết thì cũng là lúc chúng ta đi. Nhận lấy tình cảm này\dots\ tôi thấy thế nào ấy}''

Kaigou nghe vậy, đặt ly trà xuống bàn, mắt vẫn dõi theo dáng Shino đang cúi xuống nhặt đũa cho Aku. ``\vie{Người được chờ đợi thường không hay biết mình đang được chờ,}'' cô khẽ nói, giọng trầm ấm. ``\vie{Nhưng nhìn cách cả lớp cố gắng\dots\ tôi nghĩ, dù chúng ta có biến mất đi nữa, họ vẫn muốn gửi lại một chút hơi ấm, để sau này, khi nhớ lại, không phải hối tiếc vì chưa kịp nói.}'' Cô quay sang tôi, cười nhẹ: ``\vie{Đừng tính toán xứng hay không. Cậu cứ nghĩ đó là quà chia tay điđi.}''

Suzuki vừa xếp xong khay trứng, nghe vậy ngẩng lên, hai má còn phủ bột trắng. ``\vie{Em thì đơn giản hơn,}'' em cười híp mắt. ``\vie{Em chỉ thấy\dots\ mọi người hạnh phúc khi được làm điều gì đó cho người mình thích. Nếu anh cứ từ chối, họ sẽ buồn lắm đó. Nên là, Onii-sama ạ, anh cứ nhận đi, rồi sau này, khi anh đã về thế giới của mình, nhớ mùi cacao này là được.}'' Em ấy chỉ vào tủ lạnh đang kêu lách cách, trong đó bao nhiêu viên sô-cô-la còn đang chờ đông cứng, mỗi viên một trái tim, một lời chưa kịp thốt. ``\vie{Đó là cách họ níu kéo anh lại bằng vị ngọt.}''

Tôi nhắm mắt, khẽ ngẫm lại những gì mà hai chị em họ vừa nêu quan điểm. Kaigou nói đúng: đây chỉ là quà chia tay, không cần quá khệ nệ hay đắn đo làm gì. Nhưng liệu tôi có đủ dũng cảm để nhận, rồi lại đủ bạc tình để rời đi?

Ý tưởng của Suzuki cũng đâu đó gần giống cô chị: ``Anh cứ nhận đi, để họ được vui.'' Em ấy quả thật vẫn còn khá ngây thơ với đời, mà không hiểu rằng, niềm vui nhãn tiền đó sau này sẽ hóa nỗi nhớ, và ai sẽ gánh? Tôi sợ mình trở thành vết thương không tên, chậm liền trong lòng họ.

Nhưng rồi nghĩ lại, có lẽ, sự tồn tại của tôi ở đây chẳng khác gì mùi cacao: ngọt ngào rồi tan, chỉ còn lại hương vị lẩn khuất trong ký ức, và cái đắng - dù nhẹ tới đâu - cũng sẽ tới sau cùng, để lại một hậu vị trầm mặc mà nhung nhớ. Vậy thì, chắc tôi cứ mạnh dạn nhận, để sau này khi nhớ, họ còn biết mình đã yêu. Tuổi thanh xuân của con người chính là thời điểm mà tình yêu trở nên mãnh liệt nhất, và cũng là thời điểm cuối cùng mà họ có thể bày tỏ tình cảm của mình trước khi dấn thân vào guồng quay của xã hội, nơi mà thời gian cho bản thân thôi cũng đã là điều xa xỉ.

Tôi thở dài, siết nhẹ tay như trao gửi một lời hứa vô hình. Dù chỉ là kẻ qua đường, tôi sẽ mang theo hương cacao này về thế giới cũ. Để mỗi lần nhớ, tôi cũng thấy mình từng được yêu, dù chẳng thể đáp lại.

\ct

Không lâu sau cuộc trò chuyện ngắn gọn mà triết lý giữa ba anh em chúng tôi, từng bạn học quay trở lại, tiếng ghế kéo nhẹ vang lên. Từng người, từng người một quay trở về chỗ ngồi của mình, không ít trong số họ vẫn còn bàn bạc về những món sô-cô-la mà tự tay họ làm.

Ayame đập bàn, mắt giương xuống các bạn học, rồi dõng dạc: ``Buổi chiều hôm nay, trong lúc chờ sô-cô-la đông lại, mọi người có thể viết lời nhắn cho người mình định tặng.'' Cô khẽ đưa tay lên miệng, ra hiệu nhỏ nhẹ: ``Viết xong, thì bỏ vào tủ giày của họ, hoặc giấu trong ngăn bàn cũng được.''

Wata gật đầu, bổ sung thêm: ``Nói ra thì hay, nhưng giữ kín cũng chẳng sao. Quan trọng là người nhận phải cảm được mùi vị sô-cô-la mà các bạn đã dày công làm, và cả tấm lòng mà các bạn đã bỏ tâm huyết vào nữa.''

Hội trưởng Hội học sinh lách qua khe bàn, tay ôm xấp giấy A5 cắt tỉa cẩn thận, mỗi tờ đều có in hình trái tim bé xíu ở góc. Hội phó theo sau, nhẹ nhàng đặt lên bàn từng cây viết mực màu hồng phấn như thể đang bày biện một bữa tiệc thơ. Cả hai không nói gì, chỉ mỉm cười, nhưng tôi bắt gặp ánh mắt họ lướt qua từng khuôn mặt: những ánh mắt chờ đợi, hơi lo lắng, như người gieo hạt giữa mùa đông không biết đất có nở hoa hay không.

Tôi ngồi tránh nhẹ sang bên, cố tình tránh khỏi dòng người đang xếp hàng. Một tờ giấy vẫn được nhét vào tay, còn nóng hơi cơ thể Ayame. Trên bề mặt mỏng manh, hình trái tim in chìm như một vết chàm nhỏ, gợi nhắc điều gì đó mà tôi vẫn cố quên: rằng mình cũng được phép yêu thương, dù chỉ trong hai tháng ngắn ngủi.

Tôi thở dài, kéo ghế ngồi xuống chỗ ngồi. Bên tai vẫn văng vẳng tiếng khóc rít của những hồn ma đêm qua, những người chưa kịp nói ``thích'' đã vội tan biến. Valentine ở thế giới này rơi vào tháng Năm, khi mùa hè chưa gõ cửa, khi ve chưa kịp ra rả; nó lệch nhịp so với trái tim tôi, vốn quen đập theo mùa đông tháng Hai. Tôi vẫn chưa hết choáng vì sự trật đường ray ấy, như kẻ lạc đường bước vào nhà ga lúc nửa đêm, thấy đèn pha leo lét mà không biết tàu nào đưa mình về.

Có lẽ nên viết cho Shino, người đã từng chịu vòng lặp lời nguyền chỉ vì muốn có một người bạn; hoặc Sakura, người đã đích thân hồi sinh thế giới này, dù chưa trọn vẹn. Nhưng nghĩ lại, tôi chỉ là cầu nối đưa họ trở lại thế giới này; còn tình cảm, họ đã dành trọn cho nhau rồi.

Những người còn lại thì\dots\ cũng không khá khẩm hơn. Mari thì thường đỏ mặt, nhưng có thể chỉ vì cô ấy dễ ngại; Touka thì mỉm cười nhưng cũng mỉm với cả lớp; Ayame luôn trêu chọc, Wata luôn dịu dàng (tôi không biết ranh giới giữa quan tâm và thương nhớ nằm ở đâu), hay Azuki ưa phiêu lưu, có lẽ cũng không để tâm lắm. Hai tháng qua, tôi như kẻ đi trên cầu kính: thấy dưới chân là mặt nước rung rinh bóng người, nhưng không dám bước xuống vì sợ vỡ.

Chẳng khác gì như bố mẹ tôi từng nói: tôi là một thằng tồ. Một thằng tồ trong chuyện tình cảm. Nhiều lúc, các cô gái ở \hll\ đã ra dấu hiệu, nhưng tôi - theo góc nhìn của họ - thường lờ tịt đi, hoặc không hiểu họ đang làm vì gì, vì ai. Với tôi, chuyện tình cảm là một cái gì đó quá xa vời, khi đến tôi đôi lúc còn không hiểu được tình cảm của họ dành cho tôi ra sao, vả lại cũng không hứng thú gì mấy cái chuyện yêu đương này mấy.

Một khoảng lặng kéo dài, đủ để tiếng ghế kéo, tiếng cười khúc khích, tiếng mực kêu *XOẸT* trên giấy biến thành nhạc nền xa xôi. Tôi nhìn thấy Kanade viết xong, gập tờ giấy lại, mắt ánh lên niềm tin chắc nịch; thấy Inari cắn đầu bút, mặt đỏ như gấc; thấy Suzune vẽ thêm một ngôi sao nhỏ bên cạnh dòng chữ, rồi ôm lấy khay trứng như thể cả thế giới vừa gọi tên cô ấy. Còn tôi, vẫn trơ trọi một dòng đầu: ``Gửi\dots''

``\dots Thôi. Nếu mình không thể quyết định được viết cho ai, vậy thì\dots''

Tôi thở ra, cảm thấy hơi ấm cacao còn vương trên đầu ngón tay. Tôi không viết tên người nhận, cũng không ký tên mình; chỉ viết, như kẻ ghi chép lại giấc mơ sợ quên, như người đánh điện tín giữa đêm cho một kẻ không tên. Những dòng chữ trôi đi không địa chỉ, như bưu thiếp thả xuống sông, mong dòng nước đưa tới nơi cần tới, dù chính tôi cũng chưa biết đó là đâu.

Khi tôi gập tờ giấy lại, tiếng giấy rúc nhẹ như tiếng lá rơi. Tôi chợt nghĩ: có lẽ không cần biết người nhận là ai; chỉ cần mình còn dám viết, còn dám gửi, còn dám nhớ, là mình đã không hoang phí mùa hoa này, dù hoa nở trên đất lạ, và mùa sẽ tàn khi tôi rời đi.

Khi mọi người gần như đã hoàn thành phần của mình, lớp học chỉ còn tiếng bút mực kêu xoẹt và tiếng gấp giấy nhẹ nhàng. Tôi vừa gập tờ thư lại, thở phào nhẹ nhõm, thì bất chợt\dots

``Kuon-kun\~{}''

Tiếng gọi của Ayame cắt ngang dòng suy tưởng. Chậc. Nghe cái giọng kiểu châm chọc này là biết sắp có điềm rồi.

``\dots Gì?'' Tôi quay lại, tay còn giữ tờ giấy.

``Cậu viết xong chưa?''

``\dots Rồi.''

``Đọc đi.''

``\dots Hả?'' Tôi biết ngay mà.

Ngay khi cô nàng Quỷ dứt câu, cả lớp quay lại. Có người ngừng bút, có người nghiêng đầu, có người cười khúc khích. Ánh mắt tò mò bắt gặp tôi như muốn đọc cả suy nghĩ.

``\dots Không cần đâu.'' Tôi lùi lại nửa bước.

``Đọc đi\~{}'' Ayame cười, mắt híp lại như mèo vừa bắt được chuột.

``Đọc đi!'' Azuki hưởng ứng, giơ nắm đấm lên như cổ vũ.

``Đọc đi.'' Kaigou nói, giọng bình thản nhưng không kém phần quyết đoán.

``\dots Sao cô cũng theo họ vậy?'' tôi nhìn cô nàng mắt xanh, cảm giác như bị bỏ rơi giữa trận.

Và cứ thế, hàng loạt tiếng ``Đọc đi'' của cả lớp cứ thế lấn át cà phòng, như thể đang thúc giục tôi chia sẻ cảm nhận của mình.

Tôi thở dài, rõ ràng không còn đường lui. ``\dots Rồi. Đọc thì đọc.''

Tôi mở tờ giấy ra, cả lớp lập tức im lặng, chờ đợi những lời bày tỏ cảm nghĩ của tôi sẽ được dành cho ai. Một vài người khẽ dịch ghế gần lại. Tôi đọc, giọng không to nhưng đủ để cả lớp nghe:

\txtbx{
  Kính gửi các bạn, những người đã cùng tôi nếm vị cacao nồng ấm trong buổi chiều tháng Năm lạ lùng này,\\
  Tôi thật sự không biết nên đặt tên người nhận nào ở đầu thư, vì e rằng bất kỳ cái tên riêng nào cũng sẽ làm kém đi sự cân bằng mong manh vừa được gieo giữa chúng ta.\\
  Shino-san, Azuki-san, Nanase-san, Shiina-san, Aku-san - năm người đầu tiên tôi nghĩ đến - đã dạy tôi rằng sự hiền hòa, mạo hiểm, ngay thẳng, trầm lặng hay hậu đậu đều có thể hòa quyện trong cùng một lớp nước cacao sóng sánh.\\
  Shion-san, Koishin-san, Sakura-san, Touka-san, Suzune-san thì lần lượt gieo vào nồi sô-cô-la ấy chút điên rồ của tuổi dậy thì, vẻ tsundere bối rối, bí ẩn như khói, trí tuệ sắc lạnh và niềm vui nhí nhảnh không bao giờ cạn.\\
  Honoka-san, Mitsuki-san, Kaoru-san, Inari-san, Hanabi-san, những người khiến tôi hiểu thế nào là sắc màu hoạt bát, dịu nhẹ, tốt bụng, e dè và tự tin. Họ đã cùng nhau tạo nên thứ hương vị mà tôi e rằng sẽ chẳng bao giờ gặp lại nơi phương trời khác.\\
  Hotaru-san, Saki-san, Yae-san, Yuki-san, Mari-san, bộ năm đáng tin, mê hoa, đầy khát vọng thắng thế thao, yêu ca hát và bay bổng hội họa, đã khẽ khàng gõ vào thành khuôn nhịp nhàng như bản giao hưởng chỉ có thể cất lên trong buổi chiều tháng Năm này.\\
  Kana-san, Kanade-san, Akane-san, Yuka-san, Saya-san, Uzuki-san - những người hoạt bát, kín đáo, nhốn nháo vui tính, mê tòm tem chòm sao, đam mê nhiếp ảnh - đã chứng minh rằng khi mọi người cùng hướng về một điểm, kể cả thời gian cũng phải khựng lại vài nhịp.\\
  Và cuối cùng, Miku-san, Ayame-san, Wata-san, Kaigou-chan và Suzuki, những người trưởng thành dù còn phần trẻ con, tinh nghịch, êm dịu, và hai chị em đồng hành, là những người cuối cùng hoàn tất bức tranh, khiến tôi hiểu ra rằng điều kỳ diệu nhất không nằm ở viên sô-cô-la hoàn hảo, mà ở cái cách mỗi người đều dám đổ trái tim mình vào khuôn.\\
  Vì vậy, nếu phải nói - thưa cả lớp - thì tôi xin thú thật: khoảng thời gian được ở cùng mọi người, trong căn phòng ngập mùi cacao và tiếng cười này, đã trở thành điều đủ đặc biệt để tôi mang theo, dù mai kia chúng ta có phải chia xa.\\
  Xin cảm ơn, và nguyện mỗi viên sô-cô-la các bạn cầm trên tay đều chảy ngọt ngào theo đúng nhịp tim mình mong muốn.''}

Không khí như nghẹn lại. Cánh cửa sổ hé mở, gió lùa vào làm tung nhẹ rèm cửa, nhưng không ai buồn đỡ. Tôi còn nghe rõ tiếng giấy mình vừa gập lại đang rung nhè nhẹ trong lòng bàn tay. Một giây, hai giây\dots\ cả lớp khựng lại như bị ai bấm nút tạm dừng.

Tiếng thở dài đầu tiên thoát ra từ khóe môi Azuki. Cô nàng ôm trán, mắt trố tròn nhìn tôi: ``Kuon-san, cậu vừa\dots\ \textbf{vừa tỏ tình với cả lớp luôn á!?}'' Giọng cô run run, giữa sửng sốt và cảm động không phân biệt nổi.

Shion đứng phắt dậy, ngón tay run rẩy chỉ thẳng vào mặt tôi như thể vừa phát hiện kẻ phá vỡ thế giới quan: ``Cậu… cậu là Dark Lord của mùa Valentine sao!? Lời nguyền thứ nhất: khiến trái tim mọi cô gái đồng loạt rơi vào hỗn loạn!?'' Cô lắp bắp, hai má đỏ bừng, nhưng mắt lấp lánh hào hứng của kẻ vừa tìm ra cốt truyện mới.

Suzune thì ôm bụng cười khanh khách, tóc bờm xoăn tung bay: ``Ôi trời, mình tưởng cậu sẽ viết cho một người, ai dè Kuon-san lại\dots\ chọn luôn cả lớp! Kuon-san bá đạo quá đi!''

Honoka đập bàn, má bánh bao phồng lên: ``Gì chứ, tớ còn chưa kịp đắp chiếu cho mẻ sô-cô-la của mình, đã bị cậu đánh úp bằng cả bài diễn văn như thế này!?'' Nhưng nói xong cô lại cười toe toét, véo nhẹ tay Nanase như muốn khoe: ``Nghe thấy không, tớ cũng được nhắc tên kìa!''

Nanase gãi thái dương, vẻ bất đắc dĩ: ``Tỏ tình kiểu học thuật vậy\dots\ cũng chỉ có cậu ta mới nghĩ ra. Nhưng mà\dots'' cô khựng lại, khóe miệng nhếch nhẹ, ``\dots nghe không tệ nhỉ.''

Mari lén lút đứng sau lưng Touka, thì thào như sợ ai nghe thấy: ``C-cậu ta nói cả lớp đều đặc biệt á\dots'' mặt cô đỏ ửng, hai tay siết chặt gấu váy.

Touka từ tốn đáp lại, giọng có chút hứng thú: ``Không kể sót một người nào luôn. Cậu cũng có tài để ý người khác đấy nhỉ?'' Dứt lời cô khẽ cười, mắt nhìn tôi như vừa ghi thêm một chỉ số thú vị vào bảng thí nghiệm.

Koishin ngồi thụp xuống bàn, hai má bỗng dưng nóng ran, cắn môi không dám ngẩng. Còn Shiina thì vẫn đứng yên như tượng, chỉ có đôi mắt khẽ chuyển từ mảnh giấy sang tôi, ánh nhìn bình thản như vừa đọc xong một trang sách dày.

Saya vỗ tay vào trán, cười khúc khích: ``Tưởng cậu ta lạnh như băng, ai dè lại là kiểu đa thê bằng văn chương thế kia\~{}''

Mitsuki lắc đầu, mái tóc hồng-tím bồng bềnh rung rung: ``Nghe xong muốn đổ thêm đường vào sô-cô-la quá, kẻo khéo có người lại bảo chúng ta ngọt chưa đủ na!''

Saki đưa tay lên miệng, khẽ cười: ``Mình cũng được nhắc tên trong bài thơ dài thế kia rồi. Có lẽ nên trồng thêm một luống hoa hồng mới để đáp lễ mới được.''

Miku thì xoa cằm, mắt lim dim: ``Cậu ta nói `mỗi người đều đặc biệt\dots' Nghe mắc cỡ quá đi, nhưng mà tim mình vẫn đập loạn nhịp là sao?''

``\dots Thế nên tôi mới không muốn đọc đấy,'' tôi thở dài, thả tờ giấy xuống mặt bàn, cam chịu những lời đàm tiếu từ các bạn nữ trong lớp.

Kaigou quay mặt đi, tay đưa lên che miệng, nhưng tôi vẫn kịp thấy khóe môi cô nhếch lên: ``\dots Tự chuốc thôi.'' Giọng cô trầm, mang theo chút cười nhạt, như thể đã đoán trước kết cục này nhưng vẫn mắc cười khi cái cậu con trai lạc lõng lại chọn cách gom cả vườn hoa về một giỏ.

Suzuki đứng cạnh tủ lạnh, hai tay còn dính bột, mặt đỏ như gấu bông vừa được phơi nắng. Em ấy lí nhí: ``\dots Onii-sama\dots\ T-Thật sự\dots\ em chưa bao giờ được gửi những lời yêu thương như vậy bao giờ cả\dots\ Em không biết phải phản ứng thế nào nữa\dots''

Tôi thở dài, cảm thấy buổi chiều tháng Năm bỗng ngắn lại đáng kể. Có lẽ, khi mọi người đã cùng nhau đỏ mặt, cười rộ, rồi lại im lặng nhìn nhau, thì khung cảnh ấy đã đủ để tôi nhớ cả đời.

\ct

Sau khi tôi gập bức thư tỏ tình với cả lớp lại, không khí trong phòng học như được nâng lên một nấc mới. Các cô gái không còn e thẹn nữa, họ bắt đầu rủ nhau lên kế hoạch ``giai đoạn hai'' của ngày lễ. Một nhóm kéo nhau ra cửa sổ, tranh thủ phơi nắng những tấm giấy gói hoa văn; nhóm khác lục tung ngăn bàn tìm dây ruy-băng, hạt cườm và thậm chí cả bình xịt nước hoa mini. Tiếng cười rộ lên, xen lẫn tiếng ``đừng giật chỗ tôi'' hay ``cái nơ này phải xoắn hai vòng mới đẹp'', khiến căn phòng trở thành xưởng thủ công khổng lồ.

Lúc kim đồng hồ chỉ hai giờ rưỡi, Ayame vỗ tay gọi mọi người tập trung. Cô nàng cầm trên tay một xấp thẻ nhỏ in sẵn hình trái tim, mặt sau để trống để ghi địa điểm và giờ hẹn. ``\textbf{Nào, ai cũng phải có một cuộc hẹn riêng chứ!}'' Hội trưởng hét lớn, mắt liếc qua tôi một cách đầy ẩn ý. Wata đứng bên cạnh bổ sung bằng giọng dịu dàng nhưng không kém phần hào hứng: ``Hãy nhét tờ thư vào tủ giày của người mà bạn thầm mến, rồi đến đúng địa điểm ghi trên thẻ. Đừng quên mang theo sô-cô-la của mình nhé!'' Cả lớp reo hò, như thể vừa được phát lệnh xuất kích.

Ba giờ chiều, hành lang trường bỗng trở nên đông nghẹt. Những cô gái lén lút đi giày nhẹ, len lỏi qua dãy tủ sắt, nhét thư vào khe cửa rồi bỏ chạy như chim non vừa học bay. Không chỉ lớp tôi, những lớp khác cũng làm tương tự như vậy; những chàng trai từ tuổi mới lớn ở các lớp bên cũng hào hứng với màn tỏ tình đầu tiên trong những năm học cấp ba. Tôi đứng tựa lan can, vừa xem vừa cảm thấy một luồng điện nhỏ chạy dọc sống lưng: mỗi bước chân của họ như đang kéo mình vào một vở kịch mà kịch bản đã được viết sẵn, chỉ thiếu mình chưa biết vai chính là ai.

Bốn giờ mười lăm, tiếng chuông báo giải lao vừa dứt, cả lớp lại quay trở về bàn như đàn chim non kéo về tổ. Không khí lúc này khác hẳn buổi sáng: người người hớn hở, kẻ kẻ rạng rỡ, tiếng cười râm ran hòa vào mùi cacao còn vương trên tóc, trên tay áo. Kaigou kéo ghế ngồi cạnh tôi, tay vẫn cầm cốc trà, mắt liếc qua cửa sổ: ``Mặt trời hạ xuống tây rồi, còn mấy viên trong tủ chắc cũng cứng lại thôi.''

Suzuki ngồi phía trước, quay lại cười híp mắt: ``Em vừa kiểm tra, khay của anh đã bóng mượt như gương rồi đấy.'' Nói xong em khẽ đưa tay chỉ vào tôi, rồi quay sang Kaigou nháy mắt, như thể hai chị em đang giữ một bí mật chỉ riêng họ hiểu. Tôi chỉ biết cười trừ, mắt dõi theo dòng người ra vào, tim thì lúc nào cũng như đang đếm từng phút.

Năm giờ chiều, tiếng chuông cuối cùng reo vang. Cửa tủ lạnh được mở tung, khói lạnh bốc lên như sương sớm. Mọi người ùa tới, tay đeo găng nilon, mắt sáng lên như trẻ con được mở quà. Từng khay sô-cô-la được nhấc ra, lóng lánh dưới ánh đèn huỳnh quang. Có người thở phào vì viên không bị nứt, có người cắn môi vì cạnh bị sứt; nhưng rồi tất cả đều được xếp vào hộp, như những mảnh trái tim vừa được gọi tên.

Trên bàn, giấy gói màu hồng phấn, xanh nhạt, tím oải hương được trải ra thành dải. Dây ruy-băng óng ả, hạt cườm lấp lánh, tem nhãn hình trái tim bé xíu được xếp thành hàng. Có bạn dùng kẹp giấy tạo hình cánh bướm, có bạn cắt tỉa giấy decoupage thành đóa hoa năm cánh.

Chiếc nơ được thắt lại ba lần, bị tháo ra hai lần, rồi lại được chỉnh sửa đến lần thứ tư. Mỗi người một vẻ: người thì cẩn thận ghi tên mình bằng mực vàng, kẻ thì dán thêm viên kẹo nhỏ bên cạnh như mật ngọt bí mật. Tôi ngồi đó, nhìn họ, như kẻ lạc vào xưởng thủ công của thiên thần, nơi mỗi món quà là một lời thì thầm chưa kịp nói.

Ayame và Wata bất ngờ xuất hiện trước mặt tôi. Hai cô nàng cùng cầm một chiếc hộp giấy hồng nhạt, trên nắp gắn nơ bướm to tướng. ``Chúc cậu may mắn, Kuon-kun!'' Hội trưởng cười tinh quái, đẩy nhẹ hộp vào tay tôi. Hội phó khẽ nghiêng đầu: ``Đừng có mà chạy trốn nhé!'' Nói rồi hai cô quay đi, vai rung lên vì cười, như thể vừa thả một quả bóng bay vào bầu trời và chờ xem nó sẽ bay về phía nào.

Tôi nhìn theo bóng họ khuất sau cánh cửa, tim đập chậm nhưng mạnh. Tôi thở dài, không biết nên cảm thấy lo lắng hay chờ đợi. Nhưng có một điều chắc chắn: từ lúc viết xong bức thư đến giờ, cả lớp 1-C đã biến ngày lễ tình nhân thành một bản giao hưởng mà nhạc công duy nhất nam đang giữ phách chờ nhạc trưởng giơ tay.

Tôi bước ra hành lang, tiếng giày vang lộp cộp trên nền gạch lạnh. Ánh đèn huỳnh quang chiếu xuống dãy tủ sắt, loáng thoáng bóng học sinh thò tay vào khe cửa rồi lại rút ra, như những chú cá nhỏ mút mát mồi. Tôi dừng trước tủ giày của mình và khẽ nắm tay nắm kim loại lạnh ngắt.

``\dots Chắc cũng kha khá đây.''

Tiếng lẩm bẩm vừa thoát khỏi môi, tôi đã nghe thêm tiếng thở dài của chính mình.  Cửa tủ kẹt một tiếng ken két, một mùi kim loại bỗng phả ra. Ngay lập tức, một luồng giấy màu hồng phấn, xanh nhạt, tím oải hương ùa ra như thác nước mùa xuân.  Chúng va vào tay tôi, rơi xuống tràn cả lên nền gạch.

``\dots Chuẩn bài luôn.''

Tôi đứng chôn chân, nhìn đống thư tình đang nhúc nhích như muốn bò đi. Mỗi tờ gấp nếp vuông vắn, mỗi tờ xoắn trái tim, mỗi tờ dán thêm kẹo hồng lấp lánh. Tên tôi - Hikawa Kuon - được viết bằng mực dạ quang, có chỗ nắn nót như thư pháp, có chỗ nguệch ngoạc như vừa run tay. Tất cả đều đến từ lớp 1-C, lớp mà chỉ toàn là nữ.

Ngay lúc tôi còn đang lúng túng ôm đống thư, một nhóm học sinh lớp trên đi ngang qua. Cô bạn tóc tém có mái bằng đeo kính tròn chợt dừng bước, mắt trố tròn: ``Ôi trời, nhìn kìa! Cậu ta còn nhiều thư hơn cả hộp thư tình trường mình tổng hợp!'' Cậu bạn đứng cạnh cô - người vừa bị từ chối hồi sáng - đập vai anh bạn thân, cười trừ: ``Thôi xong, đừng so sánh nữa, em ấy được quá trời người hâm mộ kìa.'' Nhưng mắt cậu vẫn dán vào đống giấy màu, như thể muốn đếm xem có bao nhiêu lá.

Một chị lớp năm ba ở gần đó bất ngờ rút điện thoại ra, lia máy quay về phía tôi: ``Các cậu ơi, đây chính là hiện tượng! Nam sinh duy nhất ở lớp toàn nữ nhận được cả rổ thư Valentine!'' Cô bạn cùng lớp cô ấy che miệng cười khúc khích: ``Chắc cậu ta có bùa yêu đấy!'' rồi cả hai ré lên khi thấy một lá thư trên tủ rơi xuống, va vào giày tôi.

Ở cuối hành lang, ba anh bạn cầu thủ bóng rổ lớp 2-C đang đứng tựa tường. Người cao nhất - đội trưởng - huýt sáo một tiếng dài, nhìn theo tôi với vẻ ngưỡng mộ xen lẫn ghen tị: ``Đệ nhất đào hoa năm nay chắc thuộc về cậu ta rồi.'' Cậu bạn gầy gò bên cạnh gãi đầu: ``Mình còn chưa nhận được lá nào\dots\ Hay là sang xin chữ ký của cậu ta cho may mắn?'' Người còn lại thì cười cười, đập nhẹ vào vai đội trưởng: ``Thôi, về tập thể lực đi, đừng có mà ngắm gái!''

Một nữ sinh lớp 1-D reo lên khi thấy tôi, rồi quay sang bạn cùng lớp thì thào: ``Kìa, anh rể tương lai kìa! Nghe nói anh ấy vừa đọc thư tình cả lớp 1-C luôn đó!'' Cô bạn cùng lớp cô ấy mắt tròn xoe: ``Thật hả? Thế thì chắc cả trường này không ai vượt qua được anh ấy mất!'' Cả hai cười rúc rích, rồi chạy vội về lớp như sợ tôi nghe thấy.

\dots Thật sự luôn, từ khi nào tôi trở thành một người ``nổi tiếng bất đắc dĩ'' tới mức này? Tôi cũng không hiểu tại sao một mình tôi lại được nhận từng ấy thư đấy.

Tôi cúi nhặt hết lên, chợt thấy mình đang giữ một bản đồ mảnh giấy. Mỗi một bức thư là một câu chuyện nho nhỏ hé lộ một góc nào đó của các bông hồng lớp tôi.

``Nhiều thế này cơ à\dots''

Đứng trước dãy tủ giày, tôi đảo mắt nhìn qua vô vàn những lá thư sặc sỡ màu sắc. Có những bức mang phong cách thể thao, có những bức trang trí họa tiết tinh xảo.

Tôi nuốt khan, cảm giác như vừa bước vào trò chơi hẹn hò mà người chơi chỉ có mình tôi. Nhưng đồng hồ chỉ còn một tiếng rưỡi trước khi mặt trời tắt nắng trên hành lang. Tôi không thể tách mình làm mười, cũng không thể phi thân đến mọi nơi được.

Vậy\dots\ tôi nên chọn bức thư nào đây?

\begin{table}[H]
  \label{choice}
  \centering
  \begin{tabular}{|>{\centering\arraybackslash}p{0.23\linewidth}|>{\centering\arraybackslash}p{0.23\linewidth}|>{\centering\arraybackslash}p{0.23\linewidth}|>{\centering\arraybackslash}p{0.23\linewidth}|}\hline
    \hyperref[shino]{Bức thư màu xanh đậm, sticker gấu nhỏ}      & \hyperref[yae]{Bức thư xám bạc, có gắn một hình chữ thập}    & \hyperref[sakura]{Bức thư màu hồng nhạt, họa tiết anh đào}   & \hyperref[inari]{Bức thư trắng, sticker đuôi cáo tuyết}      \\\hline
    \hyperref[hanabi]{Bức thư màu cam, nơ đen hình mặt mèo}      & \hyperref[touka]{Bức thư màu xanh nhạt, đính hoa trắng}      & \hyperref[akane]{Bức thư màu nâu, đính dấu chân cún}         & \hyperref[kaoru]{Bức thư màu xanh cổ vịt, dán hình con vịt}  \\\hline
    \hyperref[uzuki]{Bức thư màu xanh tím nhạt, dán hình cà rốt} & \hyperref[kanade]{Bức thư trắng, cánh thiên thần}            & \hyperref[aku]{Bức thư hồng-xanh lam, dán tay cầm chơi game} & \hyperref[saki]{Bức thư màu xanh lá-hồng, dán hình bông hoa} \\\hline
    \hyperref[shiina]{Bức thư xám trắng, hình đầu sư tử bạc}     & \hyperref[azuki]{Bức thư hồng gradient, dán nốt nhạc}        & \hyperref[kana]{Bức thư tím đậm, sticker ác ma}              & \hyperref[suzune]{Bức thư màu cam vàng, dán hình hải cẩu}    \\\hline
    \hyperref[yuka]{Bức thư tím oải hương, dán hình cơm nắm}     & \hyperref[mari]{Bức thư màu đỏ thẫm, dán hình bảng màu}      & \hyperref[saya]{Bức thư màu vàng, dán hình chòm sao}         & \hyperref[shion]{Bức thư màu đen, dán hình mũ phù thủy}      \\\hline
    \hyperref[hotaru]{Bức thư màu cam-xanh đậm, dán nơ xanh}     & \hyperref[miku]{Bức thư màu đỏ-đen, dán hình con sói đen}    & \hyperref[ayawata]{Hai bức thư đính nhau}                    & \hyperref[mitsuki]{Bức thư màu tím-hồng, dán hình cục kẹo}   \\\hline
    \hyperref[honoka]{Bức thư đa sắc, dán chỏm tóc hồng-đen}     & \hyperref[koishin]{Bức thư vàng chanh, chi chít ruy băng đỏ} & \hyperref[nanase]{Bức thư xanh nước biển, họa tiết ca-rô}    & \hyperref[yuki]{Bức thư đen-xám, dán hình rô-bốt}            \\\hline
    \multicolumn{4}{|c|}{\hyperref[canon]{Không chọn một bức thư nào trong số này}}                                                                                                                                                                           \\ \hline
  \end{tabular}
\end{table}

% I. Misora Shino
\newpage
\label{shino}

Ánh mắt tôi lại vô tình dừng lại ở một bức thư nằm khiêm tốn ở một góc tủ. Một bức thư màu xanh nước biển đậm, góc phải bên dưới dán một chiếc sticker gấu nhỏ xíu trông có phần rụt rè nhưng lại vô cùng đáng yêu. Làm tôi liên tưởng tới linh vật của Sora - cũng là một con gấu nâu này.

Không hiểu sao, tôi lại thấy bức thư này rất đỗi quen thuộc. Cảm giác e ấp, cẩn trọng ấy khiến tôi nhớ đến một bóng hình mỏng manh nhưng từng ôm trong mình một sức mạnh lớn lao vô ngần.

Tôi khẽ bóc lớp niêm phong. Bên trong chỉ có một dòng chữ nắn nót, viết bằng mực xanh nhạt:

\txtbx{``Sau giờ học, mình đợi cậu ở cổng trường nhé.''\\
  \begin{flushright}-Misora Shino-\end{flushright}}

Tôi mỉm cười, gập bức thư lại cẩn thận rồi cho vào túi áo. Cầm theo phần sô-cô-la của mình, tôi rảo bước qua hành lang trường đang dần vắng bóng học sinh, hướng thẳng về phía cổng chính.

\ct

\pov{Misora Shino}

Gió chiều tháng Năm mơn man lướt qua những tán lá xà cừ trước cổng trường Aogami. Tôi đứng nép mình sau một trụ đá lớn, hai tay nắm chặt lấy hộp sô-cô-la được gói ghém cẩn thận. Trái tim tôi đập thình thịch liên hồi, cứ như thể nó sắp sửa nhảy vọt ra khỏi lồng ngực.

``\textit{Phải làm sao đây\dots}''

Tôi thầm nghĩ, khẽ cắn môi. ``\textit{Dù đã cố gắng lấy hết can đảm để mời cậu ấy, nhưng giờ phút này mình vẫn không đủ tự tin để trực tiếp trao nó\dots}''

Thói quen rụt rè, thiếu tự tin của một đứa con gái nhút nhát hồi trước lại bắt đầu trỗi dậy. Tôi vốn dĩ luôn tự ti, luôn ngưỡng mộ những người nổi bật trong lớp, và thường hay tự trách bản thân vì không thể nào bắt chuyện với họ một cách tự nhiên. Dù đã có những người bạn tốt ở bên cạnh, nhưng để chủ động bày tỏ tình cảm với một người con trai, đó vẫn là một rào cản quá lớn đối với tôi.

Nhưng rồi, một mảnh ký ức chợt xẹt qua tâm trí tôi.

Đó là khoảnh khắc trong miền ký ức ở trường Trung học Aogami cũ, khi tôi đối mặt với ảo ảnh của chính mình, với những bóng ma quá khứ cố gắng níu kéo tôi lại trong vỏ bọc an toàn. Khi đó, tôi đã không chọn trốn tránh. Tôi đã tự mình dùng mảnh kính cắt đứt vòng lặp hư ảo đó, đánh đổi cả sinh mạng và dùng chính sức mạnh của bản thân để gieo hạt giống, cứu lấy thị trấn này, cứu lấy bạn bè mình. Lần đó, tôi đã làm được. Tôi đã vượt qua giới hạn của sự sợ hãi.

``\textit{Nhưng nếu không trao cho cậu ấy như thế này, chắc chắn mình sẽ hối hận.}''

Trận chiến ngày ấy khốc liệt đến thế tôi còn dám bước tới, cớ sao bây giờ tôi lại lùi bước chỉ vì một lời thổ lộ? Chỉ mong mỏi, chỉ ước ao thôi thì chưa đủ. Tôi phải tự mình hành động. Vì tôi muốn thay đổi bản thân yếu đuối này. Vì\dots\ tôi muốn được gần gũi với cậu ấy hơn nữa.

Tiếng bước chân lộp cộp vang lên từ phía trong sân trường. Bóng dáng quen thuộc của Kuon dần hiện ra qua khung cổng sắt. Cậu ấy đang đi về phía này.

Cơ thể tôi khẽ run rẩy. Tay tôi siết chặt lấy hộp quà. ``\textit{Hành động đi, Shino. Chính cậu đã từng cứu cả Aogami cơ mà!}''

Cậu con trai bước qua cổng, đôi mắt lướt tìm xung quanh. Ngay khi cậu ấy chuẩn bị quay bước, tôi nhắm nghiền mắt, dồn hết sức lực và cất tiếng gọi:

``\textbf{N-Này, đợi đã!}''

\begin{wrapfigure}[17]{r}{7cm}
  \includegraphics[width=7cm]{../../../../Image/Book?/valentine/shino.png}
\end{wrapfigure}

Cậu ấy giật mình, lập tức dừng bước và quay lại. Khi thấy tôi đang đứng thở dốc, hai gò má ửng đỏ, cậu ấy khẽ nở một nụ cười hiền.

``Tôi đến đúng điểm hẹn rồi chứ?'' Kuon cất giọng trầm ấm, tay khẽ lấy từ trong túi ra bức thư màu xanh nước biển.

Tôi hít một hơi thật sâu, cố ngăn cho giọng mình khỏi run, hai tay run rẩy đưa chiếc hộp sô-cô-la truffle nhỏ nhắn về phía cậu.

``C-Cái này\dots\ Mình làm sô-cô-la này chỉ dành riêng cho cậu\dots\ Mình muốn cậu nhận nó.''

Kuon không nói gì, chỉ nhẹ nhàng đón lấy chiếc hộp từ tay tôi. Ngón tay cậu khẽ chạm vào tay tôi, mang theo một hơi ấm kỳ lạ khiến mặt tôi càng lúc càng nóng bừng. Cậu cẩn thận mở hộp, nhìn những viên nấm truffle nhỏ xinh được tôi tỉ mỉ nắn nót từng chút một.

``Cảm ơn cậu, Shino-san. Nó trông rất đẹp, chắc chắn cậu đã bỏ rất nhiều tâm huyết vào đây.'' Cậu khẽ mỉm cười, rồi đưa tay vào túi xách, lấy ra một chiếc hộp nhỏ hình chiếc lá mùa thu. ``Đây\dots\ là phần của tôi. Dù không biết có hợp khẩu vị của cậu không, nhưng hy vọng cậu sẽ thích.''

Tôi ngỡ ngàng nhận lấy món quà. Một chiếc sô-cô-la mang hình dáng chiếc lá, giản dị nhưng lại tinh tế vô cùng, hệt như chính con người của cậu ấy vậy.

``M-Mình có thể thử bây giờ được không?'' tôi rụt rè hỏi.

``Tất nhiên rồi.''

Tôi khẽ đưa viên sô-cô-la lên miệng cắn một miếng nhỏ. Vị cacao đắng nhẹ hòa quyện cùng chút ngọt ngào của đường nâu và hương vani lan tỏa nơi đầu lưỡi. Nhưng ngọt ngào nhất có lẽ là hơi ấm lan tràn trong lồng ngực tôi lúc này. Tôi nhắm mắt, mỉm cười hạnh phúc, hai gò má đã đỏ lựng từ bao giờ.

``Ngon lắm\dots\ Kuon-san à.'' Giọng tôi lí nhí nhưng chứa đựng sự chân thành từ tận đáy lòng. ``Thật sự rất ngon.''

Kuon cũng lấy một viên sô-cô-la truffle của tôi cho vào miệng, rồi khẽ gật gù hài lòng. ``Của cậu cũng vậy. Vị ngọt rất dịu, giống hệt như cậu vậy, Shino-san.''

Cơn gió chiều mơn man vuốt ve mái tóc tôi, cuốn theo những cánh hoa mỏng manh bay lên bầu trời trong xanh. Trong khoảnh khắc ấy, tôi biết rằng nỗ lực thay đổi của bản thân đã không hề uổng phí. Tôi không còn là một bóng ma cô độc trong vòng lặp thời gian, cũng không còn là một cô gái nhút nhát chỉ biết đứng nhìn từ xa nữa. Tôi đã tự tay nắm lấy hạnh phúc của chính mình.

Ngày Lễ tình nhân này, được cậu ấy ``chọn'' và được đứng cạnh cậu ấy thế này\dots\ quả thực là phép màu tuyệt vời nhất.

\begin{table}[H]
  \centering
  \setlength{\arrayrulewidth}{1pt}
  \begin{tabular}{|>{\centering\arraybackslash}p{0.4\linewidth}|>{\centering\arraybackslash}p{0.4\linewidth}|}
    \hline
    \hyperref[choice]{Quay trở về lựa chọn} & \hyperref[ending]{Đi tới phần cuối} \\ \hline
  \end{tabular}
\end{table}

% II. Shiromi Yae
\newpage
\label{yae}

Tôi cầm bức thư màu xám bạc lên. Điểm nhấn của nó là một hình chữ thập nhỏ được dập nổi cẩn thận ở giữa nắp phong bì, mang lại cảm giác năng động và gọn gàng, hoàn toàn không giống với những bức thư tình thông thường.

Khẽ mở lớp niêm phong, tôi thấy dòng chữ viết tay dứt khoát nhưng cũng không kém phần uyển chuyển:

\txtbx{``Gặp mình ở hành lang dãy phòng học năm nhất nhé.''\\
  \begin{flushright}-Shiromi Yae-\end{flushright}}

Không chút chần chừ, tôi đút thư vào túi, cầm lấy phần sô-cô-la của mình và tiến về phía hành lang được chỉ định.

\ct

\pov{Shiromi Yae}

Tôi dựa lưng vào tường hành lang dãy phòng học năm nhất, đôi chân khẽ đá nhẹ vào không khí. Trong tay, tôi giữ khư khư hộp sô-cô-la hạnh nhân mà mình đã cất công làm từ sáng ngày hôm nay.

Bình thường, tôi tự tin mình là một người năng động, mạnh mẽ, chẳng bao giờ chùn bước trước bất kỳ chuyện gì, kể cả trong những môn thể thao đòi hỏi sức lực. Nhưng sao lúc này, chỉ việc đợi một người con trai tới để tặng quà lại khiến tim tôi đập loạn nhịp thế này? Có phải bởi lúc mà cậu ấy đã đồng hành cùng tôi trong hội thao vừa rồi, dù là một người không hề thích thể chát, nhưng Kuon-san đã thể hiện một màn trình diễn tuyệt vời khiến tôi sững sờ không?

Cũng có thể lắm chứ.

``\textit{Chỉ cần đưa cho cậu ấy rồi nói: `Này, sô-cô-la này, ăn đi!'}'' Tôi nhẩm lại kịch bản trong đầu. ``\textit{Không, nghe thô quá. Hay là: `Kuon-san, đây là thành quả của nghệ thuật nấu ăn đỉnh cao của tớ!'}''

Lại không ổn. Cứ thế này thì đến lúc gặp, tôi sẽ cắn lưỡi vì bối rối mất. Nanase mà thấy cảnh này, chắc cậu ấy sẽ lại buông lời mỉa mai tôi cho xem.

Đang lúc tôi còn vò đầu bứt tai, tiếng bước chân nhịp nhàng vang lên. Kuon xuất hiện ở đầu hành lang, vẫn với dáng vẻ điềm tĩnh thường ngày. Cậu ấy đưa mắt nhìn quanh rồi bước thẳng về phía tôi.

``Ồ. Yae-san.''

Tôi hít một hơi thật sâu, đứng thẳng người dậy, cố tỏ ra tự nhiên nhất có thể.

``A, Kuon-san! Cậu tới rồi.'' Tôi gật đầu chào, giấu túi sô-cô-la ra sau lưng.

``Cậu hẹn tôi ra đây, có chuyện gì sao, Yae-san?'' cậu ấy hỏi, giọng nhẹ nhàng.

``Ừm\dots\ thì\dots'' Tôi đảo mắt đi nơi khác, bắt đầu kịch bản lảng tránh của mình. ``Cậu\dots\ đã mua sô-cô-la ở đâu chưa?''

``Mua sao?'' cậu ấy có vẻ hơi ngạc nhiên. ``Hôm nay ở trường mọi người tự làm nhiều lắm, nên tôi không nghĩ mình cần phải ra ngoài mua đâu.'' Nghĩ lại, đó có lẽ là một câu hỏi khá ngớ ngẩn.

``Vậy à\dots\ thế thì tốt.'' Tôi hắng giọng, đưa chiếc túi từ sau lưng ra trước mặt cậu ấy. ``Nếu cậu chưa có ai tặng, à nhầm, nếu cậu chưa mua, vậy cậu có muốn mua phần của tớ không? Tớ thích nấu ăn lắm, nên đảm bảo chất lượng của nó không đùa được đâu. Chắc chắn sẽ ngon lắm đấy!''

Kuon nhận lấy túi sô-cô-la, khóe môi hơi nhếch lên. ``Ra vậy. Cảm ơn cậu nhé, Yae-san. Sô-cô-la này\dots\ là sô-cô-la tình yêu (honmei) sao?''

Câu hỏi thẳng thắn của cậu ấy như một tiếng sấm rền giữa trời quang. Mặt tôi nóng bừng bừng, tôi lúng túng huơ huơ hai tay lên không trung.

\begin{wrapfigure}[18]{l}{7cm}
  \includegraphics[width=7cm]{../../../../Image/Book?/valentine/yae.png}
  \caption{(Giả sử cô nàng đang diện đồng phục Aogami thời hiện đại.)}
\end{wrapfigure}

``H-Hả? Tình yêu đích thực gì chứ!? Cậu đừng có mà suy diễn lung tung vậy chứ!'' tôi hốt hoảng, vội vã tìm một chủ đề khác để lấp liếm sự bối rối này. ``À\dots\ à mà này, nghe tớ hỏi đây! Shiina-san lại đi đánh nhau với mấy đứa trường khác nữa rồi phải không? Thật tình, chuyện này chỉ làm tôi lo lắng mãi thôi! Cậu ấy lúc nào cũng thích dùng nắm đấm giải quyết vấn đề\dots''

Kuon không đáp lời ngay, cậu ấy chỉ nhìn tôi, rồi bật cười thành tiếng.

``C-Cậu cười cái gì thế!?'' tôi phồng má, cảm thấy như thể bị nhìn thấu hết ruột gan.

``Không có gì. Chỉ là tôi thấy cậu lúc lúng túng trông cũng rất đáng yêu.'' Kuon đáp, đoạn đưa cho tôi một hộp quà nhỏ. ``Của tôi đây. Hy vọng một người sành ăn như cậu sẽ không chê.''

Tôi khẽ chớp mắt, đưa tay nhận lấy. ``Cậu\dots\ cậu cũng làm sao?''

``Dĩ nhiên. Cậu thử đi.''

Tôi cẩn thận mở hộp. Bên trong là viên sô-cô-la hình chiếc lá được đúc rất mượt mà. Tôi cắn một miếng nhỏ. Hương vị cacao lan tỏa, có chút đắng dịu nhưng hậu vị lại ngọt ngào, cân bằng hoàn hảo.

Mắt tôi sáng lên. Mọi sự ngại ngùng ban nãy dường như bay biến hết. Tôi mỉm cười sảng khoái, vỗ cái đét lên vai Kuon.

``Tuyệt thật đấy, Kuon-san! Không ngờ tài nấu nướng của cậu cũng khá đến thế. Trông thế mà khéo tay phết nhỉ!''

Kuon khẽ xoa vai, cười trừ: ``Cảm ơn vì lời khen. Lần này tôi làm thế này mà được như vậy cũng tốt rồi.''

Tôi mở hộp của mình ra, lấy một viên sô-cô-la hạnh nhân bọc đường đưa cho cậu ấy. ``Nào, cậu cũng thử của tớ đi!''

Cậu ấy nhấm nháp thử, rồi gật đầu vẻ tán thưởng. ``Ngon thật. Hạnh nhân rất giòn, sô-cô-la lại không quá ngọt. Đúng là sản phẩm của người đam mê nấu ăn có khác.''

Nghe những lời khen thật lòng của cậu ấy, lòng tôi trào dâng một cảm giác ấm áp khó tả. Tôi khẽ cười, nhìn ra ngoài cửa sổ hành lang.

Dù ban đầu tôi đã cố gắng lấp liếm bằng những câu chuyện không đâu, nhưng dường như Kuon đã hiểu rõ tâm ý của tôi từ trước. Đứng bên cạnh cậu ấy lúc này, tôi cảm thấy sự năng động thường ngày lại trở về, nhưng lần này nó được tiếp thêm bởi một niềm vui ngọt ngào chưa từng có.

Được cùng cậu ấy chia sẻ khoảnh khắc Lễ tình nhân này\dots\ quả thực là một điều tuyệt vời.

\begin{table}[H]
  \centering
  \setlength{\arrayrulewidth}{1pt}
  \begin{tabular}{|>{\centering\arraybackslash}p{0.4\linewidth}|>{\centering\arraybackslash}p{0.4\linewidth}|}
    \hline
    \hyperref[choice]{Quay trở về lựa chọn} & \hyperref[ending]{Đi tới phần cuối} \\ \hline
  \end{tabular}
\end{table}

% III. Shinomiya Sakura
\newpage
\label{sakura}

Ánh mắt tôi dừng lại ở một phong thư màu hồng nhạt vô cùng nổi bật, trên góc còn điểm xuyết họa tiết những cánh hoa anh đào đang bay lơ lửng. Màu sắc dịu dàng và nét trang trí mang đậm chất truyền thống ấy nhanh chóng giúp tôi nhận ra chủ nhân của nó.

Tôi khẽ mở nắp phong bì, bên trong là một tấm thiệp nhỏ nhắn với những dòng chữ mềm mại, cẩn thận:

\txtbx{``Sau giờ học, hãy đến bậc thang dẫn lên đền Aogami ngay.''\\
  \begin{flushright}-Shinomiya Sakura-\end{flushright}}

Câu văn của cô nàng có phần hơi áp đặt - đúng tính cách của Sakura-san. Khẽ mỉm cười, tôi cất bức thư vào túi áo và cầm lấy phần sô-cô-la của mình, hướng bước chân về phía con đường mòn quen thuộc dẫn lên đền thờ.

\ct

\pov{Shinomiya Sakura}

Ánh tà dương nhuộm đỏ cả một góc trời, rọi những vệt nắng xiên qua từng kẽ lá xà cừ, rớt xuống những bậc thang đá rêu phong dẫn lên đền Aogami. Tôi ngồi ôm gối trên một bậc thềm ở giữa đoạn đường, chiếc túi giấy đựng hộp sô-cô-la đặt ngay ngắn bên cạnh.

Hít một hơi thật sâu không khí trong lành của buổi chiều tà, tôi đưa mắt nhìn xuống con đường dốc. Mặc dù chỗ này hơi vắng vẻ, nhưng tôi lại thích sự tĩnh lặng của nó. Nó mang đến cho tôi cảm giác yên bình, khác hẳn với những tiếng ồn ào ở sân trường.

Nhưng sâu thẳm trong lòng, sự yên bình ấy lại đang bị khuấy động bởi những nhịp đập rộn ràng của trái tim.

Đợi người con trai mình thích đến để tặng sô-cô-la trong ngày Valentine\dots\ Cảm giác này thật lạ lẫm, nhưng cũng thật tuyệt vời. Nhớ lại khoảng thời gian trước kia, tâm trí tôi đã từng bị giam cầm trong những ảo ảnh hư vô, bị bó buộc bởi danh xưng ``Miko'' và những sứ mệnh vô hình. Chính Kuon-san đã dùng đôi tay ấy để kéo tôi thoát khỏi vũng lầy, đưa tôi trở về với hiện thực sống động này, giúp tôi có thể trải qua những ngày tháng bình dị như bao nữ sinh trung học khác.

Và giờ đây, tôi đang ngồi ở chính nơi này, chờ đợi cậu ấy với tư cách là một thiếu nữ bình thường mang trong mình một trái tim biết rung động.

``Sakura-san.''

Một giọng nói trầm ấm vang lên cắt ngang dòng suy nghĩ của tôi. Từ dưới chân dốc, Kuon đang chậm rãi bước lên, trên môi thoáng một nụ cười quen thuộc.

``A\dots\ Kuon-san!'' Tôi giật mình đứng phắt dậy, vội vàng chải lại những nếp nhăn trên váy.

Kuon dừng lại ở bậc thang ngay dưới chỗ tôi đứng, khẽ gật đầu chào: ``Xin lỗi vì đã để cậu phải đợi. Ở đây lâu chưa?''

``Kh-Không sao đâu! Tôi cũng vừa mới tới thôi!'' tôi vội xua tay, dù thực ra đã ngồi ở đây cả nửa tiếng đồng hồ rồi. Bầu không khí bỗng chốc trở nên tĩnh lặng, chỉ còn nghe thấy tiếng lá cây xào xạc theo gió. Tôi bối rối đảo mắt xung quanh, rồi quyết định sử dụng ``kịch bản'' mà mình đã dày công chuẩn bị.

``À\dots\ C-Cậu nghe chuyện này chưa\dots?'' tôi ngập ngừng lên tiếng, hai tay giấu sau lưng đan chặt vào nhau.

``Gì vậy?''

``Hình như\dots\ có lời đồn là nếu cậu ăn sô-cô-la do một thiếu nữ tự làm thì sẽ được bình an, không bệnh tật tai ương gì đó\dots''

Cậu ta hơi nhướng mày, ánh mắt mang theo vài phần tò mò xen lẫn ý cười: ``Ồ, nghe thú vị đấy. Thế cậu nghe được điều đó từ đâu vậy, Sakura-san?''

\begin{wrapfigure}[18]{r}{7cm}
  \includegraphics[width=7cm]{../../../../Image/Book?/valentine/sakura.png}
  \caption{(Giả sử cô nàng đang diện đồng phục Aogami thời hiện đại.)}
\end{wrapfigure}

Câu hỏi vặn lại của cậu ấy khiến tim tôi lỡ mất một nhịp. ``C-Cái đó\dots\ Bí mật! Cậu tuyệt đối không được nói cho ai biết đâu đấy!'' Tôi luống cuống đáp lời, hai gò má đã nóng ran như lửa đốt. Chắc chắn là cậu ấy đã nhìn thấu cái cớ vụng về này của tôi rồi.

Dù vậy, tôi vẫn thu hết can đảm, đưa chiếc hộp đựng hộp sô-cô-la ra trước mặt cậu ấy. ``V-Vậy nên\dots\ tôi cũng làm cho cậu một ít! Không, ý tôi là\dots\ tôi không chắc về cái lời đồn kia đâu, nhưng mà\dots\ cậu nhận nhé?''

Kuon nhìn hộp sô-cô-la được thắt nơ cẩn thận, rồi khẽ mỉm cười đưa tay nhận lấy. Ngón tay cậu sượt nhẹ qua tay tôi, mang theo một hơi ấm dịu dàng.

``Tất nhiên là tôi sẽ ăn rồi. Cảm ơn cậu nhé, Sakura-san.''

Tôi ngước mắt nhìn cậu, ngỡ ngàng. ``Ể-Ể? Cậu sẽ ăn sao? V-Vậy à\dots\ Cảm ơn cậu!'' Niềm vui sướng trào dâng khiến tôi không kìm được mà bật cười khúc khích.

Kuon mở chiếc hộp ra, cẩn thận nhón lấy một viên sô-cô-la nhỏ bỏ vào miệng. Cậu ấy nhai chậm rãi, rồi gật đầu vẻ hài lòng. ``Ngon lắm. Vị ngọt rất vừa vặn. Chắc chắn năm nay tôi sẽ được bình an vô sự nhờ phúc của `thiếu nữ' đây rồi.''

Cậu ấy mỉm cười trêu chọc khiến mặt tôi càng đỏ hơn, nhưng trong lòng lại ngọt ngào vô cùng. ``Mồ\dots\ cậu lại trêu tôi nữa!''

Kuon bật cười, rồi cũng lấy từ trong túi áo ra một hộp quà nhỏ. ``Đây là phần của tôi. Dù không có `phép màu' như của cậu, nhưng hy vọng nó sẽ mang lại chút niềm vui.''

Tôi vui vẻ nhận lấy hộp sô-cô-la hình chiếc lá từ tay cậu. ``Wow\dots\ Dễ thương quá! Tôi ăn thử luôn nhé?''

Kuon gật đầu. Tôi khẽ cắn một miếng. Hương vị cacao nồng nàn quyện cùng sự ngọt ngào tan chảy nơi đầu lưỡi. Thật sự rất tinh tế, hệt như sự dịu dàng mà cậu ấy luôn dành cho tôi.

``Ngon tuyệt vời luôn! Kuon-san khéo tay thật đấy\~{}'' tôi hớn hở khen ngợi.

``Cũng thường thôi. Cậu cứ quá lời làm gì\dots''

Chúng tôi cùng nhau ngồi lại trên bậc thềm đá, nhâm nhi những viên sô-cô-la ngọt ngào dưới ánh hoàng hôn rực rỡ. Gió chiều thổi qua mang theo hương hoa cỏ thoang thoảng. Lặng lẽ ngắm nhìn sườn mặt góc cạnh của Kuon, tôi lại thầm cảm ơn cậu. Cảm ơn vì đã đưa tôi trở về, cảm ơn vì đã cho tôi được trải qua những cung bậc cảm xúc tuyệt vời này.

Lễ tình nhân năm nay, đối với tôi, chính là ngày hạnh phúc nhất.

\begin{table}[H]
  \centering
  \setlength{\arrayrulewidth}{1pt}
  \begin{tabular}{|>{\centering\arraybackslash}p{0.4\linewidth}|>{\centering\arraybackslash}p{0.4\linewidth}|}
    \hline
    \hyperref[choice]{Quay trở về lựa chọn} & \hyperref[ending]{Đi tới phần cuối} \\ \hline
  \end{tabular}
\end{table}

% IV. Shirayuki Inari
\newpage
\label{inari}

Tôi chọn lấy một phong bì màu trắng tinh khôi. Nổi bật trên nền giấy trắng ấy là một hình dán bé xíu mang hình chiếc đuôi cáo tuyết xen lẫn hai màu trắng đen, được điểm xuyết thêm một ngôi sao nhỏ màu trắng nhắn ở góc.

Tôi khẽ bóc lớp niêm phong, bên trong là nét chữ ngay ngắn, mềm mại và nắn nót, không thể lẫn vào đâu:

\txtbx{``Tan học, cậu đợi mình ở hành lang dãy phòng học năm nhất nhé. Cùng đi về nào!''\\
  \begin{flushright}-Shirayuki Inari-\end{flushright}}

Cất bức thư vào túi, tôi mỉm cười và hướng bước chân về phía điểm hẹn.

\ct

\pov{Shirayuki Inari}

Ánh hoàng hôn nhuộm vàng hành lang dãy phòng học năm nhất, hắt những cái bóng dài in trên mặt sàn gỗ. Tôi khẽ siết chặt chiếc hộp trong tay, nhịp tim đập thình thịch như muốn nhảy vọt ra ngoài.

Dù chính mình là người đã rủ Kuon-kun cùng đi về, nhưng rốt cuộc tôi lại phải để cậu ấy chờ vì cuộc họp toàn thể ban chấp hành lớp kéo dài hơn dự kiến. Lúc nãy tôi đã phải chạy thục mạng từ phòng họp đến đây, chỉ sợ rằng cậu ấy đã không kiên nhẫn nổi mà bỏ về trước.

Nhưng khi thấy bóng dáng quen thuộc đang đứng tựa lưng vào tường nơi góc hành lang, lòng tôi bỗng nhẹ nhõm hẳn đi. Ơn trời, cậu ấy vẫn ở đây!

``Kuon-kun!'' tôi cất tiếng gọi, bước vội đến chỗ cậu ấy, hơi thở vẫn còn đứt quãng. ``X-Xin lỗi vì đã để cậu đợi!''

Cậu ấy quay sang, mỉm cười dịu dàng, khẽ lắc đầu: ``Không sao đâu. Tôi cũng chỉ vừa tới thôi. Có chuyện gì mà cậu lại gấp gáp thế?''

``M-Mình xin lỗi vì chính mình là người rủ cậu về cùng mà lại bắt cậu chờ\dots\ Buổi họp ban chấp hành lớp kéo dài hơn dự kiến một chút\dots'' tôi cúi gầm mặt, hai tay giấu chiếc hộp ra sau lưng.

Kuon-kun chỉ cười nhẹ, không hề tỏ ra khó chịu. Bầu không khí xung quanh bỗng chốc trở nên tĩnh lặng. Giờ này học sinh trong trường đã về gần hết, hành lang vắng lặng chỉ còn tiếng gió lùa qua khe cửa sổ.

``Nhưng mà\dots\ ở đây không có ai cả, nên có lẽ mình có thể nói ra chuyện này cho đỡ hồi hộp hơn\dots'' tôi hít một hơi thật sâu, lấy hết can đảm ngước lên nhìn cậu ấy.

``Chuyện gì vậy?'' cậu ấy hơi nghiêng đầu, ánh mắt mang theo vẻ tò mò.

``À, thì\dots\ cậu thấy đấy\dots\ ừm\dots'' tôi lắp bắp, cảm giác như hai má đang nóng rực lên. ``Cậu\dots\ cậu có thích đồ ngọt không?''

Kuon-kun hơi ngớ người ra một lúc, rồi khẽ gật đầu. ``Bình thường. Sao vậy?''

Nghe vậy, tôi mừng rỡ như bắt được vàng, vội vàng đưa chiếc hộp sô-cô-la được thắt nơ đỏ xinh xắn ra trước mặt cậu ấy.

``M-Mình đã vất vả làm chúng đấy! X-Xin cậu hãy nhận món quà này!'' tôi nhắm tịt mắt lại, giơ hai tay ra phía trước, tim đập thót lên từng nhịp.

\begin{wrapfigure}[17]{l}{7cm}
  \includegraphics[width=7cm]{../../../../Image/Book?/valentine/inari.png}
\end{wrapfigure}

Một khoảng lặng ngắn ngủi trôi qua, rồi tôi cảm nhận được một bàn tay ấm áp nhận lấy chiếc hộp từ tay mình. Mở mắt ra, tôi thấy Kuon-kun đang nhìn chiếc hộp sô-cô-la với một nụ cười rạng rỡ.

``Cảm ơn cậu, Inari-san. Nó được bọc rất đẹp, chắc hẳn cậu đã dành nhiều tâm huyết cho nó lắm. Tôi nhất định sẽ ăn thật ngon miệng.''

Nghe những lời chân thành ấy, mọi sự lo lắng và ngượng ngùng trong tôi dường như tan biến hết. ``Ư-Ừm! Cậu nhớ ăn hết nhé!''

Cậu ấy gật đầu, rồi thò tay vào túi áo, lấy ra một chiếc hộp nhỏ hình chiếc lá mùa thu đưa cho tôi. ``Đây là phần của tôi. Dù không biết có hợp khẩu vị của cậu không, nhưng hy vọng cậu sẽ thích.''

Tôi tròn mắt ngạc nhiên, hai tay đón lấy hộp sô-cô-la. ``Của\dots\ của cậu làm sao? Mình có thể thử luôn được không?''

``Tất nhiên rồi.''

Tôi khẽ mở hộp, lấy ra một viên sô-cô-la nhỏ nhắn bỏ vào miệng. Vị đắng nhẹ của cacao hòa quyện cùng vị ngọt ngào thanh tao dần tan ra nơi đầu lưỡi. Thật sự rất ngon!

``Ngon quá!'' tôi kêu lên, ánh mắt sáng rỡ nhìn Kuon-kun. ``Cậu khéo tay thật đấy!''

Cậu ấy bật cười, xoa đầu tôi nhè nhẹ. ``Cảm ơn cậu. Cậu thích là tốt rồi.''

Cảm nhận hơi ấm từ bàn tay cậu ấy truyền qua đỉnh đầu, mặt tôi lại một lần nữa đỏ bừng lên. Lặng lẽ sánh bước bên cạnh cậu ấy, chúng tôi cùng nhau rời khỏi dãy hành lang ngập tràn ánh chiều tà.

Năm nay, có lẽ tôi đã nhận được một món quà Lễ tình nhân tuyệt vời nhất rồi.

\begin{table}[H]
  \centering
  \setlength{\arrayrulewidth}{1pt}
  \begin{tabular}{|>{\centering\arraybackslash}p{0.4\linewidth}|>{\centering\arraybackslash}p{0.4\linewidth}|}
    \hline
    \hyperref[choice]{Quay trở về lựa chọn} & \hyperref[ending]{Đi tới phần cuối} \\ \hline
  \end{tabular}
\end{table}

% V. Natsuno Hanabi
\newpage
\label{hanabi}

Tôi bị thu hút bởi một phong bì màu cam rực rỡ. Điểm nhấn của bức thư là một chiếc nơ đen viền xanh lá, ở giữa đính một hình mặt mèo trắng - giống hệt với chiếc dây buộc tóc mà cô bạn năng động nọ thường dùng.

Tôi khẽ mở thư. Bên trong là nét chữ to, tròn trịa và tràn đầy năng lượng:

\txtbx{``Sau giờ học, hãy lên sân thượng ở dãy A nhé! Chờ mình ở đó!''\\
  \begin{flushright}-Natsuno Hanabi-\end{flushright}}

Cất bức thư đi, tôi hướng lên tầng cao nhất của trường. Mở cánh cửa kim loại nặng trịch, những cơn gió hoàng hôn lùa qua khiến tôi cảm thấy dễ chịu. Trong lúc chờ đợi, tôi bước đến bên hàng rào lưới thép, tĩnh lặng ngắm nhìn khung cảnh mặt trời đang dần lặn xuống phía sau dãy núi xa.

\ct

\pov{Natsuno Hanabi}

``Này này! Cậu đứng một mình thẩn thơ trên sân thượng lúc hoàng hôn thế này, thanh xuân quá nhỉ?''

Tôi cất tiếng gọi lớn, phá vỡ bầu không gian yên tĩnh. Khi Kuon giật mình quay lại, cậu ấy thấy tôi đang ngồi xổm ngay cạnh bức tường gần lối ra vào, hai tay ôm khư khư một chiếc túi quà nhỏ được buộc dây ruy-băng đỏ. Tôi khẽ ngước lên nhìn cậu ấy, nở một nụ cười nghịch ngợm thường thấy.

\begin{wrapfigure}[17]{r}{7cm}
  \includegraphics[width=7cm]{../../../../Image/Book?/valentine/hanabi.png}
\end{wrapfigure}

Cậu ta bật cười, bước lại gần. ``Ra là cậu đã ở đây từ trước rồi sao? Cậu làm tôi giật mình đấy. Thế cậu làm gì mà lại ngồi xổm ở góc đó vậy?''

Tôi từ từ đứng dậy, phủi nhẹ bụi trên váy. ``Hahaha, đùa một chút thôi mà, cậu đừng có cáu chứ!''

``Cáu gì đâu. Chỉ là thấy hơi lạ khi một người lúc nào cũng ồn ào như cậu lại trốn trong góc tĩnh lặng thế này thôi.''

``Ể, mình á? Mình á\dots'' tôi hơi ngập ngừng, đưa mắt nhìn về phía những dải mây màu cam đỏ vắt ngang bầu trời. ``Đôi khi mình cũng muốn được thẩn thơ một mình trong ánh hoàng hôn chứ\dots''

Kuon hơi chau mày, ánh mắt lộ rõ vẻ lo lắng. ``Cậu ổn chứ, Hanabi-san? Cậu đang có gì muốn tâm sự sao? Sao trông cậu bồn chồn vậy?''

Thấy vẻ mặt nghiêm túc của cậu ấy, tôi không nhịn được mà phì cười. ``Nghe\dots\ hahaha\dots\ không giống mình chút nào nhỉ?'' tôi xua tay liên lịa. ``À, đùa thôi đùa thôi! Cảm ơn vì đã lo lắng cho mình nhé! Thực ra, mình chỉ lén đi theo cậu từ lúc ở tủ giày lên tận sân thượng này để đưa cho cậu cái này thôi!''

Nói rồi, tôi chìa chiếc túi quà buộc dây đỏ ra trước mặt Kuon. ``Này, nhận lấy đi nhé!''

Cậu ta chớp mắt ngạc nhiên, rồi mỉm cười đưa tay đón lấy chiếc túi từ tay tôi. ``Cảm ơn cậu, Hanabi-san. Cậu tự làm sao?''

``Dĩ nhiên rồi! Bên trong là sô-cô-la hạnh nhân dâu tây đó! Mình đã phải vật lộn với cái bếp suốt cả buổi sáng hôm nay đấy, nên hãy ăn thử và cho mình biết cảm nhận ngay đi!''

Kuon nhẹ nhàng mở chiếc túi, lấy ra một viên sô-cô-la màu hồng nhạt và cho vào miệng. Cậu ấy nhai chậm rãi, ánh mắt khẽ nhắm lại như đang thưởng thức. Tôi hồi hộp nhìn chằm chằm vào biểu cảm của cậu ấy, hai tay bất giác siết chặt lại vào nhau.

``Ngon lắm,'' cậu mắt ra, gật đầu tán thưởng. ``Sô-cô-la rất mềm, vị dâu tây chua chua ngọt ngọt kết hợp với độ giòn của hạnh nhân rất hoàn hảo. Cậu có năng khiếu nấu ăn đấy, Hanabi-san.''

Nghe thấy lời khen ngợi ấy, một niềm vui sướng khó tả trào dâng trong lồng ngực tôi. ``Thật sao? Thật sự rất ngon sao? Yay! Bõ công mình đã làm hỏng mất mấy mẻ đầu tiên!'' tôi nhảy cẫng lên ăn mừng.

Kuon cười nhẹ, rồi lấy từ trong túi áo ra một chiếc hộp nhỏ đưa cho tôi. ``Đây là phần của tôi. Dù không màu sắc như của cậu, nhưng hy vọng nó sẽ hợp khẩu vị.''

Tôi háo hức nhận lấy chiếc hộp, bên trong là những viên sô-cô-la hình chiếc lá được làm vô cùng tinh xảo. Tôi liền cho ngay một viên vào miệng. Vị đắng nhẹ của cacao ngay lập tức tan chảy, nhường chỗ cho vị ngọt hậu thanh tao vô cùng dễ chịu.

``Ưm! Ngon quá đi mất! Kuon, cậu mua ở tiệm nào mà đỉnh thế?'' tôi thốt lên, hai mắt sáng rực.

``Là tôi tự làm đấy. Sáng nay làm với cả lớp luôn.''

``Hả?! Thật luôn á? Đỉnh quá vậy!'' tôi đấm nhẹ vào vai cậu ấy đầy ngưỡng mộ. ``Cậu giỏi thật đấy! Lần sau phải dạy mình vài bí quyết làm sô-cô-la mới được!''

``Đừng có trông mong gì quá nhiều từ tôi,' cậu ta nhún vai đáp lại, ``dù gì đây cũng là lần tôi làm mà.'' Nhưng rồi cậu cũng gật đầu đồng ý.

Cơn gió sân thượng lùa qua làm tung bay mái tóc tôi. Dưới ánh tà dương rực rỡ, tôi lén nhìn Kuon đang nhâm nhi viên sô-cô-la của mình, trong lòng ngập tràn một thứ cảm giác ấm áp khó gọi tên.

Chỉ cần được nhìn thấy nụ cười của cậu ấy, mọi sự vất vả của ngày hôm qua dường như đều tan biến hết. Lễ tình nhân này, chắc chắn sẽ là một kỷ niệm mà mình không bao giờ quên.

\begin{table}[H]
  \centering
  \setlength{\arrayrulewidth}{1pt}
  \begin{tabular}{|>{\centering\arraybackslash}p{0.4\linewidth}|>{\centering\arraybackslash}p{0.4\linewidth}|}
    \hline
    \hyperref[choice]{Quay trở về lựa chọn} & \hyperref[ending]{Đi tới phần cuối} \\ \hline
  \end{tabular}
\end{table}

% VI. Yukihara Touka
\newpage
\label{touka}

Tôi chọn lấy một phong bì màu xanh nhạt. Nổi bật trên nền giấy là một bông hoa trắng nhỏ được ép phẳng phiu. Nét chữ trên thư vô cùng thanh lịch và gọn gàng, đúng chuẩn phong cách của một thần đồng:

\txtbx{``Sau giờ học, hãy ra bờ sông nhé. Mình sẽ đợi cậu ở đó.''\\
  \begin{flushright}-Yukihara Touka-\end{flushright}}

Bờ sông\dots\ nơi mà chúng tôi từng cùng nhau thả bộ. Dù bây giờ đang là tháng Năm, thời tiết buổi sáng khá ấm áp, nhưng khi chiều muộn buông xuống, những cơn gió thổi từ mặt sông lên vẫn mang theo cái se lạnh của những ngày cuối xuân.

Tôi cất bức thư vào túi và nhanh chóng rời khỏi trường, hướng về phía triền đê quen thuộc, nơi tôi, Kaigou và Suzuki hay đi qua trên đường tới trường.

\ct

\pov{Yukihara Touka}

Gió chiều mơn man thổi qua những vạt cỏ ven sông, cuốn theo vài cánh hoa dại li ti bay lơ lửng trong không trung. Tôi đứng tựa lưng vào lan can dọc theo con đường đi dạo, khẽ kéo cao cổ áo khoác để giữ ấm.

Dù đã là tháng Năm, nhưng thời tiết lúc chạng vạng vẫn luôn mang theo cái lạnh buốt sắc sảo. Cơ thể tôi vốn khá mẫn cảm với sự thay đổi của thời tiết, chỉ cần một cơn gió lạnh lướt qua cũng đủ khiến tôi khẽ rùng mình, dù đang mặc cả áo cardigan. Nhưng hôm nay, tôi vẫn quyết định chọn nơi này làm điểm hẹn.

Tiếng bước chân giẫm lên cỏ khô vang lên từ phía sau. Tôi quay lại, và như mong đợi, Kuon đã đến. Cậu ấy cũng đang co ro trong chiếc áo khoác đồng phục, hai tay đút sâu vào túi quần.

``Xin lỗi vì đã gọi cậu ra tận đây vào lúc chiều muộn thế này nhé.'' tôi khẽ mỉm cười, bước lại gần cậu ấy.

Kuon khẽ rùng mình một cái rồi lắc đầu. ``Không sao. Chỉ là thời tiết thay đổi đột ngột quá. Buổi sáng thì rõ ấm, thế mà giờ lại lạnh ngắt thế này.'' Cậu ấy khịt mũi, xoa xoa hai bàn tay vào nhau. Có vẻ như cả tôi và cậu ấy đều có điểm chung là rất dễ bị cảm khi thời tiết thay đổi đột ngột.

``Mình thích nơi này vào mùa này lắm.'' tôi quay mặt nhìn về phía dòng nước tĩnh lặng đang phản chiếu ánh đỏ rực của tà dương. ``Nước trong vắt, không khí lại trong lành\dots\ Nhưng tiếc là trời lạnh quá nên chúng ta không ở lại lâu được.''

``Ừm\dots\ Nếu cậu thích, cuối tuần này vào buổi trưa ấm áp hơn, chúng ta có thể ra đây đi dạo.'' Kuon đề nghị, ánh mắt vẫn nhìn theo dòng chảy của con sông.

Trái tim tôi khẽ lỡ một nhịp. Lời mời tự nhiên ấy khiến tôi cảm thấy vô cùng vui sướng, nhưng tôi cố giấu đi bằng cách hắng giọng.

``À mà này\dots\ của cậu đây.'' tôi đưa tay vào chiếc túi vải mang theo bên mình, lấy ra một hộp quà nhỏ được bọc giấy xanh nhạt tinh tế. ``Ý mình là\dots\ sô-cô-la.''

Kuon ngạc nhiên đón lấy. ``Cảm ơn cậu, Touka-san. Cậu tự làm sao?''

``Đương nhiên rồi. Mình đã mất công tìm hiểu tỷ lệ pha chế ca-cao và đường sao cho hoàn hảo nhất đấy. Mình đã rất vất vả làm nó, nên cậu phải ăn hết đi nhé?'' tôi hơi khoanh tay lại, ngước nhìn cậu ấy với vẻ mong chờ.

\begin{wrapfigure}[17]{l}{7cm}
  \includegraphics[width=7cm]{../../../../Image/Book?/valentine/touka.png}
\end{wrapfigure}

Kuon phì cười trước điệu bộ của tôi, rồi nhẹ nhàng mở hộp quà. Bên trong là những viên sô-cô-la được đổ khuôn hình những bông hoa nhỏ nhắn, vô cùng đẹp mắt. Cậu ấy cẩn thận lấy một viên cho vào miệng, chậm rãi nhắm mắt lại để thưởng thức.

``Thế nào?'' tôi không giấu nổi sự hồi hộp, bước lên phía trước một bước.

``Rất tuyệt,'' Kuon mở mắt, gật đầu tán thưởng. ``Sô-cô-la tan chảy ngay trên đầu lưỡi, độ ngọt vừa phải, và hậu vị ca-cao rất sâu. Không hổ danh là thần đồng, cậu làm gì cũng hoàn hảo cả.''

Nghe những lời khen ngợi ấy, hai má tôi bỗng nóng ran lên. Dù đã được nhiều người khen ngợi về thành tích học tập, nhưng lời khen từ Kuon lại mang đến một cảm giác hoàn toàn khác biệt hơn hản.

``Đợi đã, Touka-san. Tai cậu đỏ bừng lên kìa. Cậu đang sốt sao?'' Kuon lo lắng nghiêng người tới, định đưa tay lên trán tôi.

Tôi vội lùi lại một bước, luống cuống lấy hai tay ôm lấy đôi tai đang nóng bừng của mình. ``C-Cái gì\dots\ tai mình đỏ bừng lên à?\dots\ Ch-Chắc là do lạnh thôi!''

``Thật không? Nhưng nãy giờ cậu đâu có run?'' Kuon nhướng mày nghi hoặc.

Để dập tắt sự tò mò của cậu ấy, tôi liền đưa tay đón lấy hộp quà nhỏ mà Kuon vừa lấy từ trong túi ra.

``A, đây là của cậu làm sao? Cảm ơn nhé, mình sẽ ăn nó ở nhà!'' tôi quay ngoắt đi, giấu đi khuôn mặt đang đỏ ửng, bước nhanh về phía con đường lớn. ``Trời ạ, thôi, mình về sớm đây! Cứ đứng mãi ở đây cậu sẽ bị cảm mất!''

Tôi nghe thấy tiếng cậu ấy bật cười ở phía sau, rồi tiếng bước chân chạy theo để đuổi kịp tôi. Dưới bầu trời hoàng hôn se lạnh của tháng Năm, hai chiếc bóng đổ dài trên mặt đê, sóng bước bên nhau. Dù gió lạnh vẫn thổi, nhưng trong lòng tôi lại ấm áp vô cùng.

\begin{table}[H]
  \centering
  \setlength{\arrayrulewidth}{1pt}
  \begin{tabular}{|>{\centering\arraybackslash}p{0.4\linewidth}|>{\centering\arraybackslash}p{0.4\linewidth}|}
    \hline
    \hyperref[choice]{Quay trở về lựa chọn} & \hyperref[ending]{Đi tới phần cuối} \\ \hline
  \end{tabular}
\end{table}

% VII. Inukai Akane
\newpage
\label{akane}

Một chọn một chiếc phong bì màu nâu nhạt. Điểm nhấn của bức thư là một hình dán bé xíu mang hình dấu chân cún vô cùng đáng yêu. Nét chữ bên trong có phần vội vã nhưng vẫn đọc được rõ ràng:

\txtbx{``Sau giờ học, gặp mình ở khuôn viên sau trường nhé!''\\
  \begin{flushright}-Inukai Akane-\end{flushright}}

Khuôn viên sau trường vốn dĩ rất ít người qua lại, đặc biệt là vào giờ tan tầm khi học sinh đã vội vã ra về hoặc tham gia câu lạc bộ. Tôi cất bức thư vào túi và hướng bước chân đến điểm hẹn.

\ct

\pov{Inukai Akane}

Tôi đi đi lại lại dưới tán cây lớn sau trường, hai bàn tay đan chặt vào nhau, thỉnh thoảng lại đưa mắt nhìn quanh quất. Nhịp tim tôi đập thình thịch liên hồi.

Bình thường, tôi và Yuka lúc nào cũng ở bên cạnh nhau, không rời nửa bước. Cô ấy là người thấu hiểu tôi nhất, là nơi tôi luôn tìm thấy sự an tâm. Nhưng hôm nay, cậu ấy lại có một buổi hẹn nhận sô-cô-la ngày Lễ tình nhân ở đâu đó rồi. Thiếu vắng bóng dáng quen thuộc ấy, tôi bỗng cảm thấy bồn chồn và trống trải đến lạ.

Nhưng quan trọng hơn cả, người mà tôi sắp gặp mặt\dots\ lại là Kuon. Từ khi bước chân vào trường Aogami đến nay, cậu ấy là người con trai đầu tiên mà tôi chủ động bắt chuyện. Sự thật thì tôi đã từng nói chuyện với cậu ấy ở đâu đó, nhưng tôi hoàn toàn không nhớ đó là khi nào và ở đâu. Tuy nhiên, mỗi lần vô tình chạm mặt, tôi lại có cảm giác quen thuộc một cách kỳ lạ. Có lẽ tôi đã bị thu hút bởi sự ấm áp và dịu dàng của cậu ấy.

Nghe tiếng bước chân lạo xạo trên mặt cỏ tiến lại gần, tôi giật bắn mình quay lại. Là cậu ấy!

``À\dots! Gặp cậu ở đây\dots\ Thật là tình cờ nhỉ\dots!'' tôi lúng túng buột miệng, quên mất rằng chính mình là người đã gửi thư hẹn cậu ấy ra đây.

Kuon khẽ chớp mắt, rồi bật cười nhẹ: ``Tình cờ sao? Không phải cậu gọi tôi ra đây à?''

``T-Thật à\dots? Ừm\dots\ Mình có chút việc phải làm trên đường về\dots'' tôi nói lắp bắp, lảng tránh ánh mắt của cậu ấy, hai tay siết chặt lấy quai cặp.

``Cậu ổn chứ, Akane-san? Trông cậu cứ bồn chồn sao ấy.'' Kuon nghiêng đầu nhìn tôi, ánh mắt mang theo sự quan tâm xen lẫn tò mò. ``Bình thường có bao giờ thấy cậu đi một mình đâu. Yuka-san đâu rồi?''

Nghe nhắc đến Yuka, tôi khẽ bĩu môi. ``Cậu ấy đi nhận sô-cô-la rồi\dots\ Còn mình thì\dots\ N-Này nhé!? Ý cậu là mình trông hơi lạ, hơi khác thường sao?'' tôi vội vàng ngẩng lên, hai má nóng ran. ``\dots Mình, mình không giấu gì đâu nhé!''

``Đã ai nói gì đâu mà nhắng nhít thế,'' Kuon bật cười.

Chẳng cần nói gì thêm, chỉ ánh mắt cậu ấy thôi cũng đủ để nhìn thấu mọi sự bối rối của tôi. Sự dịu dàng và kiên nhẫn ấy khiến nhịp tim tôi lại càng đập nhanh hơn. Cậu ấy luôn đối xử bình đẳng và ấm áp với mọi người, có lẽ đó là lý do mà một người khó mở lòng như tôi lại có thể tự nhiên trò chuyện cùng cậu ấy.

``Nhưng mà\dots\ haha\dots\ Mình đoán là trông mình kỳ lạ lắm nhỉ.'' tôi gãi gãi má, cuối cùng cũng chịu thở phào một hơi, để lộ ra sự yếu đuối của mình.

``Không đâu. Cậu cứ tự nhiên đi.'' Kuon nhẹ nhàng động viên.

Nhận được sự khích lệ ấy, tôi hít một hơi thật sâu, đưa tay vào cặp sách và lấy ra một chiếc hộp sô-cô-la được bọc cẩn thận.

``Ừm, ừm\dots\ thực ra, mình có cái này cho cậu!'' tôi đưa hai tay nâng chiếc hộp ra trước mặt cậu ấy, nhắm tịt mắt lại. ``Mình quên đưa nó cho cậu ở trường\dots\ À không, là mình không dám đưa lúc ở đông người\dots''

\begin{wrapfigure}[17]{r}{7cm}
  \includegraphics[width=7cm]{../../../../Image/Book?/valentine/akane.png}
\end{wrapfigure}

Kuon ồ lên một tiếng nhỏ, rồi vui vẻ đón lấy chiếc hộp. ``Cảm ơn cậu, Akane-san. Cậu đã tự làm nó à?''

``Ừm! Mình đã cố gắng làm nó cho cậu đấy. Nên cậu phải ăn thử luôn đi!'' tôi hé một mắt ra nhìn, háo hức chờ đợi phản ứng của cậu ấy.

Kuon cẩn thận mở chiếc hộp, lấy ra một viên sô-cô-la nhỏ nhắn và cho vào miệng. Cậu ấy nhai chậm rãi, rồi gật đầu tán thưởng. ``Rất ngon. Sô-cô-la có độ ngọt vừa phải, lại còn có thêm hạt dẻ nữa. Akane-san thật đấy.''

Nghe thấy lời khen ấy, chiếc đuôi vô hình sau lưng tôi như muốn ngoáy tít lên vì sung sướng. ``Thật sao? Cậu thích nó chứ? Hay quá!'' mọi sự bồn chồn và lo lắng ban nãy dường như đều bay biến sạch.

Kuon mỉm cười, rồi lấy từ trong túi áo ra một hộp nhỏ khác đưa cho tôi. ``Tôi cũng có cái này cho cậu đây, Akane-san. Dù không được đẹp như của cậu, nhưng hy vọng cậu sẽ thích.''

Tôi mở to mắt ngạc nhiên, nhận lấy chiếc hộp sô-cô-la. ``Woah! Của cậu làm sao? Mình ăn luôn được không?''

``Cứ tự nhiên đi. ''

Tôi nhanh chóng cho một miếng vào miệng. Vị đắng nhẹ của cacao xen lẫn vị ngọt ngào của sữa tươi tan chảy trên đầu lưỡi. ``Ngon quá đi mất! Kuon-kun, cậu giỏi thật đấy!''

Cậu ấy khẽ bật cười, vươn tay xoa nhẹ lên mái tóc tôi. ``Cậu thích là tốt rồi.''

Cảm nhận hơi ấm từ bàn tay cậu ấy truyền qua đỉnh đầu, tôi bất giác đứng im, khẽ dụi đầu vào tay cậu ấy như một thói quen. Ánh nắng chiều tà xuyên qua kẽ lá, rọi xuống khuôn mặt dịu dàng của Kuon. Dù Yuka không có ở đây, nhưng ngay lúc này, tôi lại cảm thấy vô cùng hạnh phúc và an tâm.

Có lẽ, Lễ tình nhân năm nay cũng không tệ chút nào.

\begin{table}[H]
  \centering
  \setlength{\arrayrulewidth}{1pt}
  \begin{tabular}{|>{\centering\arraybackslash}p{0.4\linewidth}|>{\centering\arraybackslash}p{0.4\linewidth}|}
    \hline
    \hyperref[choice]{Quay trở về lựa chọn} & \hyperref[ending]{Đi tới phần cuối} \\ \hline
  \end{tabular}
\end{table}

% VIII. Hoshimachi Kaoru
\newpage
\label{kaoru}

Trở lại dãy tủ giày, tôi lướt qua những bức thư còn lại và chọn lấy một phong bì màu xanh cổ vịt đầy cá tính. Nổi bật trên nền giấy là một hình dán con vịt trắng nhỏ xíu trông khá ngộ nghĩnh. Làm tôi liên tưởng tới việc Suba-chan ở thế giới gốc toàn bị trêu là ``vịt'' chỉ vì cái chất giọng khản đặc giống như Vịt D*nald - mà cũng từ buổi ASMR thảm họa năm đó chứ đâu!

Nét chữ bên trong mạnh mẽ, dứt khoát:

\txtbx{``Sau giờ học, ra khuôn viên trường gặp mình nhé! Đừng có bùng đấy!''\\
  \begin{flushright}-Ukai Kaoru-\end{flushright}}

Tôi mỉm cười trước phong cách đặc trưng của cô bạn được coi là thủ lĩnh tinh thần của lớp. Cất bức thư vào túi, tôi bước về phía khuôn viên trường - nơi có những hàng cây xanh mướt và những băng ghế đá tĩnh lặng.

\ct

\pov{Ukai Kaoru}

Tôi đứng tựa lưng vào một gốc cây, chân không ngừng nhịp nhịp trên mặt đất. Tay tôi siết chặt một hộp quà màu hồng thắt ruy-băng đỏ rực. Bình thường tôi chẳng bao giờ thấy run khi đứng trước đám đông hay dẫn dắt cả lớp, thế mà không hiểu sao lúc này tim tôi lại đập như đánh trống trận.

``Kaoru-san!''

Tiếng gọi khiến tôi giật bắn mình. Kuon đang đi tới, trên môi là nụ cười quen thuộc.

``A, Kuon-san! Cậu đến rồi à?'' Tôi cố giữ giọng mình thật bình thản, dù cảm thấy hai má đang bắt đầu nóng lên.

``Cảm ơn vì đã đợi. Mà này\dots'' Kuon chỉ tay vào bức thư tôi gửi. ``Tại sao cậu lại chọn hình con vịt cho bức thư vậy? Trông nó ngộ nghĩnh thật đấy.''

Tôi khựng lại một giây, rồi gãi gãi đầu cười trừ: ``À, cái đó á? Thì\dots\ họ của mình là Ukai mà, nghe cũng gần giống 'Ukaduck' nhỉ? Với lại mình thấy nó lạch bạch dễ thương giống tính cách mình lúc đang vội ấy mà! Đừng có cười đấy!''

Kuon bật cười sảng khoái: ``Không đâu, tôi thấy nó rất hợp với cậu. Tưởng cậu chọn con vịt này bởi vì cái giọng khản đặc như D*nald đấy chứ?''

``Cái cậu này\dots!'' tôi hơi bĩu môi, còn cậu ấy thì dường như thích thú hơn khi thấy vẻ mặt hờn dỗi của tôi.

Bầu không khí bỗng chốc trở nên nhẹ nhàng hơn. Tôi hít một hơi thật sâu, đưa hộp quà màu hồng ra trước mặt cậu ấy, nhưng miệng thì vẫn không ngừng tuôn ra những lời bào chữa.

``Hả\dots? Cái này á? K-không, đừng có hiểu lầm nhé! Chuyện là một đứa bạn của mình muốn tặng sô-cô-la tự làm cho người nó thích, nên mình đã giúp nó một tay thôi\dots''

Kuon rướn mày nhìn tôi, vẻ mặt đầy thú vị. ``Thật á?''

``Nhưng mà mình lỡ tay phấn khích quá nên làm hơi bị nhiều! Vừa nãy mình cũng định đi phát cho cả lớp mỗi người một ít, nên mình nghĩ mình cũng nên tặng cậu một hộp luôn cho đỡ phí công làm thôi\dots\ Ừm, vậy thì\dots''

\begin{wrapfigure}[17]{l}{7cm}
  \includegraphics[width=7cm]{../../../../Image/Book?/valentine/kaoru.png}
\end{wrapfigure}

Tôi cảm thấy mình càng nói càng rối. Nhìn vào ánh mắt dịu dàng nhưng cũng đầy trêu chọc của Kuon, tôi bỗng cảm thấy những lời nói dối vụng về này thật đáng xấu hổ.

``\dots\ Thật ra thì, không phải thế đâu.'' Tôi cúi mặt xuống, giọng nói nhỏ hẳn đi nhưng đầy chân thành. ``Mình đã chuẩn bị hộp này vì thực sự muốn tặng riêng cho cậu thôi. Cậu nhận nó nhé?''

Khoảnh khắc ấy, tôi cảm nhận được một bàn tay ấm áp nhận lấy hộp quà từ tay mình.

``Cảm ơn cậu, Kaoru-san. Tôi rất sẵn lòng.'' Kuon mỉm cười rạng rỡ. Cậu ấy mở hộp quà và lấy ra một viên sô-cô-la. ``Ngon lắm. Vị đắng vừa phải, ăn khá vừa miệng.''

Thấy cậu ấy thích, tôi bỗng cảm thấy nhẹ nhõm vô cùng. ``Thật sao? May quá! Mình cứ sợ nó sẽ hơi quá đắng cơ.''

Kuon sau đó cũng lấy ra một chiếc hộp nhỏ từ trong túi đưa cho tôi. ``Đây là phần của tôi. Dù không được bọc ruy-băng rực rỡ như của cậu, nhưng hy vọng cậu sẽ thích.''

Tôi tròn mắt ngạc nhiên, đón lấy món quà với niềm vui khó tả. Tôi nhanh chóng cho một miếng vào miệng. Vị ngọt ngào lan tỏa khiến tôi không tự chủ được mà nheo mắt cười.

``Ưm! Ngon quá! Cảm ơn cậu nhé, Kuon-san!''

Dưới tán cây xanh, chúng tôi cùng nhau trò chuyện và thưởng thức sô-cô-la. Cảm giác này\dots\ còn tuyệt vời hơn cả việc uống một ly sữa ấm mỗi sáng để mong cao thêm nữa. Có lẽ, Lễ tình nhân năm nay chính là kỷ niệm đẹp nhất của tôi ở trường Aogami.

\begin{table}[H]
  \centering
  \setlength{\arrayrulewidth}{1pt}
  \begin{tabular}{|>{\centering\arraybackslash}p{0.4\linewidth}|>{\centering\arraybackslash}p{0.4\linewidth}|}
    \hline
    \hyperref[choice]{Quay trở về lựa chọn} & \hyperref[ending]{Đi tới phần cuối} \\ \hline
  \end{tabular}
\end{table}

% IX. Mizuta Uzuki
\newpage
\label{uzuki}

Lướt qua những bức thư trên tủ giày, tôi chọn lấy một phong bì màu xanh tím nhạt khá lạ mắt. Nổi bật trên nền giấy là một hình dán củ cà rốt nhỏ xinh. Nét chữ bên trong có vẻ hơi bay bổng nhưng đầy cá tính:

\txtbx{``Sau giờ học, gặp mình ở phòng học trống ở khu CLB nhé! Có việc quan trọng cần nhờ cậu đây.''\\
  \begin{flushright}-Mizuta Uzuki-\end{flushright}}

Tôi tự hỏi cô bạn nhiếp ảnh gia tài năng này lại định bày trò gì đây. Cất bức thư vào túi, tôi bước về phía khu nhà của câu lạc bộ - nơi giờ này có khi chẳng còn ai hoạt động nữa, và vì thế, rất thích hợp với những người muốn có một không gian riêng tư.

\ct

\pov{Mizuta Uzuki}

Tôi đang loay hoay căn chỉnh ống kính máy ảnh, miệng khẽ lẩm bẩm tính toán thông số ánh sáng. ``\textit{Chỗ này đặt góc nào mới chuẩn nhỉ\dots\ Có nên dịch ra một chút để thêm ánh sáng không ta?}''

Nghe tiếng cửa lớp mở ra, tôi ngước lên và thấy Kuon-san đang đứng đó.

``A, Kuon-san, cậu đến rồi à!'' tôi vẫy vẫy tay ra hiệu cho cậu ấy lại gần.

``Cậu gọi tôi ra đây có việc gì vậy, Uzuki-san? Lại định nhờ làm mẫu ảnh nữa sao?'' Kuon bước tới, nhìn đống thiết bị của tôi với vẻ tò mò.

Tôi gật đầu, chìa ra một chiếc hộp màu đen được thắt ruy băng đỏ. ``Đúng rồi đó! Mình đang định lấy ngày Lễ tình nhân làm chủ đề cho bộ ảnh mới, nên cậu hãy giữ chặt cái hộp này và đứng vào góc kia nhé!''

Kuon nhận lấy chiếc hộp, nét mặt có chút hơi do dự: ``Ờm\dots\ Được thôi. Nhưng mà--''

Tôi vội cắt lời: ``Được rồi\dots\ đừng nhúc nhích nhé\dots\ Một, hai, ba!''

*TÁCH!*

Tôi kiểm tra màn hình máy ảnh. ``\dots\ Ừm, thành quả rất tốt. Ánh sáng và góc độ đều ổn cả.''

Cậu ấy bước lại gần xem ảnh. ``Trông cũng ổn đó, nhưng mà này\dots\ Cậu có nghĩ là bức ảnh mình đang trao cái hộp cho cậu sẽ trông chân thật và 'đúng chất' Valentine hơn không?''

Tôi hơi sững người lại trước lời gợi ý của cậu ấy. ``Hả? Ý cậu là bức ảnh mình trao sô-cô-la cho cậu chắc chắn đẹp hơn à?''

Kuon gật đầu: ``Thì chủ đề là Lễ tình nhân mà, phải có sự tương tác trao nhận mới đúng chứ.''

Tôi khẽ đỏ mặt, đưa tay chỉnh lại gọng kính. ``Ừ-ừm\dots\ mình đoán phiên bản đó cũng cần thiết\dots\ hay gì đó. Vậy thì đổi lại đi, cậu cầm máy ảnh, mình sẽ thực hiện hành động trao quà.''

Cậu ta cầm lấy máy ảnh, tôi cầm lại cái hộp, hít một hơi thật sâu rồi đưa ra trước mặt cậu ấy với một nụ cười có phần hơi gượng gạo vì ngượng.

``\dots\ Ừ, thế này được không?''

``Cậu muốn một biểu cảm nghiêm túc hơn à?'' Kuon trêu chọc khi thấy tôi cứ nhìn lảng đi chỗ khác.

``Mồ! Cậu thật là\dots\ Tiếp theo mình sẽ nghĩ kỹ hơn về biểu cảm!'' Tôi giật lại máy ảnh, mặt đỏ bừng.

Không gian trong phòng học trống bỗng trở nên tĩnh lặng đến lạ thường. Tôi nhìn cái hộp trên bàn, rồi lại nhìn Kuon. Một ý nghĩ chợt lóe lên trong đầu.

``Vậy là xong rồi chứ gì? Nếu không còn gì nữa thì--''

``À, đợi chút.'' Tôi gọi khi thấy cậu ấy định rời đi.

\begin{wrapfigure}[17]{r}{7cm}
  \includegraphics[width=7cm]{../../../../Image/Book?/valentine/uzuki.png}
\end{wrapfigure}

``Sao vậy?''

Tôi cầm cái hộp, lấy hết can đảm đưa lại cho cậu ấy. ``Cái này\dots\ dù sao thì mình cũng đã đưa cho cậu một lần để làm mẫu rồi, nên giờ cậu cứ giữ lấy mà mang về đi. Coi như là thù lao cho buổi làm mẫu hôm nay nhé!''

Kuon ngẩn người ra một giây, rồi bật cười rạng rỡ, nhận lấy chiếc hộp. ``Cảm ơn cậu nhé, Uzuki-san. Mình sẽ thưởng thức nó thật ngon miệng.'' Cậu ấy mở ra và nếm thử một viên. ``Ngon lắm.''

Nói rồi, Kuon lấy từ trong túi ra một hộp nhỏ hình vuông được bọc giấy trang nhã đưa cho tôi. ``Đây là phần của tôi dành cho cậu. Nó có hình chiếc lá, hy vọng cậu sẽ thấy thích.''

Tôi ngạc nhiên đón lấy hộp quà. Vừa mở ra, hương thơm ngọt ngào của cacao đã lan tỏa. Tôi lấy một viên hình chiếc lá cắn nhẹ. Vị ngọt lịm tan chảy nơi đầu lưỡi.

Cùng lúc đó, thấy Kuon đang nhấm nháp phần sô-cô-la của mình với vẻ mặt vô cùng tận hưởng, tôi nhanh tay đưa máy ảnh lên và bấm máy.

*TÁCH!*

Cậu ấy giật mình nhìn tôi. Còn tôi thì chỉ khẽ mỉm cười, giấu khuôn mặt đang đỏ ửng sau thân máy ảnh. Bức ảnh ghi lại nụ cười rạng rỡ ấy, cùng hương vị sô-cô-la của cậu ấy đang tan trên đầu lưỡi, chắc chắn sẽ là thước phim quý giá nhất trong bộ sưu tập cũng như ký ức của tôi.

``Hứ! Phải ăn hết đấy nhé, không là mình bắt cậu làm mẫu ảnh thêm mười buổi nữa đó!''

\begin{table}[H]
  \centering
  \setlength{\arrayrulewidth}{1pt}
  \begin{tabular}{|>{\centering\arraybackslash}p{0.4\linewidth}|>{\centering\arraybackslash}p{0.4\linewidth}|}
    \hline
    \hyperref[choice]{Quay trở về lựa chọn} & \hyperref[ending]{Đi tới phần cuối} \\ \hline
  \end{tabular}
\end{table}

% X. Amano Kanade
\newpage
\label{kanade}

Tôi lấy một phong bì màu trắng tinh khôi trong xấp phong bì này. Nổi bật trên nền giấy là một hình dán đôi cánh thiên thần nhỏ nhắn. Nét chữ bên trong gọn gàng, có chút rụt rè:

\txtbx{``Sau giờ học, đến phòng phát thanh gặp mình nhé. Mình\dots\ có chuyện muốn nói.''\\
  \begin{flushright}-Amano Kanade-\end{flushright}}

Phòng phát thanh\dots\ Giờ này chắc hẳn Kana đang chuẩn bị cho buổi phát thanh cuối ngày của cô ấy. Tôi cất bức thư vào túi và nhanh chóng đi lên lầu.

Khi bước đến cửa phòng phát thanh, tôi ngạc nhiên khi thấy cánh cửa đang hé mở, nhưng không hề nghe thấy giọng của Kana như mọi khi. Thay vào đó, một giọng nói trong trẻo nhưng đầy bối rối vang lên từ bên trong.

\ct

\pov{Amano Kanade}

Tôi ngồi trước micro trong phòng phát thanh, hai tay vân vê chiếc túi màu xanh được thắt nơ vàng cẩn thận. Nhìn đồng hồ nhích từng giây, lòng tôi càng thêm bồn chồn.

``Kana-san đến muộn quá\dots'' tôi thở dài. Cô bạn thân của tôi bảo tôi cứ lên phòng phát thanh đợi trước, thế mà giờ này vẫn chưa thấy bóng dáng đâu.

Cuối cùng, đã đến lúc lên sóng mà mình vẫn chưa kịp trao sô-cô-la cho Kuon. Rõ ràng là mình đã chuẩn bị rất kỹ, nhưng khi đối diện với cậu ấy ở lớp, mình lại chẳng thốt nên lời.

``Giá mà\dots'' tôi tự lẩm bẩm, ánh mắt nhìn chằm chằm vào chiếc micro chưa bật nguồn trước mặt. ``Lẽ ra mình có thể nói rõ ràng hơn trước micro trong phòng phát thanh\dots''

Tôi hít một hơi thật sâu, nhắm mắt lại và thử tưởng tượng cậu con trai ấy đang đứng trước mặt mình.

``Phù\dots'' tôi thu hết can đảm, hắng giọng. ``Mình là Amano Kanade. Mình\dots\ mình làm sô-cô-la này chỉ dành riêng cho cậu thôi!''

``Cảm ơn cậu, Kanade-san. Mình rất vui.''

Một giọng nói trầm ấm đột ngột vang lên ngay từ phía cửa. Tôi giật bắn mình, mở choàng mắt ra và nhìn thấy Kuon đang đứng đó tự lúc nào.

``H-Hả?! S-Sao cậu lại ở đây?'' tôi luống cuống đứng bật dậy, suýt chút nữa làm rơi cả chiếc túi sô-cô-la. ``Dù sao thì đây cũng là phòng phát thanh mà\dots!''

Kuon mỉm cười, bước vào trong và nhẹ nhàng khép cửa lại. ``Tôi nhận được thư của cậu hẹn ở tủ giày, nên mới đến đây. Mà nhắc mới nhớ, Kana-san đâu rồi?''

Tôi chớp mắt, đầu óc quay cuồng. Bức thư\dots? Kana bảo sẽ giúp mình gửi thư\dots\ ``Có lẽ\dots\ cậu ấy đã gọi cậu đến đây sao?'' tôi lầm bầm, chợt nhận ra mình đã mắc bẫy của cô bạn thân. Kana cố tình đến muộn để tạo không gian riêng cho hai đứa!

Và điều tồi tệ nhất là\dots

``Cậu\dots\ cậu đã nghe thấy hết rồi sao?'' tôi ôm lấy hai má đang nóng rực như lửa đốt, muốn tìm một cái hố để chui xuống. Không ngờ lời tập dượt ngớ ngẩn vừa rồi lại bị cậu ấy nghe thấy toàn bộ!

Kuon gật đầu, ánh mắt ánh lên sự dịu dàng. ``Ừ. Và tôi rất muốn nhận lấy món quà ấy, nếu cậu không phiền.''

\begin{wrapfigure}[17]{l}{7cm}
  \includegraphics[width=7cm]{../../../../Image/Book?/valentine/kanade.png}
\end{wrapfigure}

Dù sự xấu hổ vẫn bủa vây, nhưng nhìn nụ cười chân thành của cậu ấy, tôi cảm thấy một luồng dũng khí ấm áp len lỏi trong tim. Tôi bước tới, đưa chiếc túi xanh thắt nơ vàng cho cậu ấy.

``Nhưng\dots\ cuối cùng thì, mình cũng có thể tự miệng nói cho cậu biết.'' tôi khẽ mỉm cười, trái tim đập rộn ràng. ``Xin hãy nhận lấy.''

Kuon cẩn thận nhận chiếc túi, rồi mở ra. Bên trong là những viên sô-cô-la được nắn nót tạo hình cẩn thận. Cậu ấy nếm thử một viên và khẽ gật đầu. ``Sô-cô-la rất ngon, độ ngọt thật dịu nhẹ. Giống như chính cậu vậy, Kanade-san.''

Câu nói ấy khiến tôi đỏ mặt tía tai.

Sau đó, cậu cũng lấy từ trong túi ra một chiếc hộp nhỏ đưa cho tôi. ``Đây là phần của tôi. Dù không tinh tế bằng món quà của cậu, nhưng tôi đã tự tay làm nó.''

Tôi nhận lấy, cảm giác hạnh phúc trào dâng. Nếm thử một viên, vị đắng hòa quyện cùng vị ngọt ngào như tan chảy trong miệng. ``Ngon quá\dots\ Cảm ơn cậu, Kuon-san.''

Chúng tôi cùng nhau đứng trong phòng phát thanh, tận hưởng khoảnh khắc ngọt ngào của ngày Lễ tình nhân. Dù Kana không đến, nhưng tôi thầm cảm ơn cô ấy vì đã tạo ra cơ hội tuyệt vời này.

\begin{table}[H]
  \centering
  \setlength{\arrayrulewidth}{1pt}
  \begin{tabular}{|>{\centering\arraybackslash}p{0.4\linewidth}|>{\centering\arraybackslash}p{0.4\linewidth}|}
    \hline
    \hyperref[choice]{Quay trở về lựa chọn} & \hyperref[ending]{Đi tới phần cuối} \\ \hline
  \end{tabular}
\end{table}

% XI. Mizuno Aku
\newpage
\label{aku}

Ánh mắt tôi dừng lại ở một phong bì có màu sắc khá độc đáo: sự kết hợp giữa hồng nhạt và xanh lam nhạt, gợi nhắc ngay đến mái tóc hai màu đặc trưng của Mizuno Aku, hay Aku-tan ở thế giới thực. Nổi bật trên nắp phong bì là một miếng sticker hình tay cầm chơi game retro vô cùng dễ thương.

Nét chữ bên trong có vẻ hơi run, dường như người viết đã phải đấu tranh tư tưởng rất nhiều:

\txtbx{``Sau giờ học\dots\ cậu có thể ghé qua hành lang vắng ở tầng 3 khu nhà cũ không? Mình\dots mình có chuyện muốn nói. Làm ơn đừng cười mình nhé!''\\
  \begin{flushright}-Mizuno Aku-\end{flushright}}

Khu nhà cũ vốn ít người qua lại, nhất là vào giờ này. Tôi cất bức thư vào túi và bắt đầu leo cầu thang lên.

\ct

\pov{Mizuno Aku}

Tôi đứng tựa lưng vào tường hành lang, hai tay siết chặt lấy hộp quà nhỏ được bọc giấy bóng kiếng màu tím. Lồng ngực tôi phập phồng không yên, đầu óc thì cứ vẩn vơ tới những viễn cảnh tồi tệ nhất.

``\textit{Tại sao mình lại làm chuyện này chứ\dots?}'' tôi thầm than thở.

Bình thường, tôi chỉ là một đứa hậu đậu luôn bị cô bạn thân Shion sai bảo chạy việc vặt. Đi lấy tài liệu, mua nước ngọt, hay bị lôi vào những trò ma thuật hắc ám quái đản của cậu ấy\dots\ tôi vốn dĩ đã quen với việc núp bóng người khác. Nhưng hôm nay, tôi muốn tự mình đối diện với \emph{cậu ấy}. Một mình. Không có Shion, không có game trong thế giới ảo, chỉ có tôi và tâm tư này, ở ngay đời thực.

Tiếng bước chân vang lên từ phía cầu thang. Tim tôi như muốn nhảy ra khỏi lồng ngực. Tôi vội vàng đứng thẳng dậy, nhưng vì quá luống cuống, gót giày của tôi bỗng trượt nhẹ trên mặt sàn gỗ nhẵn bóng.

``Á\dots!''

Tôi mất thăng bằng, cả người đổ nhào về phía trước. Theo bản năng, tôi nhắm chặt mắt lại, chuẩn bị tinh thần cho một cú tiếp đất đau đớn.

Nhưng cơn đau đã không đến. Thay vào đó, tôi cảm nhận được một đôi tay cứng cáp nhưng cũng rất dịu dàng đỡ lấy bả vai mình. Một mùi hương thanh khiết, ấm áp bao trùm lấy tôi.

``Cậu không sao chứ, Aku-san? Đi đứng phải cẩn thận chứ.''

Giọng nói của Kuon vang lên ngay sát bên tai. Tôi hé mắt nhìn, thấy khuôn mặt cậu ấy đang ở rất gần, ánh mắt lộ rõ vẻ lo lắng. Mặt tôi nóng bừng lên trong tích tắc, nóng đến mức tôi có cảm giác như đầu mình đang bốc khói vậy.

``K-K-Kuon-san! C-Cảm ơn cậu\dots\ mình xin lỗi, mình hậu đậu quá!'' tôi lắp bắp, vội vàng đứng thẳng dậy và lùi lại một bước, hai tay múa may loạn xạ vì ngượng.

Kuon bật cười nhẹ, vẻ mặt trở nên thoải mái hơn: ``Không sao, cậu ổn là tốt rồi. Mà này, bức thư cậu gửi\dots''

``A! Đúng rồi!'' Tôi sực nhớ ra mục đích chính, vội vàng đưa hộp quà ra trước mặt, đầu cúi thấp đến mức cằm chạm vào ngực. ``Cái này\dots\ mình\dots\ mình đã tự tay làm nó! Dù mình vốn không giỏi mấy việc tỉ mỉ này lắm, nhưng mình đã thử làm đi làm lại nhiều lần\dots\ Xin cậu hãy nhận lấy!''

Khoảng lặng trôi qua trong vài giây khiến tôi thấy dài như cả thế kỷ. Rồi tôi cảm nhận được sức nặng trên tay mình biến mất.

``Cảm ơn cậu, Aku-san. Tôi rất trân trọng nó.'' Cậu nhẹ nhàng mở hộp quà. Bên trong là những viên sô-cô-la có hình thù hơi\dots\ lạ mắt. Tôi định nặn hình những nhân vật trong game nhưng kết quả lại hơi méo mó.

Kuon lấy một viên cho vào miệng. Tôi nín thở chờ đợi, tim đập thình thịch.

``Rất tuyệt.'' Cậu ấy gật đầu, nụ cười rạng rỡ hiện rõ trên môi. ``Vị ngọt rất thanh, lại có chút giòn giòn của hạt ngũ cốc bên trong nữa. Cậu đã vất vả rồi.''

Nghe lời khen ấy, bao nhiêu lo lắng trong tôi tan biến sạch. Tôi cảm thấy vui sướng đến mức muốn nhảy cẫng lên, nhưng lại cố kìm lại để không bị vấp ngã lần nữa.

Kuon sau đó cũng lấy ra một hộp nhỏ từ túi áo. ``Đây là phần của tôi dành cho cậu. Hy vọng nó sẽ giúp cậu có thêm năng lượng để\dots\ không bị Shion-chan sai vặt nữa.''

Tôi ngạc nhiên đón lấy món quà, cảm giác ấm áp lan tỏa khắp cơ thể. Tôi nhanh chóng lấy một viên nếm thử. Vị sô-cô-la tan chảy, ngọt ngào và dịu nhẹ như chính sự hiện diện của cậu ấy vậy.

``C-Cảm ơn Kuon-san nhiều lắm! Mình sẽ\dots\ mình sẽ cố gắng hơn!'' tôi cười tươi, dù đôi chỗ trên mặt vẫn còn vương chút đỏ hồng vì ngượng ngùng.

Đứng ở hành lang cũ kỹ ngập tràn ánh nắng chiều, tôi nhận ra rằng, đôi khi bước ra khỏi vùng an toàn và đối diện với cảm xúc của chính mình cũng không đáng sợ đến thế. Nhất là khi người đối diện là một người ấm áp như cậu ấy.

\begin{table}[H]
  \centering
  \setlength{\arrayrulewidth}{1pt}
  \begin{tabular}{|>{\centering\arraybackslash}p{0.4\linewidth}|>{\centering\arraybackslash}p{0.4\linewidth}|}
    \hline
    \hyperref[choice]{Quay trở về lựa chọn} & \hyperref[ending]{Đi tới phần cuối} \\ \hline
  \end{tabular}
\end{table}

% XII. Ibara Saki
\newpage
\label{saki}

Tôi lướt qua những bức thư trong ngăn tủ, và một phong bì với sự kết hợp dịu dàng giữa sắc xanh lá cây của cỏ non và hồng phấn của những nụ hoa chớm nở đã thu hút tôi. Ngay giữa trung tâm là một hình dán bông hoa nhỏ xíu, tỏa ra một cảm giác thanh bình và ấm áp.

Nét chữ bên trong vô cùng ngay ngắn, thanh thoát:

\txtbx{``Sau giờ học, hãy ghé qua khu bồn hoa phía sau trường nhé. Mình có điều muốn chia sẻ cùng cậu.''\\
  \begin{flushright}-Ibara Saki-\end{flushright}}

Khu bồn hoa\dots\ nơi Saki thường dành hàng giờ để chăm sóc, và nếu tôi nhớ không nhầm, cô ấy cũng là thành viên thuộc Hội chăm sóc cây cảnh của trường, nên việc ở lại để chăm cây vào giờ này cũng không quá ngạc nhiên. Tôi cất bức thư vào túi và thong thả bước về phía vườn trường.

\ct

\pov{Ibara Saki}

Tôi khẽ nghiêng người, tay cầm bình nhỏ tưới nước cho những khóm hoa đang vươn mình trong nắng chiều tà. Mùi đất ẩm quyện cùng hương hoa dìu dịu tạo nên một bầu không khí thật dễ chịu.

Ở góc này là khóm hoa bỉ ngạn trắng tinh khiết mà tôi đã dày công chăm sóc, biểu tượng cho những ký ức thanh cao. Nhưng không chỉ có vậy, năm nay tôi còn trồng thêm cả những dải hoa oải hương tím ngắt mang lại sự bình yên, và cả những bông cúc họa mi trắng nhỏ li ti tượng trưng cho sự chân thành. Gần đó, những bông uất kim hương cũng đang bắt đầu khoe sắc. Mỗi loài hoa đều có một ngôn ngữ riêng, và tôi yêu việc lắng nghe hơi thở của chúng từng ngày.

Nghe tiếng bước chân lạo xạo trên lối đi nhỏ, tôi mỉm cười mà không cần quay lại. Mùi hương và nhịp bước chân này\dots\ tôi biết đó là ai.

``Chào cậu, Kuon-san. Cậu đến thật đúng lúc, khóm cúc họa mi này vừa nở rộ này.'' tôi quay lại, khẽ vén lọn tóc mái, mỉm cười chào cậu ấy.

Kuon bước lại gần, ánh mắt cậu ấy lướt qua những luống hoa đầy màu sắc. ``Đẹp thật đấy, Saki-san. Cậu thực sự rất mát tay. Ở đây có nhiều loài hoa lạ quá.''

``Cảm ơn cậu. Vì gia đình mình kinh doanh tiệm hoa mà, nên mình luôn muốn trường học cũng ngập tràn hương sắc như vậy. Cậu thấy khóm hoa oải hương này thế nào? Mùi hương của nó giúp giảm căng thẳng rất tốt đấy, cực kỳ hợp cho những lúc ôn thi căng thẳng.'' tôi say sưa kể về đặc tính của từng loài hoa, niềm đam mê rực cháy trong ánh mắt.

Kuon chăm chú lắng nghe, thỉnh thoảng lại gật đầu tán thưởng. Sự hiện diện điềm tĩnh của cậu ấy khiến tôi cảm thấy bình yên vô cùng.

``À\dots\ suýt nữa thì mình quên mất việc chính,'' tôi hơi đỏ mặt, khẽ đặt bình tưới xuống và lấy từ trong chiếc giỏ mây đặt bên cạnh một hộp quà nhỏ được trang trí bằng những bông hoa khô ép phẳng.

``Cái này\dots\ là món quà của mình dành cho cậu. Mong cậu nhận lấy.''

Cậu ấy ngạc nhiên đón lấy món quà. ``Cảm ơn cậu, Saki-san. Nó đẹp quá, đến mức mình chẳng nỡ mở lớp gói quà này ra nữa.''

``Cậu cứ mở đi, sô-cô-la bên trong mình cũng thêm một chút hương hoa nhài tự nhiên đấy.'' tôi khẽ cười khích lệ.

Cậu ấy cẩn thận mở hộp và thưởng thức. Ánh mắt cậu ấy sáng lên đầy thích thú. ``Tuyệt vời. Một hương vị rất thanh tao, không quá ngọt mà lại mang đến cảm giác rất nhẹ nhàng. Đúng là phong cách của cậu.''

Kuon sau đó cũng lấy từ trong túi ra một món quà nhỏ trao cho tôi. ``Của cậu đây. Hy vọng nó cũng mang lại cho cậu chút niềm vui như những bông hoa này.''

Tôi nhận lấy, trái tim như nở rộ hàng ngàn đóa hoa trong lồng ngực. Dưới ánh hoàng hôn vàng vọt, hương hoa và vị sô-cô-la quyện vào nhau, tạo nên một khoảnh khắc Valentine thật dịu dàng và khó quên.

\begin{table}[H]
  \centering
  \setlength{\arrayrulewidth}{1pt}
  \begin{tabular}{|>{\centering\arraybackslash}p{0.4\linewidth}|>{\centering\arraybackslash}p{0.4\linewidth}|}
    \hline
    \hyperref[choice]{Quay trở về lựa chọn} & \hyperref[ending]{Đi tới phần cuối} \\ \hline
  \end{tabular}
\end{table}

% XIII. Fukami Shiina
\newpage
\label{shiina}

Tôi cầm lên một phong bì màu xám trắng, trông vô cùng giản đơn nhưng lại toát lên một vẻ lạnh lùng, dứt khoát. Ở góc phong bì là một hình dán đầu sư tử bằng bạc nhỏ xíu, ánh lên tia sáng sắc lẹm. Làm tôi liên tưởng ngay tới biệt danh ``White Lionness'' của Botan ở thế giới thực.

Nét chữ bên trong ngắn gọn, đi thẳng vào vấn đề:

\txtbx{``Tan học ra phòng thể chất. Tôi đợi.''\\
  \begin{flushright}-Fukami Shiina-\end{flushright}}

Phòng thể chất\dots\ nơi chúng tôi từng có trận đấu bóng né nảy lửa trong khuôn khổ hội thao vừa rồi. Tôi cất bức thư vào túi và rảo bước về phía khu nhà thể thao.

\ct

\pov{Fukami Shiina}

Tôi ngồi trên băng ghế gỗ dài bên rìa sân bóng né, đôi chân khẽ đung đưa. Không gian rộng lớn của phòng thể chất giờ đây thật tĩnh lặng, chỉ còn lại dư âm của những tiếng hò reo từ trận đấu lần trước.

Trong tay tôi là một hộp quà vuông vức, được bọc giấy đen bóng đơn giản. Tôi đưa tay lên chạm nhẹ vào vết xước nhỏ trên gò má, cảm giác hơi xót nhưng lại khiến tôi thấy tỉnh táo hơn. Đám nhóc trường bên kia quả thực không biết lượng sức mình, cứ thích tìm đến tôi để ``thử lửa''.

Tiếng cửa phòng thể chất mở ra, cắt ngang dòng suy nghĩ của tôi. Kuon bước vào, ánh mắt cậu ấy lướt nhanh khắp phòng rồi dừng lại ở chỗ tôi.

``Ồ, cậu đến rồi à,'' tôi cất giọng, cố giữ cho âm điệu thật bình thản.

Kuon bước lại gần, rồi khẽ khựng lại khi nhìn thấy khuôn mặt tôi. ``Shiina-san? Mặt cậu\dots\ tại sao lại có vết xước thế kia?''

Tôi xoay nhẹ cổ tay, nhún vai một cái thật nhẹ: ``Có đứa muốn ra oai, thế là cho thể hiện luôn. Vết xước này cũng có nhằm nhò gì.''

Cậu ta thở dài, ánh mắt lộ rõ vẻ không hài lòng nhưng cũng mang theo sự quan tâm. ``Cậu lúc nào cũng vậy\dots\ Lần sau đừng có tự mình lao vào rắc rối như thế nữa.''

``Biết rồi mà, nói mãi.'' tôi hắng giọng, đưa hộp quà ra phía trước. ``Của cậu này. Coi như là\dots\ thù lao vì đã làm đồng đội của tôi trong trận bóng né lần trước.''

Kuon ngạc nhiên đón lấy. ``Cảm ơn cậu, Shiina-san. Tôi không ngờ cậu lại tự mình chuẩn bị sô-cô-la đấy.''

``Cũng chẳng có gì to tát đâu,'' tôi lảng mắt đi chỗ khác, cảm thấy hai tai mình hơi nóng lên.

Cậu ấy mở hộp, lấy một viên sô-cô-la đen nguyên chất nếm thử. ``Đắng thật đấy\dots\ nhưng mà rất đậm đà. Rất giống phong cách của cậu.''

``Hừm, thích là tốt rồi.'' tôi khẽ cười, rồi nhận lấy món quà đáp lễ từ Kuon. ``Cảm ơn cậu nhé. Tôi sẽ thưởng thức nó sau.''

Dưới mái vòm cao vút của phòng thể chất, chúng tôi đứng đó một lúc lâu. Dù vết xước vẫn còn hơi nhức, nhưng tôi cảm thấy tâm trạng mình lúc này thực sự không tệ chút nào.

\begin{table}[H]
  \centering
  \setlength{\arrayrulewidth}{1pt}
  \begin{tabular}{|>{\centering\arraybackslash}p{0.4\linewidth}|>{\centering\arraybackslash}p{0.4\linewidth}|}
    \hline
    \hyperref[choice]{Quay trở về lựa chọn} & \hyperref[ending]{Đi tới phần cuối} \\ \hline
  \end{tabular}
\end{table}

% XIV. Utachi Azuki
\newpage
\label{azuki}

Tôi phát hiện ra một phong bì vô cùng bắt mắt với hiệu ứng chuyển màu từ hồng đậm sang hồng nhạt đầy tinh tế. Ngay góc phong bì là một hình dán nốt nhạc nhỏ xíu được ép nhũ vàng, trông như thể nó đang chuẩn bị phát ra những giai điệu du dương ngay khi tôi chạm vào.

Nét chữ bên trong vô cùng mềm mại và bay bổng, tựa như một dòng suối âm thanh:

\txtbx{``Sau giờ học, hãy ghé qua phòng CLB Âm nhạc một chút nhé. Mình muốn dành tặng cậu một giai điệu đặc biệt.''\\
  \begin{flushright}-Utachi Azuki-\end{flushright}}

Phòng CLB Âm nhạc\dots\ một không gian vốn dĩ đã quá quen thuộc với một người yêu thích âm nhạc như Azuki-san. Tôi cất bức thư và hướng bước chân về phía khu nhà đa năng.

\ct

\pov{Utachi Azuki}

Tôi ngồi bên chiếc đàn piano trong phòng CLB, đôi tay lướt nhẹ trên những phím đàn trắng đen. Những giai điệu trong trẻo vang vọng trong không gian tĩnh lặng của buổi chiều tà, như đang kể lại những câu chuyện về những chuyến phiêu lưu mà tôi hằng mong ước từ thuở nhỏ.

Âm nhạc luôn là lẽ sống, là người bạn đồng hành trung thành nhất của tôi. Nhưng từ khi gặp Kuon-san, tôi nhận ra rằng thế giới này còn có những bản nhạc sống động hơn cả những nốt nhạc trên trang giấy. Đó là tiếng cười của cậu ấy, là ánh mắt khích lệ, và cả sự hiện diện ấm áp khiến tâm hồn tôi luôn cảm thấy bình yên.

Tiếng cửa phòng khẽ mở. Tôi dừng tay, quay lại và thấy Kuon-san đang đứng đó, ánh mắt lộ rõ vẻ say mê khi lắng nghe âm nhạc.

``Giai điệu lúc nãy tuyệt thật đấy, Azuki-san. Hình như tôi chưa từng nghe bản nhạc này bao giờ,'' cậu ấy bước lại gần, nụ cười hiền hậu nở trên môi.

``Đó là một bản nhạc mình mới sáng tác đấy. Mình vẫn chưa đặt tên đâu, nhưng mình muốn dành tặng nó cho người đã mang đến cảm hứng cho mình,'' tôi khẽ cười, cảm nhận hơi ấm lan tỏa trong lồng ngực khi thấy cậu ấy chăm chú lắng nghe.

``Vậy sao? Thật vinh dự cho tôi quá,'' Kuon ngồi xuống chiếc ghế bên cạnh, tạo ra một khoảng không gian gần gũi đến kỳ lạ.

``Kuon-san này, cậu biết đấy, mình luôn khao khát những chuyến hành trình mới. Nhưng hôm nay, mình muốn chia sẻ một điều ngọt ngào và tĩnh lặng hơn cả âm nhạc với cậu.'' Tôi lấy ra một hộp quà nhỏ có hình dáng giống như một chiếc đàn piano mini vô cùng tinh xảo.

``Món quà Valentine này\dots\ mình đã gửi gắm vào đó tất cả những tâm tư của mình. Mong cậu hãy nhận lấy tâm ý này.''

Cậu ấy ngạc nhiên đón lấy món quà. ``Cảm ơn cậu, Azuki-san. Một chiếc hộp quà thật đặc biệt, đúng chất nghệ sĩ luôn.''

Cậu ấy mở hộp và nếm thử sô-cô-la bên trong. Ánh mắt cậu ấy khẽ thay đổi. ``Vị ngọt rất thanh, và có một chút hương vị gì đó rất tươi mới\dots\ như thể đang được đi dạo giữa một cánh đồng hoa dưới nắng sớm vậy.''

``Đó là hương vị của tự do và những khởi đầu mới đấy.'' tôi cười tươi, hạnh phúc khi thấy cậu ấy hiểu được tâm ý của mình.

Kuon sau đó cũng lấy ra một món quà nhỏ trao cho tôi. ``Của cậu đây. Hy vọng giai điệu này cũng khiến cậu cảm thấy hạnh phúc.''

Tôi nhận lấy món quà, cảm thấy tâm hồn mình như đang bay bổng theo những nốt nhạc của niềm vui. Dưới ánh hoàng hôn vàng vọt tràn vào qua khung cửa sổ phòng nhạc, tôi biết rằng bản nhạc tuyệt vời nhất mà mình từng viết chính là khoảnh khắc chúng tôi cùng đứng bên nhau lúc này.

\begin{table}[H]
  \centering
  \setlength{\arrayrulewidth}{1pt}
  \begin{tabular}{|>{\centering\arraybackslash}p{0.4\linewidth}|>{\centering\arraybackslash}p{0.4\linewidth}|}
    \hline
    \hyperref[choice]{Quay trở về lựa chọn} & \hyperref[ending]{Đi tới phần cuối} \\ \hline
  \end{tabular}
\end{table}

% XV. Tokiwa Kana
\newpage
\label{kana}

Tôi chú ý đến một bức thư có màu tím đậm, sang trọng và đầy cá tính. Ngay nắp phong bì là một hình dán ác ma nhỏ với đôi cánh và cái đuôi tinh nghịch. Nhìn cái sticker này, tôi không khỏi bật cười khi nhớ tới sự tự luyến của Towa mỗi khi xưng danh ác ma ở thế giới thực, mà lần nào cũng toàn bị réo là thiên sứ.

Bên trong là nét chữ dứt khoát nhưng dường như có chút ý đồ mờ ám:

\txtbx{``Tan học, hãy ghé qua phòng phát thanh nhé. Kanade-chan đang có một thông báo cực kỳ quan trọng muốn nói với cậu đấy. Nhớ đến đúng giờ nhé!''\\
  \begin{flushright}-Tokiwa Kana-\end{flushright}}

Đến chỗ Kanade sao? Tôi hơi nhướng mày. Tại sao cô ấy lại viết thư rủ tôi đến chỗ người khác? Tính cách tốt bụng đến mức bao đồng này của cô ấy thật khiến người ta phải suy ngẫm. Tuy nhiên, thay vì đi thẳng đến phòng phát thanh, tôi quyết định đi vòng quanh trường để kiểm tra một chút, và linh cảm đã dẫn tôi đến dãy hành lang phía sau khu nhà đa năng.

\ct

\pov{Tokiwa Kana}

Tôi nấp sau góc hành lang, lén nhìn về phía cầu thang dẫn lên phòng phát thanh. ``Kế hoạch hoàn hảo! Chắc chắn Kuon-san đã đọc thư và đang trên đường đến chỗ Kanade-chan rồi. Hehe, mình quả là một ác ma tốt bụng mà!''

Dù tự nhủ như vậy, nhưng sâu thẳm trong lòng, tôi cảm thấy có chút nhói đau. Tự tay đẩy người mình thích về phía bạn thân\dots\ cảm giác này thật chẳng dễ chịu chút nào. Nhưng với tôi, hạnh phúc của bạn bè luôn quan trọng hơn tất thảy.

``\textit{Mình ổn mà\dots\ hoàn toàn ổn\dots}'' tôi thầm nhủ, khẽ siết chặt hộp quà nhỏ màu tím mà tôi vốn đã chuẩn bị cho cậu ấy, nhưng giờ có lẽ sẽ chẳng bao giờ được trao đi.

``Này, ác ma gì mà lại đứng thẫn thờ ở đây thế?''

Một giọng nói quen thuộc vang lên ngay sát bên cạnh khiến tôi giật bắn mình. Tôi quay phắt lại và thấy Kuon đang đứng đó, khoanh tay nhìn tôi với nụ cười đầy ẩn ý.

``K-K-Kuon-kun?! S-Sao cậu lại ở đây? Không phải cậu nên ở phòng phát thanh với Kanade-chan sao?'' tôi lắp bắp, lùi lại một bước, mặt nóng bừng lên vì bị bắt quả tang.

Kuon bật cười, bước lại gần hơn. ``Tôi đã đến đó rồi, nhưng Kanade-san bảo cô ấy chẳng hẹn gì tôi cả, và cũng đang chờ một người tên Kana-san đến để giúp chuẩn bị cho buổi phát sóng chiều nay. Có vẻ như ai đó đã cố tình nhầm lẫn chăng?''

Tôi cứng đờ người. ``Á\dots\ cái đó\dots\ mình\dots''

``Kana-san,'' Kuon cắt ngang, giọng nói trở nên dịu dàng hơn hẳn. ``Cậu không cần phải lúc nào cũng hy sinh bản thân mình như thế đâu. Tôi đến đây là vì muốn gặp người đã viết bức thư tím này cho mình.''

Nghe những lời ấy, bức tường phòng thủ mà tôi cố gắng dựng lên bấy lâu nay dường như sụp đổ hoàn toàn. Tôi cúi gầm mặt, đôi vai khẽ run lên. ``Cậu\dots\ thật là\dots\ chẳng giống ác ma chút nào cả\dots''

Tôi run rẩy đưa hộp quà ra trước mặt cậu ấy. ``Của cậu đây\dots\ mình đã vất vả chuẩn bị đấy. Đừng có mà chê đắng nhé!''

Kuon nhận lấy, nụ cười trên môi cậu ấy càng thêm rạng rỡ. Cậu ấy nếm thử và gật đầu tán thưởng: ``Tuyệt lắm. Vị đắng nhẹ quyện cùng vị ngọt hậu\dots\ rất giống tính cách của một ác ma nhưng lại có trái tim thiên thần như cậu vậy.''

Tôi ngước nhìn cậu ấy, nước mắt khẽ rưng rưng nhưng nụ cười đã trở lại trên môi. Valentine này, có lẽ tôi đã không thể trở thành một ác ma hoàn hảo như ý muốn, nhưng tôi đã nhận được điều ngọt ngào nhất mà mình hằng mong ước từ tận đáy lòng.

\begin{table}[H]
  \centering
  \setlength{\arrayrulewidth}{1pt}
  \begin{tabular}{|>{\centering\arraybackslash}p{0.4\linewidth}|>{\centering\arraybackslash}p{0.4\linewidth}|}
    \hline
    \hyperref[choice]{Quay trở về lựa chọn} & \hyperref[ending]{Đi tới phần cuối} \\ \hline
  \end{tabular}
\end{table}

% XVI. Karamomo Suzune
\newpage
\label{suzune}

Ánh mắt tôi lia tới một phong bì có màu cam vàng rực rỡ, trông vô cùng bắt mắt. Ngay giữa nắp phong bì là một hình dán chú hải cẩu trắng muốt đang nằm lăn kềnh, trông vừa ngộ nghĩnh vừa đáng yêu. Nét chữ bên trong tròn trịa, có chút nhí nhảnh:

\txtbx{``Kuon-kun ơi! Tan học ra sân sau trường gặp mình nhé! Có bất ngờ lớn cho cậu đây! Hehe!''\\
  \begin{flushright}-Karamomo Suzune-\end{flushright}}

Sân sau trường\dots\ nơi có những bụi cây rậm rạp và cũng là nơi Suzune hay chạy nhảy lung tung nhất. Tôi cất bức thư vào túi và mỉm cười, tự hỏi ``bất ngờ lớn'' của cô nàng linh vật lớp này sẽ là gì đây.

\ct

\pov{Karamomo Suzune}

Tôi đứng núp sau một gốc cây lớn, thỉnh thoảng lại ló đầu ra nhìn về phía lối đi. Trong tay tôi là một hộp quà lớn được bọc giấy màu cam rực rỡ, thắt nơ xanh lá cây - sự kết hợp màu sắc mà Azuki bảo là ``trông giống quả cà rốt hơn là hộp quà Valentine''.

``\textit{Hừm, Kuon-kun lâu quá đi mất! Hay là cậu ấy không thấy thư của mình?}'' tôi lẩm bẩm, chân không ngừng nhún nhảy vì phấn khích xen lẫn chút hồi hộp.

Đúng lúc đó, bóng dáng quen thuộc của Kuon xuất hiện. Tôi vội vàng nhảy ra khỏi chỗ nấp, giơ cao hộp quà và hét lớn:

``Kuon-kun! Ở đây, ở đây!''

Kuon giật mình khựng lại, rồi bật cười khi thấy điệu bộ của tôi. ``Ái chà, Suzune-san. Cậu đứng đây từ bao giờ thế?''

``Lâu lắm rồi đó nha!'' tôi chu môi, chạy lại gần cậu ấy và chìa hộp quà ra. ``Tada! Quà Valentine của cậu đây! Mình đã dành cả buổi sáng để 'nghiên cứu' và chế biến nó đấy!''

Cậu ấy nhận lấy hộp quà, nhìn nó một hồi rồi khẽ hỏi: ``Hộp quà này\dots\ trông có vẻ nặng hơn bình thường nhỉ?''

``Thì mình đã thêm vào đó rất nhiều 'tình yêu' và 'sự sáng tạo' mà!'' tôi cười hì hì, hối thúc cậu ấy: ``Mau mở ra ăn thử đi! Mình muốn biết cảm nghĩ của cậu lắm!''

Kuon từ từ mở nắp hộp. Bên trong là những viên sô-cô-la có hình thù cực kỳ đa dạng, từ hình đầu gấu, hải cẩu cho đến cả\dots\ hình mấy miếng trái cây. Cậu ấy chọn một viên có màu hơi sậm hơn bình thường và cho vào miệng.

Tôi nín thở, hai mắt mở to chờ đợi.

Vừa nhai được vài cái, lông mày của Kuon bỗng nhíu lại. Cậu ấy khựng người mất vài giây, rồi khẽ thốt lên: ``\dots Vị này\dots\ ăn kỳ kỳ thế nào ấy.''

``Hả?! Kỳ sao?!'' tôi hốt hoảng, cuống cuồng giật lấy hộp quà. ``Nó đắng quá à? Hay là quá ngọt? Hay là mình lỡ tay cho nhầm muối? Thôi chết rồi, để mình nếm thử xem nào!''

Tôi định lấy một viên định ăn thì Kuon ngăn lại. Cậu ấy mỉm cười trấn an, dù vẻ mặt vẫn có chút gì đó rất khó tả: ``Không đến mức khó ăn đâu\dots\ Chỉ là, nó có vị rất lạ. Giống như là\dots\ vị sô-cô-la trộn với\dots\ mù tạt và mật ong, chắc vậy?''

Tôi gãi gãi đầu, cười trừ: ``A\dots\ thực ra là mình có thử nghiệm một chút công thức mới. Mình nghĩ sự kết hợp giữa cay và ngọt sẽ tạo ra một cú nổ về hương vị!''

Kuon phì cười, nỗ lực nuốt trọn viên sô-cô-la ``độc bản'' ấy. Cậu ấy lau nhẹ khóe miệng, rồi bỗng nhiên nhìn tôi với ánh mắt đầy nghi hoặc: ``Suzune-san, cái sô-cô-la này\dots\ cậu không cho cái 'chất lỏng màu xanh' cậu định mang ra đấy chứ?''

Tôi giật bắn mình, lắp bắp: ``H-Hả? Cái lọ lấp lánh lượm được ở phòng hóa học ấy hả? Không có đâu! Azuki đã thu giữ nó như thể đó là tang vật vụ án rồi còn đâu! Mình chỉ kịp\dots\ ừm\dots\ ngắm nó một chút thôi mà!''

Kuon thở phào nhẹ nhõm, vỗ nhẹ lên vai tôi: ``May quá. Tôi cứ tưởng mình sắp biến thành quái vật của Frankenstein đến nơi rồi.''

``Cậu thật là\dots!'' tôi bĩu môi, nhưng lòng thầm vui sướng khi thấy cậu ấy vẫn vui vẻ dù món quà của mình có chút kỳ quặc.

Cậu ấy sau đó lấy từ trong cặp ra một hộp sô-cô-la nhỏ nhắn, gói ghém cực kỳ xinh xắn đưa cho tôi. ``Đây là phần của tôi dành cho cậu. Hy vọng nó không 'kỳ quặc' như của cậu.''

Tôi nhận lấy món quà, cảm giác hạnh phúc trào dâng lấp đầy cả lồng ngực. Tôi nhanh chóng mở ra và nếm thử. Vị ngọt lịm và béo ngậy tan chảy trên đầu lưỡi, ngon đến mức tôi muốn hét lên vì sung sướng.

``Ngon quá đi mất! Kuon-kun đúng là nhất luôn!''

Tôi không kìm được mà nhảy lên ôm chầm lấy cánh tay cậu ấy, khiến Kuon có chút lúng túng nhưng vẫn mỉm cười dịu dàng. Giữa sân sau trường ngập tràn ánh hoàng hôn, tiếng cười của tôi vang vọng, xua tan đi cả cái lạnh của buổi chiều muộn. Valentine này, dù sô-cô-la có vị kỳ quặc một chút, nhưng đối với tôi, nó chính là hương vị tuyệt vời nhất thế gian.

\begin{table}[H]
  \centering
  \setlength{\arrayrulewidth}{1pt}
  \begin{tabular}{|>{\centering\arraybackslash}p{0.4\linewidth}|>{\centering\arraybackslash}p{0.4\linewidth}|}
    \hline
    \hyperref[choice]{Quay trở về lựa chọn} & \hyperref[ending]{Đi tới phần cuối} \\ \hline
  \end{tabular}
\end{table}

% XVII. Nekomiya Yuka
\newpage
\label{yuka}

Tôi cầm lấy một bức thư có màu tím oải hương dịu mắt. Điểm nhấn trên đó là một hình dán cơm nắm (onigiri) bé xíu vô cùng dễ thương, khiến tôi nhớ tới Okayu ở thế giới thực. Nét chữ bên trong uốn lượn một cách thong thả, mang lại cảm giác rất thư thái:

\txtbx{``Kuon-kun này, tan học nếu không bận thì lên sân thượng ở dãy nhà B nha\~{} Mình đợi ở đó nha\~{}''\\
  \begin{flushright}-Nekomiya Yuka-\end{flushright}}

Sân thượng\dots\ một nơi lý tưởng để hóng gió và\dots\ ngủ. Tôi cất bức thư vào túi và thong thả đi tới khu nhà bên cạnh, rồi lên tầng cao nhất của tòa nhà.

\ct

\pov{Nekomiya Yuka}

``Mmm\dots\ Oáp\~{}''

Tôi khẽ vươn vai, chớp mắt nhìn bầu trời đang dần chuyển sang màu hổ phách. Không khí sân thượng vào lúc chiều muộn thật dễ chịu, đủ mát mẻ để làm một giấc ngủ ngắn chất lượng.

Tiếng cửa sắt khô khốc vang lên, rồi tiếng bước chân đều đặn tiến lại gần. Tôi vươn vai lên cao, miệng khẽ mỉm cười, dù vẫn còn hơi ngái ngủ.

``Chào Kuon-kun\~{} Cậu đến rồi sao? Tôi tưởng mình đã ngủ quên rồi chứ.''

Kuon bước đến bên cạnh tôi, nhìn quanh một lượt rồi bật cười. ``Cậu ngủ ở đây thật à? Không sợ bị cảm sao?''

``Tại ở đây yên tĩnh quá mà\~{}'' tôi dụi dụi mắt, nhìn cậu ấy với vẻ tinh nghịch.

``Cậu định hỏi tôi đang làm gì ở đây sao?'' tôi hỏi ngược lại khi thấy ánh mắt tò mò của cậu ấy. ``À\dots\ thực ra tôi đang đợi một người đặc biệt\dots\ để đưa cái này.''

Tôi lấy từ trong cặp sách ra một hộp quà nhỏ, bọc giấy tím đơn giản nhưng tinh tế.

``Đoán xem là ai nào\~{}''

Kuon nhướng mày, rồi chỉ tay vào chính mình: ``Là tôi sao?''

``Hehe\~{} đúng rồi, là cậu đó,'' tôi chìa hộp quà ra trước mặt cậu ấy. ``Tôi đã làm nó cho cậu\dots\ dù có thể nó trông hơi giống cơm nắm một chút. Akane bảo tôi là đồ cuồng onigiri, nhưng tôi nghĩ hình dáng này cầm rất chắc tay và dễ ăn mà, đúng không?''

Kuon đón lấy món quà, nụ cười rạng rỡ hiện rõ trên môi. Cậu ấy mở hộp, lấy ra một viên sô-cô-la được tạo hình đúng chuẩn một miếng onigiri thu nhỏ, hoàn chỉnh với cả một miếng sô-cô-la đen bọc ngoài như miếng rong biển.

``Cảm ơn cậu, Yuka-san. Nó thực sự rất sáng tạo.'' cậu ấy nếm thử và gật đầu tán thưởng. ``Vị ngọt dịu, lại có chút gì đó béo ngậy của kem tươi\dots\ Rất tuyệt.''

Thấy cậu ấy hài lòng, tôi cảm thấy trong lòng dâng lên một sự ấm áp lạ thường. Ánh mắt tôi vô tình lướt qua Akane đang lấp ló sau cánh cửa dẫn lên sân thượng, chắc hẳn cậu ấy cũng vừa kết thúc buổi hẹn của mình và đang lo lắng cho tôi.

Kuon sau đó cũng lấy ra một hộp sô-cô-la đáp lễ. ``Đây là của cậu. Hy vọng nó sẽ giúp cậu tỉnh táo hơn sau những giấc ngủ gật trong lớp.''

Tôi nhận lấy, cảm nhận hơi ấm từ món quà truyền qua lòng bàn tay. ``Cảm ơn Kuon-kun nhé. Tôi sẽ trân trọng nó.''

Dưới ánh hoàng hôn vàng vọt của sân thượng, chúng tôi cùng nhau đứng tựa vào lan can, im lặng ngắm nhìn khung cảnh thị trấn Aogami đang dần lên đèn. Không cần quá nhiều lời nói, sự hiện diện của cậu ấy bên cạnh cũng đủ khiến tôi cảm thấy bình yên và hạnh phúc nhất rồi.

\begin{table}[H]
  \centering
  \setlength{\arrayrulewidth}{1pt}
  \begin{tabular}{|>{\centering\arraybackslash}p{0.4\linewidth}|>{\centering\arraybackslash}p{0.4\linewidth}|}
    \hline
    \hyperref[choice]{Quay trở về lựa chọn} & \hyperref[ending]{Đi tới phần cuối} \\ \hline
  \end{tabular}
\end{table}

% XVIII. Akagane Mari
\newpage
\label{mari}

Ánh mắt tôi dừng lại ở một phong bì màu đỏ thẫm vô cùng quý phái. Nổi bật trên nắp phong bì là một hình dán bảng màu nhỏ xíu với đủ các sắc màu rực rỡ. Nét chữ bên trong nắn nót, mang đậm phong cách của một người nghệ sĩ:

\txtbx{``Sau giờ học, hãy ghé qua phòng Mỹ thuật một chút nhé. Tôi có thứ này muốn cho cậu xem.''\\
  \begin{flushright}-Akagane Mari-\end{flushright}}

Phòng Mỹ thuật\dots\ nơi lưu giữ những giấc mơ hội họa của Mari-san\dots\ và cũng là nơi mà cô nàng đã từng bỏ mạng chỉ để hoàn thành cái ``kiệt tác'' quái gở ấy, ít nhất là ở vòng lặp trước. Tôi cất bức thư vào túi và bước về phía dãy nhà chuyên dụng, nơi ánh nắng chiều tà thường tạo nên những mảng màu tuyệt đẹp trên những bức tường cũ.

\ct

\pov{Akagane Mari}

Tôi ngồi trước giá vẽ, tay cầm bảng màu và cây cọ nhỏ, chăm chú dặm thêm những vệt màu đỏ rực vào bức tranh phong cảnh đang dang dở. Mùi dầu pha màu quyện cùng hương gỗ cũ của căn phòng tạo nên một không gian tĩnh lặng, tách biệt hẳn với sự ồn ào bên ngoài.

Trong tâm trí tôi, những ký ức về cơn đau đầu dữ dội và những tháng ngày mất đi lý trí thỉnh thoảng vẫn hiện về như một bóng ma mờ nhạt. Nhưng khi cầm cây cọ này lên, tôi lại cảm thấy mình được sống lại một lần nữa. Tôi muốn vẽ một bức tranh có thể chạm đến trái tim của mọi người, một bức tranh không gì có thể sánh bằng.

Tiếng cửa phòng khẽ mở. Tôi quay lại và thấy Kuon đang đứng đó, ánh mắt lộ rõ vẻ ngưỡng mộ khi nhìn vào bức tranh ``Mari-san, bức tranh này\dots\ sống động quá. Cảm giác như màu sắc đang thực sự nhảy múa vậy.''

Tôi khẽ cười, đặt bảng màu xuống và lau nhẹ bàn tay dính màu vào chiếc tạp dề. ``Cảm ơn cậu, Kuon-kun. Cậu luôn là người đầu tiên hiểu được những gì mình muốn truyền tải qua từng nét vẽ.''

Tôi bước lại gần chiếc bàn gỗ lớn ở góc phòng, nơi đặt một hộp quà nhỏ được bọc giấy đỏ thẫm, trang trí bằng một chiếc nơ nhỏ làm từ dây thừng thô.

``Hôm nay là Lễ tình nhân\dots\ Tôi không giỏi nấu nướng như vẽ tranh, nhưng tôi đã cố gắng hết sức để tạo ra những 'viên màu' đặc biệt này cho cậu đó.''

Tôi chìa hộp quà ra trước mặt cậu ấy. ``Hy vọng cậu sẽ thích 'tác phẩm' này.''

Kuon nhận lấy món quà, nụ cười rạng rỡ hiện rõ trên môi. Cậu ấy mở hộp, và bên trong là những viên sô-cô-la được tạo hình giống như những tuýp màu nhỏ nhắn, mỗi tuýp lại có một màu sắc khác nhau từ lớp phủ trái cây tự nhiên.

``Tuyệt quá! Trông chúng giống hệt như những món đồ trong phòng Mỹ thuật vậy.'' cậu ấy chọn một viên màu đỏ dâu tây và nếm thử. ``Vị ngọt ngào pha chút chua nhẹ\dots\ Rất tinh tế và đầy cảm hứng, Mari-san.''

Nghe lời khen ấy, tôi cảm thấy một luồng hơi ấm lan tỏa trong lồng ngực, còn hơn cả niềm vui khi hoàn thành một bức tranh ưng ý.

Kuon sau đó cũng lấy ra một món quà nhỏ trao cho tôi. ``Đây là phần của tôi dành cho cậu. Mong rằng nó sẽ tiếp thêm màu sắc cho những bức vẽ tiếp theo của cậu.''

Tôi nhận lấy món quà, cảm thấy trái tim mình như vừa được tô thêm một mảng màu rực rỡ nhất. Dưới ánh nắng hoàng hôn len lỏi qua khung cửa sổ phòng Mỹ thuật, tôi biết rằng bức tranh đẹp nhất mà mình từng vẽ không nằm trên giá kia, mà chính là khoảnh khắc chúng tôi cùng đứng bên nhau lúc này, giữa hương vị sô-cô-la và tình cảm chân thành nhất.

\begin{table}[H]
  \centering
  \setlength{\arrayrulewidth}{1pt}
  \begin{tabular}{|>{\centering\arraybackslash}p{0.4\linewidth}|>{\centering\arraybackslash}p{0.4\linewidth}|}
    \hline
    \hyperref[choice]{Quay trở về lựa chọn} & \hyperref[ending]{Đi tới phần cuối} \\ \hline
  \end{tabular}
\end{table}

% XIX. Kisaragi Saya
\newpage
\label{saya}

Tôi chọn lấy một phong bì màu vàng tươi tắn, gợi cảm giác ấm áp như những tia nắng mặt trời cuối ngày. Nổi bật trên nắp phong bì là một hình dán chòm sao nhỏ xíu lấp lánh. Nét chữ bên trong uốn lượn, thong thả:

\txtbx{``Kuon-kun nè\~{} Nếu không bận thì hãy ra khu đất trống gần trường nhé\~{} Mình có chuyện muốn nói cùng cậu nè\~{}''\\
  \begin{flushright}-Kisaragi Saya-\end{flushright}}

Khu đất trống gần trường à\dots\ Đó cũng là nơi mà phần lớn hội thao lần trước đã tổ chức gần đó. Ngoài ra, nơi đó cũng rất thoáng đãng và là địa điểm lý tưởng để ngắm bầu trời đang dần chuyển mình. Tôi cất bức thư vào túi và thong thả bước đi, trong đầu thầm tưởng tượng ra điệu bộ nói chuyện chậm rãi quen thuộc của cô nàng yêu thiên văn này.

\ct

\pov{Kisaragi Saya}

Tôi đứng bên khu đất trống, ngửa cổ nhìn lên bầu trời đang nhuộm một màu tím sẫm xen lẫn những vệt hồng tà dương. Thời tiết hôm nay cũng hoàn hảo để có thểm ngắm bầu trời đêm. Những ngôi sao đầu tiên đang bắt đầu lấp ló hiện ra, như những viên kim cương nhỏ xíu được rải trên một tấm thảm nhung vô tận.

Tôi vốn rất yêu bầu trời đêm. Từ khi còn nhỏ, tôi đã luôn mơ mộng về những hành tinh xa xôi và những câu chuyện thần thoại ẩn sau mỗi chòm sao. Đối với tôi, mỗi vì sao đều mang trong mình một tâm sự riêng, đang thầm thì kể cho những ai chịu lắng nghe.

Nghe tiếng bước chân lạo xạo tiến lại gần, tôi khẽ quay lại, nở một nụ cười thật dịu dàng.

``Kuon-kun đã đến rồi nè\~{} Cảm ơn cậu đã dành thời gian cho mình nha\~{}''

Kuon bước tới bên cạnh tôi, cũng nhìn lên bầu trời. ``Saya-san đến từ trước rồi cơ à. Cậu lại đang ngắm sao sao?''

``Ừm nè\~{} Cậu nhìn kìa, chòm sao Thiên Lang đang bắt đầu tỏa sáng rồi đó\~{} Trông nó thật kiêu sa phải không nào\~{}'' Tôi say sưa chỉ tay lên bầu trời, bắt đầu kể cho cậu ấy nghe về những kiến thức thiên văn mà mình hằng ấp ủ.

Kuon chăm chú lắng nghe, thỉnh thoảng lại gật đầu tán thưởng. Sự hiện diện điềm tĩnh của cậu ấy khiến tôi cảm thấy bình yên đến lạ, cứ như thể chúng tôi đang trôi lửng lơ giữa dải ngân hà vậy.

Mải nói về những chòm sao trên trời, tôi chợt giật mình: ``À, suýt nữa thì mình quên mất việc chính rồi nè\~{}'' Tôi khẽ đỏ mặt, lấy từ trong chiếc túi nhỏ ra một hộp quà màu vàng được thắt nơ lấp lánh như bụi sao.

``Cái này\dots\ là món quà mình tự tay chuẩn bị cho cậu đó\~{} Mong cậu hãy nhận lấy tâm ý này của mình nha\~{}''

Kuon đón lấy món quà, nụ cười rạng rỡ hiện rõ trên môi. ``Cảm ơn cậu, Saya-san. Một chiếc hộp quà thật đẹp, trông như chứa cả một bầu trời sao bên trong vậy.''

Cậu ấy mở hộp và nếm thử. Ánh mắt cậu ấy khẽ thay đổi đầy thích thú. ``Vị ngọt thật nhẹ nhàng và thanh khiết\dots\ giống như cảm giác khi được ngắm nhìn bầu trời đêm cùng cậu vậy.''

Nghe lời khen ấy, trái tim tôi như vừa được một vì sao băng lướt qua, để lại những vệt sáng lấp lánh của niềm hạnh phúc.

Kuon sau đó cũng lấy ra một món quà nhỏ trao cho tôi. ``Của cậu đây. Hy vọng nó sẽ là `ngôi sao' ngọt ngào nhất trong bộ sưu tập của cậu ngày hôm nay.''

Tôi nhận lấy món quà, cảm thấy tâm hồn mình như đang bay bổng giữa những tầng mây. Dưới bầu trời chạng vạng đang dần hiện rõ những chòm sao, tôi biết rằng khoảnh khắc đẹp nhất không nằm ở những thiên thể xa xôi kia, mà chính là sự ấm áp khi được đứng cạnh cậu ấy lúc này.

``Valentine này, đối với mình, thực sự rất diệu kỳ luôn đó\~{}''

\begin{table}[H]
  \centering
  \setlength{\arrayrulewidth}{1pt}
  \begin{tabular}{|>{\centering\arraybackslash}p{0.4\linewidth}|>{\centering\arraybackslash}p{0.4\linewidth}|}
    \hline
    \hyperref[choice]{Quay trở về lựa chọn} & \hyperref[ending]{Đi tới phần cuối} \\ \hline
  \end{tabular}
\end{table}

% XX. Shitou Shion
\newpage
\label{shion}

Ánh mắt tôi bị thu hút bởi một phong bì màu đen tuyền, nổi bật trên đó là một chiếc sticker hình mũ phù thủy nhỏ xíu nhưng đầy bí ẩn. Nét chữ bên trong uốn lượn một cách kỳ lạ, như thể được viết bằng một loại mực ma thuật nào đó:

\txtbx{``Hỡi phàm nhân đang bị xiềng xích bởi định mệnh! Nếu muốn giải phóng linh hồn khỏi vòng lặp vô tận này, hãy tìm đến nơi cao nhất của tòa tháp tri thức khi tà dương nhuộm đỏ bầu trời. Ta sẽ đợi ngươi ở đó.''\\
  \begin{flushright}-Shitou Shion-\end{flushright}}

``Tòa tháp tri thức'' à\dots\ Thư viện chứ gì?

Tôi mỉm cười trước phong cách đặc trưng của cô bạn tự xưng là phù thủy đen này. Cất bức thư vào túi, tôi rảo bước về phía thư viện trường - nơi giờ này chắc hẳn đã không còn ai ngoài cô nàng đang chìm đắm trong thế giới giả tưởng của mình.

\ct

\pov{Shitou Shion}

Ta đứng tựa lưng vào giá sách cổ kính, đôi mắt dõi nhìn qua khung cửa sổ về phía đường chân trời đang rực cháy ánh hoàng hôn. Trong tay ta là một chiếc hộp đen huyền bí, được phong ấn bằng những sợi dây thừng ma thuật (thực chất là ruy băng lụa màu tím).

``Hừm, tên phàm nhân kia lại dám để một phù thủy quyền năng như ta phải chờ đợi sao?'' ta lẩm bẩm, bàn tay vô thức siết chặt chiếc hộp.

Lồng ngực ta bỗng cảm thấy thắt lại một cách kỳ lạ. Dù ta luôn tự xưng là Phù thủy xứ Vương Quốc Bóng Tối Vĩnh Cửu, là thực thể đứng ngoài mọi quy luật của thế giới này, nhưng đứng trước khoảnh khắc này, tim ta lại đập loạn nhịp như một nữ sinh bình thường. Thật là một loại ma thuật hắc ám đáng sợ!

Tiếng cửa thư viện khẽ rít lên khô khốc. Bước chân đều đặn vang vọng trong không gian tĩnh mịch. Là ngươi, Kuon!

Ta xoay người lại, khoanh tay trước ngực, cố gắng lấy lại vẻ uy nghiêm vốn có. Đôi môi ta khẽ nhếch lên một nụ cười đầy bí ẩn.

``Cuối cùng ngươi cũng đã đến, kẻ mang trong mình dòng máu của những vì sao bị lãng quên! Ngươi có biết việc để một phù thủy phải chờ đợi sẽ dẫn đến thảm cảnh gì không?''

Kuon khẽ bật cười, nụ cười hiền hậu của ngươi dường như đang cố tình phá tan bầu không khí huyền bí mà ta đã dày công xây dựng. ``Xin lỗi vì đã để cậu chờ, Shion-san. Chỉ là cầu thang hôm nay hơi dốc chút thôi.''

``Hừ, đó là do chướng ngại vật của định mệnh sắp đặt để thử thách ý chí của ngươi mà thôi!'' ta hất hàm, che giấu đi sự bối rối đang dâng trào trong lòng.

Kuon tiến lại gần, nhìn vào chiếc hộp trên tay ta với vẻ tò mò. ``Bức thư của cậu thú vị thật đấy. 'Tháp tri thức' cơ à?''

``Đương nhiên! Nơi này chứa đựng vô vàn bí mật cổ xưa mà phàm nhân như ngươi khó lòng thấu hiểu,'' ta hắng giọng, rồi lấy hết can đảm, đưa chiếc hộp đen về phía ngươi.

``Cầm lấy đi! Đây là tinh hoa của bóng tối mà ta đã tôi luyện suốt sáng nay. Bên trong nó chứa đựng một loại bùa chú đặc biệt giúp ngươi tăng cường 'mana' để chống chọi với thế giới đầy rẫy những biến số này.''

Kuon nhận lấy món quà, đôi mắt ngươi lấp lánh sự thích thú. ``Cảm ơn cậu, Shion-san. Đây là sô-cô-la cậu tự làm sao?''

``Đừng có gọi bằng cái tên tầm thường đó! Đó là 'Hắc ngọc của sự vĩnh hằng'!'' ta vội vàng lảng tránh ánh mắt của ngươi, cảm thấy đôi tai mình đang nóng bừng lên. ``Ta đã vất vả pha trộn tỷ lệ cacao hắc ám để tạo ra hương vị hoàn hảo nhất đấy. Ngươi mà dám chê đắng thì ta sẽ nguyền rủa ngươi suốt đời!''

Kuon khẽ mỉm cười, mở hộp và thưởng thức một viên. Ngươi nhắm mắt lại, rồi gật đầu tán thưởng: ``Tuyệt quá\dots\ Vị đắng rất đậm đà nhưng lại có hậu vị ngọt ngào sâu lắng. Đúng là món quà của một phù thủy quyền năng có khác.''

Nghe những lời khen ngợi ấy, trái tim ta bỗng rung động mãnh liệt. Dù ta luôn miệng nói về bóng tối và sự cô độc, nhưng ánh sáng ấm áp phát ra từ nụ cười của ngươi lại là thứ ma thuật duy nhất có thể xuyên thấu qua lớp giáp bảo vệ của ta.

Kuon sau đó cũng lấy ra một món quà nhỏ trao cho ta. ``Của cậu đây. Hy vọng nó sẽ là 'linh vật' giúp cậu thực hiện được những phép thuật vĩ đại nhất.''

Ta nhận lấy món quà, cảm giác hạnh phúc trào dâng lấp đầy cả lồng ngực. Ta nhanh chóng nếm thử, vị ngọt ngào tan chảy khiến ta không tự chủ được mà mỉm cười rạng rỡ.

``Ưm\dots\ không tệ! Ta rất hài lòng với vật phẩm này.'' Ta cố giữ giọng mình thật lạnh lùng, dù trong lòng đang muốn nhảy cẫng lên vì sung sướng.

Dưới ánh hoàng hôn vàng vọt tràn vào qua khung cửa sổ thư viện, chúng tôi cùng nhau đứng đó, giữa những giá sách cũ kỹ và hương vị sô-cô-la ngọt ngào. Valentine này, dù thế giới có xoay chuyển ra sao, dù vòng lặp có bắt đầu lại, ta cũng sẽ mãi ghi nhớ khoảnh khắc này - khoảnh khắc mà một phù thủy và một phàm nhân cùng chung bước dưới bầu trời đầy sao của định mệnh.

``Ngươi phải biết ơn ta đấy nhé, phàm nhân! Vì ta đã chọn ngươi là kẻ đồng hành cùng ta trong hành trình tìm kiếm sự thật vĩnh hằng này!''

\begin{table}[H]
  \centering
  \setlength{\arrayrulewidth}{1pt}
  \begin{tabular}{|>{\centering\arraybackslash}p{0.4\linewidth}|>{\centering\arraybackslash}p{0.4\linewidth}|}
    \hline
    \hyperref[choice]{Quay trở về lựa chọn} & \hyperref[ending]{Đi tới phần cuối} \\ \hline
  \end{tabular}
\end{table}

% XXI. Shiratori Hotaru
\newpage
\label{hotaru}

Tôi chọn lấy một phong bì có sự kết hợp màu sắc khá mạnh mẽ: màu cam rực rỡ hòa quyện cùng sắc xanh nước biển đậm đà. Ở giữa nắp phong bì là một chiếc nơ xanh nước biển được dán ngay ngắn, gợi nhớ đến phong cách thời trang của Shiratori Hotaru.

Nét chữ bên trong thẳng tắp và tràn đầy năng lượng:

\txtbx{``Kuon-kun nè! Tan học cậu có rảnh không? Nếu được thì ghé qua căng-tin trường nhé, mình có thứ này muốn tặng cậu!''\\
  \begin{flushright}-Shiratori Hotaru-\end{flushright}}

Căng-tin trường ư? Chắc hẳn giờ này nơi đó đang tràn ngập mùi hương của những món ăn vặt và tiếng cười nói của học sinh. Tôi cất bức thư vào túi và rảo bước về phía khu vực ăn uống.

\ct

\pov{Shiratori Hotaru}

Tôi ngồi ở một góc căng-tin, đôi chân khẽ đung đưa dưới gầm bàn. Trên mặt bàn là một hộp quà được thắt nơ vàng tinh tế. Bình thường tôi là người khá tự tin và hay nói, nhưng không hiểu sao lúc này tim tôi lại đập thình thịch như muốn nhảy ra ngoài.

``\textit{Cố lên nào Hotaru, chỉ là tặng sô-cô-la thôi mà!}'' tôi thầm tự nhủ, tay vân vê gấu áo.

``Hotaru-san!''

Tiếng gọi khiến tôi giật mình ngẩng lên. Kuon đang đi tới, trên môi là nụ cười hiền hậu khiến bao nhiêu lo lắng trong tôi dường như tan biến sạch.

``A, Kuon-kun! Cậu đến rồi à?'' Tôi cười tươi, vẫy vẫy tay ra hiệu cho cậu ấy ngồi xuống đối diện.

``Xin lỗi vì đã để cậu đợi. Căng-tin giờ này vẫn nhộn nhịp nhỉ,'' Kuon ngồi xuống, ánh mắt lướt qua hộp quà trên bàn.

``Không sao, không sao! Tôi cũng vừa mới đến thôi mà,'' tôi xua xua tay, rồi hít một hơi thật sâu để lấy lại sự bình tĩnh thường ngày. ``À\dots\ thật ra thì, tôi có cái này muốn tặng cậu nè!''

Tôi đưa hộp quà về phía cậu ấy bằng cả hai tay, ánh mắt lộ rõ vẻ mong chờ. ``Lễ tình nhân mà, đúng không? Tôi đã tự tay làm nó đấy. Dù tôi hay bị Kana mắng là nói quá nhiều, nhưng khi làm sô-cô-la này, tôi đã dành hết tâm trí vào nó để không bị sai sót gì cả!''

Kuon ngạc nhiên đón lấy món quà. ``Cảm ơn cậu, Hotaru-san. Một chiếc hộp thật đẹp.''

Cậu ấy mở hộp, bên trong là những viên sô-cô-la được tạo hình những chú gấu nhỏ và những bông hoa xinh xắn. Cậu ấy nếm thử một viên và khẽ gật đầu tán thưởng: ``Ngon quá! Vị ngọt rất đậm đà nhưng không hề bị gắt. Cậu khéo tay thật đấy, Hotaru-san.''

Nghe thấy lời khen ấy, tôi cảm thấy mặt mình nóng bừng lên vì sung sướng. ``Thật sao? Cậu thích là tôi vui rồi! Hehe, tôi đã phải thử nghiệm mấy lần mới ra được công thức này đó nha!''

Sự tự tin trở lại, tôi lại bắt đầu liến thoắng kể cho cậu ấy nghe về quá trình làm sô-cô-la, về việc tôi đã phải cẩn thận thế nào để không làm cháy cacao. Kuon vẫn kiên nhẫn lắng nghe, thỉnh thoảng lại mỉm cười và đặt ra những câu hỏi khiến cuộc trò chuyện càng thêm sôi nổi.

Cậu ta sau đó cũng lấy từ trong cặp ra một món quà nhỏ trao cho tôi. ``Đây là phần của tôi dành cho cậu. Hy vọng nó sẽ mang lại cho cậu nhiều năng lượng như chính tính cách của cậu vậy.''

Tôi ngạc nhiên đón lấy món quà, cảm giác hạnh phúc lan tỏa khắp cơ thể. ``Woah! Của cậu làm sao? Cảm ơn nhiều nhé, Kuon-kun!''

Dưới không gian ấm cúng của căng-tin trường, chúng tôi cùng nhau trò chuyện và thưởng thức sô-cô-la. Dù chung quanh có ồn ào đến đâu, tôi vẫn cảm thấy thế giới này lúc này chỉ có hai chúng tôi. Valentine này, đối với một người thích nói như tôi, điều tuyệt vời nhất chính là có một người luôn sẵn lòng lắng nghe và chia sẻ những khoảnh khắc ngọt ngào như thế này.

``Nè, ăn xong thì tôi đi dạo một chút nhé? Tôi còn nhiều chuyện muốn kể cho cậu nghe lắm đó!'' tôi cười rạng rỡ, trong lòng thầm cảm ơn ngày Lễ tình nhân đã mang đến cơ hội tuyệt vời này.

\begin{table}[H]
  \centering
  \setlength{\arrayrulewidth}{1pt}
  \begin{tabular}{|>{\centering\arraybackslash}p{0.4\linewidth}|>{\centering\arraybackslash}p{0.4\linewidth}|}
    \hline
    \hyperref[choice]{Quay trở về lựa chọn} & \hyperref[ending]{Đi tới phần cuối} \\ \hline
  \end{tabular}
\end{table}

% XXII. Kamishiro Miku
\newpage
\label{miku}

Phong bì tiếp theo tôi chọn có màu đỏ thẫm phối với đen tuyền đầy vẻ huyền bí. Một hình dán con sói đen nhỏ nhắn ở góc thư khiến tôi nhớ ngay đến Miku-san. Nét chữ bên trong nắn nót, mang theo một phong thái vô cùng điềm tĩnh:

\txtbx{``Sau giờ học, hãy ghé qua khu phố truyền thống phía nam thị trấn nhé. Có một tiệm trà cũ mà mình rất thích ở đó, mình sẽ đợi cậu tại sảnh chờ.''\\
  \begin{flushright}-Kamishiro Miku-\end{flushright}}

Khu phố truyền thống à? Nó nằm khá xa trường. Và vấn đề là, nếu tôi cứ thế đi ra bằng cổng chính, rất có thể sẽ chạm mặt những cô gái khác - những người mà tôi chưa kịp hồi đáp bức thư của họ. Một cuộc vây ráp lúc này chắc chắn sẽ là thảm họa.

Tôi quyết định đi đường vòng. Lẻn ra phía sau nhà thể chất, tôi lách qua dãy hàng rào sắt đang hé mở một khoảng nhỏ. Trước mặt tôi là khu rừng thưa nằm sát cạnh trường. Tôi men theo lối mòn nhỏ len lỏi giữa những gốc cây già cỗi, tiếng lá khô xào xạc dưới chân. Sau khoảng mười phút đi bộ, tôi đã ra được con đường chính dẫn thẳng tới khu phố cổ mà không bị bất cứ ánh mắt hình viên đạn nào bắt gặp.

\ct

\pov{Kamishiro Miku}

Tôi ngồi trên băng ghế gỗ dài ở sảnh tiệm trà truyền thống, trên người vẫn là bộ đồng phục trường nhưng đã khoác thêm một chiếc cardigan và chiếc áo blazer để trông đứng đắn hơn. Tôi khẽ lật một trang tạp chí thời trang, nhưng thực tế là nãy giờ tôi chưa đọc được chữ nào cả.

Bên cạnh tôi là một hộp quà bọc giấy đỏ phối đen vô cùng sang trọng. Tôi đã đọc trong một bài báo gần đây rằng sô-cô-la đen có tác dụng rất tốt trong việc làm đẹp da và giữ gìn nét thanh xuân. Là một người luôn hướng tới hình ảnh người phụ nữ trưởng thành, tôi nghĩ món quà này là sự lựa chọn không thể hoàn hảo hơn.

Tiếng chuông gió treo trước cửa tiệm khẽ reo lên. Kuon bước vào, trên vai vẫn còn vương một vài mẩu lá khô nhỏ.

``Xin lỗi vì đã để cậu đợi lâu, Miku-san. Đường đi có hơi\dots\ chút trắc trở,'' cậu ta vừa thở dốc vừa cười khổ.

Tôi khép cuốn tạp chí lại, mỉm cười nhẹ nhàng, cố gắng giữ cho giọng điệu thật thanh thoát và chững chạc: ``Không sao đâu. Một người phụ nữ trưởng thành luôn biết cách kiên nhẫn chờ đợi những điều xứng đáng mà. Cậu ngồi đi.''

Kuon ngồi xuống đối diện tôi. Cậu ấy đưa tay gạt một mảnh lá khô trên tóc, rồi nhìn tôi với vẻ ngưỡng mộ: ``Trông cậu hôm nay có vẻ điềm đạm hơn hẳn ở lớp nhỉ, Miku-san.''

Nghe lời khen ấy, lồng ngực tôi bỗng nhói lên một nhịp. Tôi hắng giọng, đưa hộp quà về phía cậu ấy: ``Cái này\dots\ tặng cậu. Đây là loại sô-cô-la cao cấp mà mình đã dày công chọn lọc. Nó có hàm lượng cacao rất cao, cực kỳ tốt cho sức khỏe và làn da đó.''

``Cảm ơn cậu, Miku-san. Cậu chu đáo thật đấy,'' Kuon mở hộp quà, lấy một viên nếm thử. ``Ưm, vị đắng rất thanh và sang trọng. Đúng là phong cách của cậu.''

Nhìn cậu bạn thưởng thức món quà của mình, vẻ điềm tĩnh mà tôi cố công xây dựng bỗng dưng tan biến sạch. Khi cậu ấy nhìn thẳng vào mắt tôi và nói câu: ``Cậu thực sự rất tuyệt vời, Miku-san,'' tôi cảm thấy mặt mình nóng bừng lên.

``H-Hả? T-Thật sao? Cậu thấy mình tuyệt vời sao?'' Giọng tôi bỗng cao vút lên, mất đi vẻ trầm ổn ban nãy. ``Hehe, Kuon-kun quá khen rồi\dots\ Mình cũng chỉ làm những việc bình thường thôi mà\dots\ C-Cảm ơn cậu nhé\dots''

Tôi vội vàng lấy hai tay che má, cúi gầm mặt xuống để giấu đi sự bối rối. Tại sao chỉ một câu khen đơn giản của cậu ấy lại có thể khiến tôi quay về giọng điệu của một đứa trẻ thế này chứ? Thật là thiếu chín chắn quá đi mất!

Kuon bật cười rạng rỡ, rồi lấy từ trong túi ra một món quà nhỏ trao cho tôi. ``Đây là phần của tôi. Hy vọng nó cũng giúp ích cho việc làm đẹp của cậu.''

Tôi nhận lấy món quà, cảm giác hạnh phúc lan tỏa khắp cơ thể. ``C-Cảm ơn cậu nhiều nhé, Kuon-kun\dots\ Mình sẽ trân trọng nó lắm đó\~{}''

Trong không gian tĩnh lặng của tiệm trà cổ, giữa mùi hương trầm mặc và vị đắng ngọt của sô-cô-la, tôi nhận ra rằng: đôi khi việc cố tỏ ra trưởng thành không quan trọng bằng việc được sống thật với cảm xúc của mình trước mặt người mình yêu quý. Và hôm nay, Valentine này thực sự là một kỷ niệm vô cùng ngọt ngào.

\begin{table}[H]
  \centering
  \setlength{\arrayrulewidth}{1pt}
  \begin{tabular}{|>{\centering\arraybackslash}p{0.4\linewidth}|>{\centering\arraybackslash}p{0.4\linewidth}|}
    \hline
    \hyperref[choice]{Quay trở về lựa chọn} & \hyperref[ending]{Đi tới phần cuối} \\ \hline
  \end{tabular}
\end{table}

% XXIII. Onidoro Ayame & Hitsujikado Wata
\newpage
\label{ayawata}

``Hửm?''

Phong bì tôi cầm lên có chút khác biệt so với phần còn lại. Đó thực chất là hai phong bì được đính chặt vào nhau bằng một chiếc ghim kẹp nhỏ. Một phong bì có màu đỏ-trắng rực rỡ với hình dán mặt quỷ oni tinh nghịch, phong bì còn lại có màu vàng-trắng ấm áp với hình dán chú cừu bông xốp.

Tôi mở phong bì đỏ-trắng ra trước, nét chữ bên trong đầy vẻ thách thức:

\txtbx{``Này, Kuon-kun! Nếu có đủ dũng khí thì hãy mò đến phòng Hội học sinh sau giờ học nhé! Tôi đã giăng sẵn thiên la địa võng để đón tiếp cậu rồi đó! Yo-ho-ho!''\\
  \begin{flushright}-Onidoro Ayame-\end{flushright}}

Và ngay phía sau là phong bì vàng-trắng, nét chữ dịu dàng như đang cố xoa dịu lời tuyên chiến ban nãy:

\txtbx{``Xin lỗi cậu vì sự ồn ào của Ayame-chan nhé\dots\ Cậu ấy chỉ đang quá phấn khích thôi. Bọn mình sẽ đợi cậu ở phòng Hội học sinh, có trà và bánh ngọt (và cả sô-cô-la nữa) đang chờ cậu đó. Cậu nhớ ghé qua nhé.''\\
  \begin{flushright}-Hitsujikado Wata-\end{flushright}}

Trưởng và Phó hội học sinh à? Sự kết hợp giữa một người nhây nhớt và một người hiền hòa này luôn mang lại những tình huống khó lường. Tôi cất ``cặp bài trùng'' này vào túi và bước về phía văn phòng Hội học sinh.

\ct

\pov{Onidoro Ayame}

Tôi lén nhìn qua khe cửa phòng Hội học sinh, tay giữ chặt một sợi dây cước mỏng. Trên đỉnh cửa, tôi đã đặt sẵn một xô nhỏ đựng đầy những cánh hoa giấy và ruy băng lấp lánh. Chỉ cần ai đó mở cửa bước vào, và bụp!, một cơn mưa rực rỡ sẽ trút xuống đầu họ.

``Hehe, chắc chắn Kuon-kun sẽ giật mình cho mà xem! Wata-chan, cậu nhìn xem, góc độ này đã chuẩn chưa?'' tôi hào hứng quay sang hỏi cô bạn thân.

\pov{Hitsujikado Wata}

Tôi thở dài, tay vẫn đang cẩn thận bày biện đĩa bánh quy và sô-cô-la lên bàn trà. Nhìn bộ dạng nhí nhố của Ayame-chan, tôi vừa thấy buồn cười vừa thấy lo lắng cho vị khách sắp tới.

``Ayame-chan nè, cậu chơi khăm vừa thôi chứ. Lỡ Kuon-kun giận thì sao? Với lại dọn dẹp đống hoa giấy đó cũng mệt lắm đấy,'' tôi khẽ nhắc nhở, nhưng trong lòng vẫn không kìm được sự tò mò về phản ứng của Kuon. Dù sao thì, những trò đùa của cô bạn luôn khiến không khí Hội học sinh trở nên sôi động hơn.

\pov{Onidoro Ayame}

``Giận sao được mà giận! Đây là đặc ân của Trưởng hội học sinh dành cho cậu ấy đó nha! Suỵt, có tiếng bước chân kìa!'' tôi vội vàng nín thở, kéo căng sợi dây cước.

*CẠCH*

Cánh cửa mở ra. Tôi lập tức buông tay. Xô hoa giấy đổ ập xuống, ruy băng bay tứ tung.

``Yoooo! Chào mừng phàm nhân đến với lãnh địa của quỷ! Yo-ho-ho!'' tôi nhảy ra, hai tay giơ cao đầy vẻ đắc thắng.

\pov{Hitsujikado Wata}

Tôi đứng bật dậy, định chạy ra xin lỗi nhưng rồi đứng hình khi nhìn thấy Kuon đang đứng giữa đống hoa giấy lộn xộn. Cậu ấy đưa tay gạt một mẩu ruy băng màu hồng trên đầu, vẻ mặt đầy cam chịu nhưng cũng không giấu nổi nụ cười khổ.

``\dots Chơi khăm hay đấy, Ayame-san. Không hổ danh cô nàng bày trò có khác.''

Tôi bước lại gần, nhẹ nhàng lấy một vài cánh hoa giấy còn sót lại trên vai cậu ấy: ``Xin lỗi cậu nhiều nhé, Kuon-kun. Ayame-chan cứ nhất quyết muốn 'chào đón' cậu theo cách đặc biệt này. Cậu không sao chứ?''

``À, tôi ổn,'' cậu ấy từ tốn đáp lại. ``Tôi còn từng gặp mấy thứ còn ba chấm hơn thế nhiều.''

Tôi thở một hơi dài, tưởng rằng cậu ấy sẽ giận chúng tôi chứ. Tôi liếc mắt trừng nhẹ Ayame-chan một cái, ra hiệu cho cậu ấy thôi cái trò nhây nhớt đó lại.

\pov{Onidoro Ayame}

Thấy Wata bắt đầu ra dáng ``người mẹ'' hiền hậu, tôi cũng thu bớt vẻ phấn khích lại, nhưng vẫn không quên trêu chọc thêm một câu: ``Hừm, phản xạ của cậu hôm nay có vẻ hơi chậm chạp đó nha, Kuon-kun! Cái phản xạ né tránh ở hội thao của cậu đâu rồi?''

Tôi kéo cậu ấy lại phía bàn trà, nơi Wata đã chuẩn bị sẵn mọi thứ. ``Thôi được rồi, coi như đền bù cho cậu, hãy thưởng thức sô-cô-la hảo hạng nhất của Vương quốc Quỷ này đi!''

Tôi chìa ra một hộp quà màu đỏ-trắng rực rỡ, bên trong là những viên sô-cô-la hình mặt quỷ oni tí hon trông cực kỳ tinh quái.

\pov{Hitsujikado Wata}

Tôi mỉm cười, đẩy đĩa bánh quy do mình tự nướng về phía Kuon. ``Đây là phần của mình nữa. Sô-cô-la này mình đã thêm một chút vị sữa thơm béo để trung hòa lại vị đắng của cacao. Cậu ăn thử xem có hợp khẩu vị không nhé.''

Kuon ngồi xuống, lần lượt nếm thử món quà của cả hai chúng tôi. Nhìn vẻ mặt hài lòng của cậu ấy khi thưởng thức vị đắng ngọt hòa quyện, tôi cảm thấy một niềm vui dịu dàng lan tỏa. Tính cách điềm tĩnh và bao dung của Kuon thực sự rất hợp để kìm chế sự hiếu động của Ayame.

\pov{Onidoro Ayame}

``Thế nào? Tuyệt cú mèo đúng không?'' tôi chống cằm nhìn cậu ấy, chiếc đuôi vô hình sau lưng như đang ngoáy tít vì sung sướng khi thấy cậu ấy khen ngon. ``Phải ăn hết đó nha, nếu không ta sẽ bắt cậu làm chân sai vặt cho Hội học sinh suốt cả học kỳ tới đó!''

Tôi liếc nhìn cô bạn, thấy cậu ấy cũng đang mỉm cười hạnh phúc. Dù trò chơi khăm của tôi có hơi lố một chút, nhưng cuối cùng mọi thứ vẫn diễn ra thật tuyệt vời.

\pov{Hitsujikado Wata}

Kuon sau đó cũng lấy từ trong túi ra hai hộp quà nhỏ đưa cho chúng tôi. ``Đây là quà đáp lễ cho cả hai cậu. Cảm ơn vì buổi tiếp đãi đặc biệt này nhé.''

Chúng tôi nhận lấy món quà từ tay cậu con trai ấy, cảm thấy hơi ấm từ bàn tay cậu ấy khi trao quà. Trong không gian ấm cúng của phòng Hội học sinh, giữa tiếng cười nói nhí nhố của Ayame và sự dịu dàng của Kuon, tôi nhận ra rằng đây chính là ngày Valentine ấm áp nhất mà mình từng có.

``Cảm ơn cậu nhé, Kuon-kun. Lần sau cậu lại ghé chơi nhé, nhưng hứa là sẽ không có xô hoa giấy nào nữa đâu,'' tôi khẽ cười, nhìn hai người họ lại bắt đầu trêu chọc nhau một cách thân thiết.

Dưới ánh hoàng hôn muộn rọi qua cửa sổ văn phòng, ba chiếc bóng đổ dài trên mặt sàn, tạo nên một bức tranh thật hài hòa giữa sự sôi nổi của quỷ và sự bình yên của cừu. Valentine này, quả thực rất nhộn nhịp!

\begin{table}[H]
  \centering
  \setlength{\arrayrulewidth}{1pt}
  \begin{tabular}{|>{\centering\arraybackslash}p{0.4\linewidth}|>{\centering\arraybackslash}p{0.4\linewidth}|}
    \hline
    \hyperref[choice]{Quay trở về lựa chọn} & \hyperref[ending]{Đi tới phần cuối} \\ \hline
  \end{tabular}
\end{table}

% XXIV. Himekawa Mitsuki
\newpage
\label{mitsuki}

Trong xấp bức thư trong tủ, tôi bốc đại một lá. Phong bì tôi chọn có màu tím nhạt pha lẫn sắc hồng dịu dàng, ở giữa nắp phong bì là một hình dán cục kẹo được gói cẩn thận. Một sự lựa chọn rất đỗi ngọt ngào và đậm chất ``công chúa''. Nét chữ bên trong thanh thoát, mềm mại:

\txtbx{``Kuon-kun nè, chiều nay cậu có muốn cùng mình đi dạo một chút không na? Mình sẽ đợi cậu ở đài phun nước trung tâm thương mại trước ga Aogami nhé. Có một vài thứ mình muốn chia sẻ cùng cậu.''\\
  \begin{flushright}-Himekawa Mitsuki-\end{flushright}}

Trung tâm thương mại trước ga Aogami à? Địa điểm này nằm ngay trung tâm thị trấn. Để đến được đó mà không bị ai bắt gặp, tôi buộc phải tính toán kỹ lưỡng.

Thay vì đi ra cổng chính, tôi lẻn ra phía sau dãy hành lang khu thực hành, men theo bức tường gạch cũ rồi leo qua một đoạn rào thấp dẫn ra con đường nhỏ phía đông. Từ đây, tôi rảo bước nhanh qua những con hẻm vắng, giữ khoảng cách an toàn với khu vực ven sông. Cuối cùng, bóng dáng hiện đại của khu trung tâm thương mại cũng hiện ra trước mắt.

\ct

\pov{Himekawa Mitsuki}

Tôi đứng bên cạnh đài phun nước, ánh mắt lơ đãng nhìn những tia nước tung bọt trắng xóa dưới ánh nắng chiều. Hôm nay tôi chọn một bộ váy nhẹ nhàng hơn thường lệ, phối cùng chiếc mũ nồi nhỏ xinh. Dù luôn tự tin với phong thái của một ``công chúa'', nhưng khi nghĩ đến việc lát nữa sẽ gặp Kuon, tôi lại thấy lòng mình xốn xang lạ thường.

Trong tay tôi là một hộp quà màu tím pastel, được thắt nơ lụa vô cùng tỉ mỉ.

``Mitsuki-san!''

Tiếng gọi vang lên giữa đám đông ồn ào. Tôi quay lại và thấy Kuon đang tiến về phía mình, mồ hôi lấm tấm trên trán nhưng nụ cười vẫn rạng rỡ.

``Kuon-kun! Cậu đến đúng giờ thật đấy na,'' tôi mỉm cười, khôi phục lại vẻ tự tin vốn có. ``Hôm nay trung tâm thương mại có vẻ nhộn nhịp nhỉ. Cậu có muốn cùng mình đi dạo một vòng không?''

Chúng tôi cùng nhau bước vào bên trong. Tôi tự tin dẫn dắt cậu ấy đi qua những gian hàng thời trang, những tiệm đồ lưu niệm lấp lánh. Tôi thao thao bất tuyệt kể về những bộ sưu tập mới, về những món đồ dễ thương mà mình vừa khám phá ra. Kuon vẫn kiên nhẫn bước bên cạnh, lắng nghe và thỉnh thoảng lại đưa ra những lời nhận xét rất tinh tế.

Tuy nhiên, khi chúng tôi bước vào một khu vực vắng người hơn, cuộc trò chuyện bỗng nhiên chững lại. Một khoảng lặng kỳ lạ bao trùm lấy cả hai.

Tôi bỗng thấy mặt mình nóng bừng lên. Chưa kịp để cậu ấy nói lên điều gì, tôi lập tức lên tiếng: ``À\dots\ thật ra thì\dots\ mình có thứ này muốn tặng cậu.''

Tôi chìa hộp quà về phía cậu ấy, đôi mắt khẽ lảng tránh. ``Lễ tình nhân mà. Mình đã tự tay chọn loại kẹo sô-cô-la này đó. Nó mang hương vị ngọt ngào của những giấc mơ\dots\ nanora.''

Kuon ngạc nhiên đón lấy món quà. ``Cảm ơn cậu, Mitsuki-san. Một chiếc hộp thật đẹp, đúng là phong cách của cậu.''

Cậu ấy mở hộp, lấy một viên kẹo sô-cô-la nhỏ nếm thử. Ánh mắt cậu ấy sáng lên: ``Vị ngọt thật thanh thoát và dễ chịu. Cảm ơn cậu nhiều nhé, Mitsuki-san.''

Thấy cậu ấy hài lòng, tôi cảm thấy một niềm hạnh phúc khó tả. Chúng tôi lại tiếp tục bước đi, nhưng lần này khoảng cách d dường như đã thu hẹp lại đôi chút.

Kuon sau đó cũng lấy ra một món quà nhỏ trao cho tôi. ``Đây là phần của mình dành cho cậu. Hy vọng nó sẽ mang lại cho cậu những niềm vui ngọt ngào nhất.''

Tôi đón lấy món quà, cảm thấy tâm hồn mình như đang bay bổng. Giữa không gian hiện đại và náo nhiệt của trung tâm thương mại, tôi nhận ra rằng điều tuyệt vời nhất không phải là những món đồ xa hoa, mà chính là khoảnh khắc được cùng người mình thương mến sẻ chia những điều nhỏ bé và ngọt ngào như thế này.

``Valentine này\dots\ thực sự rất tuyệt vời, Kuon-kun ạ,'' tôi khẽ nói, nụ cười dịu dàng nở trên môi dưới ánh đèn lung linh của khu trung tâm.

\begin{table}[H]
  \centering
  \setlength{\arrayrulewidth}{1pt}
  \begin{tabular}{|>{\centering\arraybackslash}p{0.4\linewidth}|>{\centering\arraybackslash}p{0.4\linewidth}|}
    \hline
    \hyperref[choice]{Quay trở về lựa chọn} & \hyperref[ending]{Đi tới phần cuối} \\ \hline
  \end{tabular}
\end{table}

% XXV. Omori Honoka
\newpage
\label{honoka}

Tôi lượm bừa một bức thư trong xấp phong bì đó. Và ngay khi tôi cầm nó lên, tôi lập tức hơi nhăn mặt.

``\dots Ờm\dots\ Cái quái gì thế này?''

Bức thư tôi cầm lên có lẽ là bức thư rực rỡ và kỳ dị nhất trong tủ. Phong bì là một sự pha trộn hỗn loạn giữa vàng, đen, hồng và xanh lam. Nhưng điểm khiến tôi đứng hình chính là thứ được dán trên nắp phong bì: một chỏm tóc thật màu hồng-đen được thắt lại trông giống như cánh quạt trực thăng. Không phải hình dán, là \emph{tóc thật}.

Chỉ có thể là Honoka; cô nàng có những ý tưởng nghe qua rất điên rồ, và cái này có thể là một trong những thứ như vậy. Tôi mở thư, và đúng như dự đoán, những dòng chữ nhảy múa như đang reo hò:

\txtbx{``Kuon-kun! Mau đến lớp học trống ở dãy nhà cũ ngay và luôn nhé! Tớ có một 'siêu phẩm' muốn cho cậu xem nè! Đừng bắt 'bậc thầy giải trí' này phải đợi lâu đó nha!''\\
  \begin{flushright}-Omori Honoka-\end{flushright}}

Lớp học trống à? Nơi này gợi lại cho tôi khá nhiều kỷ niệm, nhất là vào cái đêm gan dạ thảm kịch mà Honoka là người cầm đầu. Tôi cất bức thư (với chỏm tóc vẫn còn dính trên đó) vào túi và bước về phía dãy nhà cũ im lìm.

\ct

\pov{Omori Honoka}

Tôi ngồi vắt vẻo trên bàn giáo viên, đôi chân đung đưa không ngừng. Nhìn quanh căn phòng trống trải đầy bụi bặm này, tôi bỗng thấy một cảm giác quen thuộc đến lạ lùng. Cứ như thể tôi đã từng đứng đây, trong một hoàn cảnh tương tự, để chờ đợi một điều gì đó trọng đại lắm vậy.

``A, Kuon-kun! Cậu đến chậm quá đó nha!'' tôi nhảy xuống đất ngay khi thấy bóng dáng cậu ấy xuất hiện ở cửa.

Kuon nhìn tôi, rồi nhìn vào căn phòng: ``Cậu lại chọn cái nơi hẻo lánh này à, Honoka-san? Lại định bày trò gì nữa đây?''

``Trò gì là trò gì! Hôm nay là ngày đặc biệt mà!'' tôi cười lớn, lôi từ trong chiếc túi đeo chéo ra một hộp quà rực rỡ không kém gì bức thư ban nãy. ``Tèn tén ten! Siêu phẩm sô-cô-la 'Ngọt Ngào Vô Tận' do chính tay Honoka này chế tác!''

Tôi ấn hộp quà vào tay cậu ấy, miệng không ngừng liến thoắng kể về việc tôi đã phải đun chảy bao nhiêu cân đường và mật ong để tạo ra hương vị này. Kuon mở hộp, nhìn những viên sô-cô-la bóng loáng, nét mặt có chút do dự: ``Nghe tên thôi đã thấy máy chạy thận trước mắt rồi\dots''

Cậu ấy nếm thử một viên. Ngay lập tức, khuôn mặt cậu ấy biến sắc, mắt trợn ngược: ``Ngọt\dots\ ngọt quá! Honoka-san, cậu bỏ bao nhiêu đường vào đây vậy?''

``Hehe, bí mật quân sự đó nha! Nhưng mà ngọt mới là biểu tượng của tình yêu chân thành chứ, đúng không?'' tôi nháy mắt tinh nghịch. Dù miệng nói đùa, nhưng nhìn Kuon cố gắng nuốt viên sô-cô-la siêu ngọt đó, tôi cảm thấy trong lòng dâng lên một sự ấm áp khó tả. Đằng sau những lời nói ồn ào này, tôi thực sự mong cậu ấy có thể cảm nhận được chút dư vị ngọt ngào giữa cuộc đời đầy biến động của chúng tôi.

Bỗng dưng, tôi nhìn vào túi áo của Kuon, nơi chỏm tóc hồng-đen đang ló ra. ``À đúng rồi! Cho tớ xin lại chỏm tóc đó đi, lúc nãy vội quá tớ cắt hơi quá tay, giờ nhìn gương thấy thiếu thiếu sao á!''

Tôi vươn tay định lấy, nhưng Kuon đã nhanh chóng né sang một bên, giơ bức thư lên cao: ``Đâu có được, tặng rồi là của tôi chứ. Với lại để đây trông cũng giống\dots\ cánh máy bay trực thăng, hay phết.''

``Này! Đừng có trêu tớ chứ! \textbf{Trả đây!!}'' tôi phồng má, nhảy lên định giật lại, cả hai rượt đuổi nhau quanh những bộ bàn ghế cũ, tiếng cười vang vọng khắp căn phòng vắng.

Sau một hồi náo loạn, tôi đứng lại, thở hổn hển nhưng mắt lại sáng rực một ý tưởng điên rồ. Tôi tiến lại gần Kuon, ghé sát tai cậu ấy thì thầm:

``Nè Kuon-kun, sô-cô-la ngọt thế này mà chỉ ăn ở đây thì phí quá. Cậu có muốn cùng mình lẻn lên sân thượng trung tâm thương mại tối nay để\dots\ bắn pháo hoa giấy và hét thật to cho cả thị trấn nghe không? Mình vừa kiếm được một xấp pháo lậu xịn lắm!''

Tôi nhìn cậu ấy với ánh mắt đầy vẻ mời gọi đầy liều lĩnh. Với Honoka này, Valentine không chỉ là sự ngọt ngào, mà còn phải là một cuộc phiêu lưu không hồi kết! Cậu thấy sao, Kuon-kun?

\begin{table}[H]
  \centering
  \setlength{\arrayrulewidth}{1pt}
  \begin{tabular}{|>{\centering\arraybackslash}p{0.4\linewidth}|>{\centering\arraybackslash}p{0.4\linewidth}|}
    \hline
    \hyperref[choice]{Quay trở về lựa chọn} & \hyperref[ending]{Đi tới phần cuối} \\ \hline
  \end{tabular}
\end{table}

% XXVI. Akaba Koishin
\newpage
\label{koishin}

Tôi lấy bừa một bức trong xấp thư đó. Bức thư tôi cầm lên có màu vàng chanh vô cùng nổi bật, nhưng điều khiến tôi chú ý hơn cả là những sợi ruy băng đỏ được buộc chi chít xung quanh phong bì. Cảm giác như người gửi đã phải rất vất vả để có thể gói ghém tâm tư của mình vào bên trong.

Nét chữ bên trong có chút run rẩy nhưng đầy quyết đoán:

\txtbx{``Tên ngốc Kuon kia! Sau giờ học, cấm cậu không được đi đâu mà phải đến phòng CLB Văn học ngay rõ chưa! Tôi\dots\ tôi có chuyện muốn tính sổ với cậu! Đừng có mà trốn đó!''\\
  \begin{flushright}-Akaba Koishin-\end{flushright}}

Tính sổ ư? Nghe giọng điệu này, tôi có thể đoán ngay được khuôn mặt đang đỏ bừng của Koishin lúc viết những dòng này. Tôi khẽ mỉm cười, cất bức thư vào túi và hướng về phía phòng CLB Văn học nằm ở cuối hành lang dãy khu câu lạc.

\ct

\pov{Akaba Koishin}

Tôi đứng ngồi không yên trong phòng CLB, đôi tay cứ vô thức xoắn lấy tà áo. Trên bàn là một hộp sô-cô-la được gói ghém cực kỳ cẩn thận, dù tôi đã phải làm đi làm lại hàng chục lần mới ưng ý.

``Tên ngốc đó\dots\ sao mãi mà chưa đến chứ? Hay là cậu ta vứt bức thư của mình đi rồi?'' tôi lẩm bẩm, lòng nóng như lửa đốt.

Tiếng cửa phòng khẽ mở. Kuon bước vào, vẫn là cái vẻ mặt điềm tĩnh đáng ghét đó.

``Ồ, Koishin-san, cậu đợi tôi lâu chưa?''

``L-Lâu cái gì mà lâu! Cậu có biết tôi đã tốn bao nhiêu thời gian quý báu để ngồi đây không hả?'' tôi đứng bật dậy, mặt đỏ bừng lên như quả cà chua chín. ``Mà ai cho cậu gọi tên tôi thân mật thế hả tên ngốc kia!''

Kuon khẽ bật cười, tiến lại gần bàn. ``Xin lỗi, xin lỗi. Tại bức thư của cậu buộc ruy băng kỹ quá, tôi phải mất hồi lâu mới mở ra được.''

``Hừ, kệ cậu chứ!'' tôi hét lên, rồi lấy hết can đảm, cầm hộp sô-cô-la dí mạnh vào ngực cậu ấy. ``Này! Cầm lấy đi! Đây là sô-cô-la tôi làm thừa nên mới cho cậu thôi nhé! Đừng có mà hiểu lầm là tôi đặc biệt làm cho cậu đấy!''

Kuon dịu dàng đón lấy món quà, ánh mắt cậu ấy nhìn tôi ấm áp đến mức khiến mọi sự phòng thủ trong tôi dường như tan chảy. ``Cảm ơn cậu, Koishin-san. Tôi sẽ trân trọng nó.''

Cậu ấy mở hộp, và bên cạnh những viên sô-cô-la, một mảnh giấy nhỏ có chép một bài thơ hiện ra. ``Ồ, trong này có cái gì nữa này.''

``Ơ, cái đó\dots\ đừng có đọc!'' tôi hốt hoảng định giật lại nhưng đã muộn.

Kuon khẽ tằng hắng, rồi bắt đầu đọc to những dòng thơ mà tôi đã thức trắng đêm để viết:

\begin{center}
  \textit{``Ánh nắng chiều vương trên trang giấy trắng,\\
    Lòng ta bỗng thấy xao xuyến lạ thường\dots''}
\end{center}

``\textbf{Dừng lại! Đã bảo là đừng có đọc mà!}'' tôi ngượng chín mặt, cảm giác như có luồng khói đang bốc ra từ đỉnh đầu. Tôi lấy hai tay che mặt, chỉ muốn tìm một cái lỗ nào đó để chui xuống ngay lập tức.

Kuon dừng lại, nhìn tôi với ánh mắt đầy ý nhị: ``Sao vậy? Cậu không muốn tôi đọc bài thơ mà cậu đã tâm huyết viết sao?''

``\textbf{K-Không phải thế!}'' tôi hét lên, nhưng rồi lại lí nhí: ``Chỉ là\dots\ nghe cậu đọc trực tiếp thế này\dots\ ngượng lắm\dots''

Kuon khẽ mỉm cười, gấp mảnh giấy lại cẩn thận. ``Bài thơ hay lắm, Koishin-san. Hơi ngắn, nhưng thế cũng đủ để khiến tôi vui rồi.''

Nghe câu nói đó, tôi cảm thấy cả người mình như đang bay bổng. Sự bướng bỉnh ban nãy hoàn toàn biến mất, chỉ còn lại sự ngọt ngào lan tỏa trong lồng ngực.

Kuon sau đó cũng lấy ra một món quà nhỏ trao cho tôi. ``Đây là phần của tôi dành cho cậu. Hy vọng nó cũng sẽ 'chạm' đến trái tim cậu như bài thơ vừa rồi.''

Tôi nhận lấy món quà, cúi gầm mặt xuống để che giấu nụ cười đang nở rộ trên môi. ``H-Hừ, chỉ là quà đáp lễ thôi đúng không? Được rồi, tôi sẽ nhận\dots\ vì cậu đã có lòng.''

Valentine này, dù tôi vẫn luôn miệng mắng cậu ấy là tên ngốc, nhưng trong thâm tâm, tôi biết mình chính là kẻ ngốc nhất khi đã hoàn toàn lún sâu vào cái bẫy ngọt ngào mang tên Hikawa Kuon mất rồi.

``\textit{Này\dots\ lần sau\dots\ đừng có đọc thơ to như thế nữa nhé\dots}'' tôi thì thầm, lòng thầm mong khoảnh khắc này sẽ kéo dài mãi mãi.

\begin{table}[H]
  \centering
  \setlength{\arrayrulewidth}{1pt}
  \begin{tabular}{|>{\centering\arraybackslash}p{0.4\linewidth}|>{\centering\arraybackslash}p{0.4\linewidth}|}
    \hline
    \hyperref[choice]{Quay trở về lựa chọn} & \hyperref[ending]{Đi tới phần cuối} \\ \hline
  \end{tabular}
\end{table}

% XXVII. Furukawa Nanase
\newpage
\label{nanase}

Tôi tìm thấy một phong bì có màu xanh nước biển với họa tiết ca-rô vô cùng lịch thiệp và chỉn chu. Ngay trên nắp phong bì là một hình dán ngôi sao chổi đang lướt đi, rạng rỡ và đầy kiêu hãnh. Nét chữ bên trong thẳng tắp, gãy gọn như chính tính cách của người gửi:

\txtbx{``Kuon-san. Tan học hãy đến cầu thang ở khu nhà C nhé. Tôi có chuyện cần trao đổi trực tiếp với cậu. Hy vọng cậu sẽ không để tôi phải chờ lâu.''\\
  \begin{flushright}-Furukawa Nanase-\end{flushright}}

Khu nhà C à? Nơi đó thường khá vắng vẻ vào giờ này. Tôi cất bức thư vào túi và rảo bước về phía khu nhà C. Tính cách thẳng thắn của Nanase-san luôn khiến tôi cảm thấy nể phục, và dường như bức thư này cũng không ngoại lệ.

\ct

\pov{Furukawa Nanase}

Tôi đứng tựa lưng vào lan can cầu thang, đôi mắt dõi nhìn ra khoảng sân trường đang dần thưa thớt học sinh. Trên tay tôi là một hộp quà màu xanh thẫm, đơn giản nhưng vô cùng tinh tế. Tôi không thích những thứ rườm rà, và trong chuyện tình cảm cũng vậy.

Nghe tiếng bước chân vang lên đều đặn, tôi không cần quay lại cũng biết đó là ai.

``Cậu đến rồi à, Kuon-san. Rất đúng giờ,'' tôi xoay người lại, nhìn thẳng vào mắt cậu ấy với nụ cười điềm tĩnh.

`` Nanase-san. Có vẻ cậu đã đợi ở đây một lúc rồi nhỉ?''

``Cũng không lâu lắm. Tôi thích không gian yên tĩnh ở đây để suy nghĩ,'' tôi bước lại gần cậu ấy, đưa hộp quà ra mà không hề có chút do dự hay e thẹn thái quá nào. ``Cái này dành cho cậu. Lễ tình nhân, tôi nghĩ một món quà nhỏ là điều cần thiết để bày tỏ sự trân trọng giữa những người bạn.''

Kuon đón lấy món quà, khẽ gật đầu cảm ơn. Cậu ấy mở hộp, lấy một viên sô-cô-la nguyên chất nếm thử.

``Vị đắng thật đậm đà\dots\ Rất đặc biệt, Nanase-san.''

``Đó là sô-cô-la nguyên chất 85\% cacao. Tôi nghĩ vị đắng của nó tượng trưng cho sự trưởng thành và thực tế của cuộc sống này. Không phải cái gì cũng cần phải ngọt lịm thì mới là tốt, đúng không?'' tôi thẳng thắn chia sẻ quan điểm của mình.

Kuon mỉm cười: ``Quả đúng là phong cách của cậu. Rất chân thực và sâu sắc.''

Thấy cậu ấy hiểu được ý nghĩa của món quà, tôi cảm thấy một sự hài lòng nhẹ nhàng. Tiện đây, tôi cũng muốn đưa ra một lời đề nghị: ``Kuon-san, cuối tuần này nhà Furukawa có tổ chức một buổi dã ngoại nhỏ ở ngoại ô. Trước đây tôi từng rủ Suzuki-san tham gia, và lần trước Kaigou-san cũng đã đi cùng chúng ta ở trung tâm thương mại rồi. Lần này, tôi muốn mời riêng cậu. Cậu thấy sao?''

Kuon có chút bất ngờ trước lời mời trực diện này. ``Lời mời từ tiểu thư nhà Furukawa à? Tôi thực sự rất vinh hạnh.''

``Vậy thì quyết định thế nhé. Tôi sẽ gửi thông tin chi tiết sau,'' tôi khẽ mỉm cười, cảm nhận được một sợi dây liên kết vô hình nhưng chắc chắn đang dần hình thành giữa chúng tôi.

Kuon sau đó cũng lấy ra một món quà nhỏ trao cho tôi. ``Đây là phần của tôi. Hy vọng nó cũng mang lại cho cậu những suy nghĩ thú vị cho ngày cuối tuần.''

Tôi nhận lấy món quà, lòng thầm đánh giá cao sự tinh tế của cậu ấy. Valentine này, đối với tôi, không cần sự ồn ào hay những lời nói hoa mỹ. Chỉ cần một cuộc trò chuyện thẳng thắn và một lời hứa cho tương lai như thế này là đã quá đủ rồi.

``Hẹn gặp cậu vào cuối tuần nhé, Kuon-san,'' tôi nói, rồi thong thả bước xuống cầu thang, để lại một dư vị đắng ngọt đầy ấn tượng trong lòng cậu bạn.

\begin{table}[H]
  \centering
  \setlength{\arrayrulewidth}{1pt}
  \begin{tabular}{|>{\centering\arraybackslash}p{0.4\linewidth}|>{\centering\arraybackslash}p{0.4\linewidth}|}
    \hline
    \hyperref[choice]{Quay trở về lựa chọn} & \hyperref[ending]{Đi tới phần cuối} \\ \hline
  \end{tabular}
\end{table}

% Kaizuka Yuki
\newpage
\label{yuki}

\pov{Hikawa Kuon}

Tôi vớ bừa một bức thư trong tủ ra.

Bức thư này có màu đen-xám với những vạch dọc đỏ đầy cá tính. Ngay chính giữa phong bì là một hình dán rô-bốt nhỏ nhắn, khiến tôi bất giác nhớ về Roboco ở thế giới thực - nơi cô ấy thực sự là một rô-bốt với công nghệ đỉnh cao\dots\ đầu hàng. Còn ở đây, Yuki-san chỉ là một cô gái bình thường với phong cách thời trang găng tay và tất cao gối đặc trưng, như để phóng tác lại phiên bản gốc ở thế giới của tôi.

Nét chữ bên trong bức thư tràn đầy sự phấn khích:

\txtbx{``Kuon-kun! Tan học đừng có về nhà sớm nhé! Mau đến quán karaoke 'Super Nova' ở gần ga đi, mình đã đặt phòng sẵn rồi đó! Bọn mình sẽ có một buổi chiều bùng nổ! Đừng có mà từ chối đó nha!''\\
  \begin{flushright}-Kaizuka Yuki-\end{flushright}}

Quán karaoke à? Địa điểm này nằm hoàn toàn bên ngoài trường. Để đảm bảo không bị bất cứ ai bắt gặp, tôi quyết định lẻn ra bằng cửa sau của dãy nhà thực hành. Thay vì đi theo hướng đài phun nước nhộn nhịp, tôi chọn con đường mòn phía ngược lại, đi vòng qua khu dân cư yên tĩnh để tiến về phía ga.

\ct

\pov{Kaizuka Yuki}

Tôi ngồi trong phòng karaoke, tay cầm micro nhún nhảy theo điệu nhạc sôi động đang phát trên màn hình. Trên bàn là một hộp sô-cô-la bọc giấy bạc lấp lánh như những linh kiện điện tử cao cấp. Tôi đã chuẩn bị sẵn danh sách những bài hát song ca hay nhất để cùng Kuon bùng nổ chiều nay.

Tiếng cửa phòng mở ra, Kuon bước vào với vẻ mặt có chút ngơ ngác trước âm thanh cực lớn trong phòng.

``Yuki-san! Cậu gọi tôi đến đây thật đấy à?''

Tôi cười rạng rỡ, đưa chiếc micro còn lại cho cậu ấy: ``Đương nhiên rồi! Lễ tình nhân mà chỉ tặng quà thôi thì chán lắm, phải có âm nhạc nữa mới đúng điệu chứ! Nào, Kuon-kun, hát cùng mình bài này đi!''

Kuon nhìn chiếc micro với vẻ mặt đầy ái ngại: ``Thôi cho tôi xin kiếu đi Yuki-san\dots\ Tôi hát dở lắm, cậu cũng biết rồi mà. Nghe xong chắc cậu không muốn tặng quà cho tôi luôn quá.''

``Không chấp nhận lời từ chối đâu\~{}'' tôi nhấn nút chọn bài và đẩy cậu ấy về phía màn hình. ``Ở đây không có ai ngoài hai đứa mình đâu, cứ hát hết mình đi! Đây là mệnh lệnh của 'High-spec Yuki' đó nha\~{}''

Thế là, mặc cho Kuon ra sức van nài, tôi vẫn ép cậu ấy phải cầm micro. Tiếng hát của cậu ấy quả thực có chút\dots\ lệch tông, nhưng nhìn vẻ mặt nghiêm túc cố gắng theo kịp nhịp điệu của cậu ấy, tôi không thể nhịn được cười. Sự vụng về đáng yêu đó khiến tôi cảm thấy cậu ấy gần gũi hơn bao giờ hết.

Sau khi đã hát hò thỏa thích đến mức khản cả cổ, tôi mới cầm hộp sô-cô-la lên và đưa cho cậu ấy.

``Nè, phần thưởng cho nỗ lực ca hát của cậu đó!'' tôi nháy mắt tinh nghịch. ``Đây là sô-cô-la 'High-spec' đặc biệt, ăn vào là đảm bảo giọng hát của cậu sẽ được nâng cấp lên level mới luôn!''

Kuon nhận lấy món quà, cười khổ: ``Cảm ơn cậu, Yuki-san. Tôi chỉ hy vọng nó không khiến tôi muốn đi hát karaoke thêm lần nào nữa trong tháng này.''

Cậu ấy nếm thử một viên và ngạc nhiên trước hương vị độc đáo của nó. ``Ngon thật đấy, vị cacao rất thanh và có chút gì đó rất hiện đại.''

Nghe lời khen ấy, trái tim tôi như được nạp đầy năng lượng. Kuon sau đó cũng tặng tôi một món quà đáp lễ, khiến buổi chiều tại quán karaoke càng thêm phần ngọt ngào.

Dưới ánh đèn màu nhấp nháy của phòng karaoke, tôi nhận ra rằng: đôi khi niềm vui không đến từ những thứ hoàn hảo, mà đến từ chính những phút giây vui vẻ, thoải mái và có phần 'lệch tông' như thế này bên cạnh người mình yêu quý. Valentine này, quả thực là một màn trình diễn tuyệt vời nhất!

\begin{table}[H]
  \centering
  \setlength{\arrayrulewidth}{1pt}
  \begin{tabular}{|>{\centering\arraybackslash}p{0.4\linewidth}|>{\centering\arraybackslash}p{0.4\linewidth}|}
    \hline
    \hyperref[choice]{Quay trở về lựa chọn} & \hyperref[ending]{Đi tới phần cuối} \\ \hline
  \end{tabular}
\end{table}

% XXIX. Canon route
\newpage
\label{canon}

Tôi đứng lặng người trước tủ thư, đôi tay lơ lửng giữa không trung như bị đóng băng bởi một sức mạnh vô hình nào đó. Ánh mắt tôi lướt qua xấp phong bì đa sắc đang nằm ngay ngắn bên trong - đỏ rực rỡ, xanh thanh khiết, vàng ấm áp, và cả những màu sắc kỳ dị đầy cá tính. Mỗi bức thư là một tâm hồn, một lời mời gọi, và trên hết là một mảnh tình cảm chân thành mà các cô gái lớp 1-C đã dành cho tôi.

Tôi biết rõ đằng sau lớp giấy mỏng manh kia chứa đựng những gì. Đó là hy vọng. Đó là niềm vui lần đầu được biết đến hương vị của sự rung động.

Nhưng cũng chính vì thế, trái tim tôi thắt lại.

``\dots''

Tôi khẽ nhắm mắt, và những lời nói của Suzuki và Kaigou lúc sáng lại vang lên trong tâm trí tôi như một bản nhạc buồn đầy ám ảnh. Chúng tôi không thuộc về nơi này. Chúng tôi là những lữ khách lạc bước từ một vũ trụ xa xôi, một thế giới mà sự tồn tại của chúng tôi vốn dĩ đã là một nghịch lý. Một khi những linh hồn ai oán trong khu rừng Aogami được siêu thoát, khi nhiệm vụ cuối cùng của bộ ba Hikawa hoàn thành\dots\ chúng tôi sẽ rời khỏi nơi này, và cũng không chắc liệu chúng tôi có thể quay trở lại nơi này được không.

Tôi cũng có quyền được tương tư chứ. Nếu không, có lẽ tôi sẽ sống mãi trong cái kén của một tên ``tồ'' trong tình yêu, một kẻ không bao giờ dám đương đầu về tình cảm yêu đương. Nhưng\dots

Nếu tôi đưa tay ra và nhận lấy một bức thư, nếu tôi đáp lại tình cảm ấy bằng một nụ cười hay một lời hứa hẹn, thì tôi đang gieo rắc điều gì vào tương lai của họ? Một sự chờ đợi vô vọng? Một nỗi đau xé lòng khi người mình thương yêu đột ngột tan biến vào hư không?

Càng trân trọng họ, tôi lại càng cảm thấy mình không có quyền được hạnh phúc bên họ.

Bàn tay tôi run rẩy, rồi từ từ nắm chặt lại thành nắm đấm. Tôi hít một hơi thật sâu, cảm nhận vị lạnh của không khí buổi chiều len lỏi vào tận phổi, cố gắng dập tắt ngọn lửa nhỏ nhoi của sự ích kỷ đang nhen nhóm trong lòng.

*CẠCH*

Tôi dứt khoát đóng cánh cửa tủ thư lại. Tiếng vang khô khốc ấy như một nhát dao cắt đứt sợi dây liên kết mà tôi hằng mong ước nhưng không dám chạm vào.

Tôi đứng tựa lưng vào tủ, nhìn lên trần hành lang mờ ảo qua làn hơi thở. Một lời thì thầm khàn đục thoát ra từ đôi môi khô khốc, mỏng manh đến mức dường như chỉ mình tôi và những linh hồn trong bức thư kia nghe thấy:

``Tôi xin lỗi\dots\ Thực sự xin lỗi mọi người.''

Tôi xin lỗi vì đã không thể là một phần trong cuộc đời của các bạn lâu hơn. Tôi xin lỗi vì đã chọn cách trốn chạy này để bảo vệ các bạn khỏi một nỗi đau còn lớn hơn gấp bội.

Tôi quay lưng bước đi, mỗi bước chân trên hành lang vắng lặng đều nặng nề như đeo chì. Tôi phải rời khỏi đây, trốn thoát khỏi cái ``thiên la địa võng'' của tình yêu thương này trước khi ý chí của mình hoàn toàn sụp đổ.

\ct

Mặt đồng hồ trên cổ tay tôi bỗng lóe lên một ánh sáng mờ. Một giọng nói vang lên trong đầu tôi, mang theo chất giọng Anh pha âm sắc Pháp đặc sệt cùng một thái độ mỉa mai quen thuộc.

``\eng{Oh la la. Lựa chọn an toàn ghê. Cậu thực sự định bỏ mặc ngần ấy tấm chân tình chỉ để đóng vai một kẻ tử vì đạo sao, Kuon-user?}''

``\eng{Ông im đi, Game,}'' tôi lầm bầm đáp lại trong suy nghĩ. ``\eng{Đây là cách duy nhất.}''

``\eng{Très bien, très bien. Nếu cậu đã chọn làm một kẻ chạy trốn, thì ít nhất cũng phải chạy cho trót. Cơ mà ấy, cậu nghĩ mình có thể dễ dàng rời khỏi đây sao? Nhìn qua cửa sổ đi.}''

Tôi thận trọng tiến đến gần cửa sổ hành lang, hé mắt nhìn xuống. Dưới sân trường, Saki đang cặm cụi tưới những luống hoa, dù tôi biết rõ giờ này chẳng phải lúc làm việc đó. Kaoru và Suzune đang ngồi trên ghế đá, đảo mắt quan sát xung quanh như những radar dò tìm. Ở phía xa, ngay tại cổng chính, bóng dáng quen thuộc của Shino đứng sừng sững như một người gác cổng định mệnh.

``\eng{Thấy chưa?}'' Game cười khẩy. ``\eng{Họ ở khắp mọi nơi luôn, tứ phía đều là các cô gái. Tình yêu của họ quả là một thứ vũ khí đáng sợ. Nên nếu muốn rời khỏi đây, thì phải có chiến lược thật bài bản, Kuon-user. Trừ khi cậu muốn bị họ bắt quả tang và có thể xé xác cậu vì tội dám cho họ leo cây.}''

Tôi nuốt khan. Những kỹ năng ẩn nấp từng giúp tôi sinh tồn ở The Void và thế giới Nijisanji giờ đây lại được dùng để trốn chạy khỏi\dots\ những cô gái tuổi teen. Thật mỉa mai thay.

Tôi bắt đầu di chuyển, hạ thấp trọng tâm, men theo những góc khuất của hành lang. Vừa bước qua ngã rẽ, tôi đã phải khựng lại, dán chặt người vào tường. Inari đang cuống cuồng chạy thục mạng qua dãy hành lang, miệng lẩm bẩm điều gì đó về việc đến muộn. Phía bên kia, Aku đang thong thả đi dạo, mắt không ngừng dáo dác tìm kiếm một bóng hình quen thuộc.

Tôi lùi lại, rẽ sang hướng cầu thang khu nhà C, nhưng lại thoáng thấy mái tóc xanh nước biển của Nanase đang đứng tựa lưng vào lan can, phong thái tĩnh lặng nhưng vô cùng sắc sảo. Hết cách, tôi đành vòng qua khu nhà cũ. Tại đây, tiếng ồn ào văng vẳng từ phòng học trống cho thấy Honoka đang ở bên trong, chuẩn bị cho một trò điên rồ nào đó.

Đi ngang qua dãy phòng câu lạc bộ, tôi như lọt vào một mê cung cảm xúc: tiếng đàn piano da diết của Azuki vang lên từ phòng Âm nhạc, bóng Koishin thấp thoáng qua cửa kính phòng CLB Văn học với khuôn mặt bồn chồn, Mari đang cặm cụi vẽ trong phòng Mỹ thuật, và ánh đèn flash liên tục chớp nháy từ phòng CLB Nhiếp ảnh của Uzuki.

Ngay cả ở khu hành lang chính, tôi cũng phải nín thở lướt qua phòng Hội học sinh, nơi Ayame và Wata đang đợi với những chiếc bẫy đã giăng sẵn, và tránh xa Thư viện nơi Shion đang ngự trị. Tiếng loa rè rè từ phòng phát thanh vọng ra giọng nói hồi hộp của Kanade, xen lẫn tiếng ồn ào từ căng-tin nơi Hotaru đang túc trực bên những khay sô-cô-la.

Mỗi một bóng dáng lướt qua, mỗi một âm thanh lọt vào tai đều như những mũi kim đâm vào trái tim tôi. Họ đều đang chờ đợi. Những tình cảm thuần khiết ấy, những hộp quà chứa đựng bao tâm tư ấy\dots\ tất cả đều đang hướng về tôi.

``\eng{Tiếp tục tiến về phía trước đi. Đừng để sự ân hận làm chậm chân cậu lại,}'' Game nhắc nhở khi tôi khẽ rùng mình, cố gắng gạt đi sự hối hận đang cắn rứt trong lồng ngực.

Tôi lẻn ra khỏi dãy nhà chính bằng cửa sau. Ẩn mình dưới tán cây lớn sau trường, tôi bắt gặp Akane đang lo lắng đứng đợi với khuôn mặt đỏ bừng. Chạy thục mạng qua khu thể chất, bóng dáng Shiina đang miệt mài tập luyện (có lẽ cô ấy muốn vận động trong khi chờ người mà cô thầm mến) lướt qua mắt tôi. Thoát ra được đến bãi đất trống bên ngoài, tôi lại phải nấp vội sau một bức tường cũ khi thấy Saya đang lững thững đá những viên sỏi nhỏ.

Cuối cùng, sau những nỗ lực luồn lách đến nghẹt thở, tôi cũng chạm chân đến bìa rừng Aogami.

Tôi dừng lại, thở hắt ra, rồi ngoái đầu nhìn về phía ngôi trường lần cuối. Trên tầng thượng lộng gió, tôi có thể nhìn thấy hai chấm nhỏ thấp thoáng - Hanabi và Yuka - đang phóng tầm mắt nhìn xuống toàn bộ khuôn viên, có lẽ cũng đang mỏi mòn tìm kiếm một người sẽ không bao giờ xuất hiện.

Nắm đấm tôi siết chặt lại. Chạy trốn khỏi kẻ thù thì dễ, nhưng chạy trốn khỏi tình cảm chân thành của những người yêu thương mình lại là hình phạt tàn khốc nhất. Mỗi ánh mắt chờ đợi vô vọng của họ như đang trút thêm một tảng đá ngàn cân lên đôi vai tôi.

Nhưng tôi không có quyền quay lại.

Tôi kéo thấp cổ áo, dứt khoát xoay người bước vào con đường rừng ẩm ướt, bỏ lại sau lưng ngần ấy những trái tim đang thổn thức vì tôi.

\ct

Tôi rảo bước trên con đường mòn xuyên qua khu rừng Aogami, hơi ẩm của đất và lá mục bốc lên sau cơn mưa sớm khiến không khí càng thêm nặng nề. Tiếng bước chân của tôi xào xạc trên thảm lá khô, dường như cũng đang cố gắng trốn chạy khỏi sự truy đuổi của lương tâm.

``\eng{Cẩn thận đó, Kuon-user. Đừng có mà lơ là cảnh giác. Nhìn lên bậc thang đá dẫn lên đền thờ kìa,}'' giọng nói của Game lại vang lên, lần này có chút gì đó nghiêm túc hơn.

Tôi khựng lại, nấp sau một gốc cây đại thụ. Ngước mắt nhìn lên những bậc thang đá rêu phong dẫn lối lên đền thờ Aogami, tôi thấy Sakura đang ngồi bó gối ở đó. Ánh mắt cô nàng đượm buồn, cứ nhìn chằm chằm xuống con đường dẫn từ trường vào rừng như đang chờ đợi một phép màu.

Tim tôi thắt lại. Sakura - người duy nhất (cùng với Shino) biết rõ về thân phận của chúng tôi, vậy mà cô ấy vẫn chọn cách đứng đây chờ đợi.

``\eng{Cậu không thể dừng lại đâu. Nếu cậu lộ diện lúc này, mọi nỗ lực trốn chạy của cậu sẽ thành công cốc đó.}'' Game nhắc nhở.

Tôi cắn chặt môi, dứt khoát quay đi, tiếp tục men theo bìa rừng để tiến ra khu vực ven sông. Từ xa, tôi đã thấy Touka đang đứng bất động bên lan can, mái tóc xanh lam khẽ bay trong gió chiều. Cô ấy đang nhìn về phía dòng nước chảy xiết, nơi mà có lẽ tôi và cô ấy đã từng có những cuộc trò chuyện về thời tiết và vận mệnh.

``\eng{Touka-san cũng đang túc trực ở đó. Đừng có đi đường chính nếu không muốn bị 'bắt bài',}'' Game tiếp tục chỉ dẫn.

Tôi vòng qua những con hẻm nhỏ, lách mình qua khu phố truyền thống. Qua khung cửa sổ của tiệm trà cũ, tôi thấy Miku đang ngồi nhâm nhi tách trà với vẻ mặt điềm tĩnh đến lạ lùng, nhưng đôi bàn tay đan chặt vào nhau đã tố cáo sự bồn chồn trong lòng cô ấy.

Ra đến gần trung tâm thương mại, tôi lén nhìn về phía đài phun nước. Mitsuki vẫn đứng đó, kiêu hãnh và rực rỡ như một nàng công chúa giữa đám đông tan tầm nhộn nhịp. Tôi vội vàng kéo sụp mũ áo, hòa mình vào dòng người đang hối hả trở về nhà, cố gắng thu nhỏ sự hiện diện của mình hết mức có thể.

Khi đi ngang qua quán karaoke 'Super Nova', một giai điệu sôi động vọng ra khiến tôi bất giác mỉm cười chua chát. Tôi có thể mường tượng được cảnh Yuki đang hào hứng chọn bài, tay cầm micro và không ngừng dáo dác nhìn ra cửa phòng chờ đợi bóng hình tôi xuất hiện để cùng cậu ấy ``bùng nổ''.

``Xin lỗi mọi người\dots\ thực sự xin lỗi\dots''

Lời xin lỗi ấy cứ lặp đi lặp lại trong đầu tôi như một bản kinh cầu nguyện. Tôi luồn lách qua những con phố quen thuộc, vượt qua những ánh đèn đường vừa bắt đầu le lói sáng. Bầu trời Aogami dần chuyển sang màu tím thẫm của buổi hoàng hôn chập choạng.

Trên đường tôi về, những cặp đôi bắt đầu xuất hiện. Họ đi sát vào nhau, tay trong tay, những túi quà nhỏ xinh được trao đi với nụ cười e thẹn. Một chàng trai trẻ vội vàng chỉnh lại chiếc nơ trên hộp sô-cô-la, còn cô gái bên cạnh che miệng cười khúc khích. Ánh mắt họ nhìn nhau đầy ắp thứ ánh sáng mà tôi đang cố gắng trốn chạy.

Tôi cúi đầu, kéo cao cổ áo, bước nhanh hơn. Nhưng càng tránh, những hình ảnh ấy lại càng bám theo. Trước tiệm bánh nhỏ, một cậu học sinh cấp ba run run trao tấm thiệp hồng cho bạn học. Cô bé đỏ mặt, nhận lấy món quà như một báu vật. Họ không cần lời nói, chỉ cần ánh mắt cũng đủ để hiểu nhau.

Còn tôi? Tôi chỉ có tiếng bước chân cô độc vang vọng trên nền đá lạnh. Mỗi lần đi qua một cặp đôi, tôi lại thấy mình như một bóng ma - hiện hữu mà vô hình, sống động mà trống rỗng. Họ trao nhau nụ cười, còn tôi trao cho mình những lời trách móc thầm lặng.

Qua công viên nhỏ, tôi thấy một đôi bạn già ngồi trên băng ghế gỗ. Bà cụ tóc bạc trắng nhẹ nhàng gỡ một chiếc lá rơi khỏi má bác cụ râu xám. Họ đã cùng nhau đi qua bao mùa Valentine, và giờ đây, khi tóc đã phai màu thời gian, họ vẫn còn nhau. Tôi bất giác dừng bước, nhưng rồi lại lắc đầu, tiếp tục đi.

Trên cây cầu bắc qua dòng sông nhỏ, một cậu bé đang run rẩy cầm hộp quà nhỏ. Tôi biết cậu ấy đang chờ đợi điều gì - một lời tỏ tình đầu đời, một khởi đầu mới. Còn tôi? Tôi đang kết thúc, đang chạy trốn khỏi những khởi đầu mà tôi không dám nắm giữ.

Ánh đèn đường bật sáng dần, bầu trời chuyển sang màu chàm sẫm. Những cặp đôi bắt đầu tản về, tay vẫn nắm chặt không rời. Tiếng cười họ vang lên đầy ắp hạnh phúc, xen lẫn trong tiếng gió đêm se lạnh. Tôi đi ngược dòng người ấy, như một kẻ lạc lõng trong chính câu chuyện của mình.

Mỗi bước chân đưa tôi xa hơn khỏi trường học, xa hơn khỏi những trái tim đang chờ đợi, nhưng lại gần hơn với ngôi nhà trống rỗng đang chực chờ. Trong mùa Valentine này, khi cả thế giới dường như đều có ai đó để sánh bước, tôi lại chọn cách đi một mình trên con đường về - không hoa, không chocolate, không nụ cười. Chỉ có những lời xin lỗi vang vọng theo từng nhịp bước, và trái tim nặng trĩu những điều không thể nói thành lời.

Cuối cùng, bóng dáng ngôi nhà của bộ ba Hikawa cũng hiện ra ở cuối con phố. Ánh đèn vàng ấm áp hắt ra từ cửa sổ khiến tôi cảm thấy một chút nhẹ lòng, nhưng nỗi hối hận vì đã bỏ mặc tình cảm của những cô gái ấy vẫn như một bóng ma lởn vởn xung quanh.

Tôi đứng trước cửa nhà, thở dốc, hai tay run rẩy chạm vào tay nắm cửa. Ngày Lễ tình nhân đầu tiên và có lẽ là duy nhất của tôi ở thế giới này, lại kết thúc bằng một cuộc đào thoát không danh giá như thế này sao?

``\vie{Tôi về rồi đây\dots}''

Tôi hít một hơi thật sâu để bình ổn nhịp tim rồi đẩy cửa bước vào. Trái ngược với sự tĩnh lặng của khu vườn bên ngoài, phòng khách đang tỏa ra một không khí ấm cúng và thơm nồng mùi sô-cô-la.

Vừa thấy tôi bước vào, Kaigou đang ngồi trên ghế sofa cùng Suzuki liền đứng dậy. Cả hai nhìn tôi với ánh mắt đầy thấu hiểu, như thể họ đã biết trước mọi chuyện sẽ diễn ra như thế này.

``\vie{Anh về rồi đấy à, Onii-sama? Chào mừng anh đã 'sống sót' trở về,}'' Suzuki mỉm cười tinh nghịch, bước lại gần và giúp tôi cởi chiếc áo khoác nặng nề.

Kaigou khoanh tay trước ngực, nhìn tôi với vẻ mặt điềm tĩnh: ``\vie{Nhìn bộ dạng lếch thếch của cậu kìa. Chắc hẳn là đã có một cuộc đào thoát ngoạn mục lắm phải không?}''

Tôi gượng cười, thả mình xuống ghế: ``\vie{Hai người đừng trêu tôi nữa. Thực sự là\dots\ tôi không biết phải đối diện với họ như thế nào cả.}''

Cô nàng buộc tóc đuôi ngựa thở dài, bước lại gần và đặt lên bàn một hộp quà nhỏ được gói ghém vô cùng tinh tế. ``\vie{Cậu không chọn ai cả, vì cậu biết mình không thuộc về thế giới này. Chúng tôi hiểu tại sao cậu làm thế. Nếu là tôi, có lẽ tôi cũng sẽ chọn cách trốn chạy như một kẻ hèn nhát vậy thôi.}''

Suzuki cũng nhanh nhảu lấy ra một chiếc hộp khác màu hồng nhạt: ``\vie{Đúng thế, vậy nên\dots\ anh đừng quá dằn vặt. Thà là một nỗi đau ngắn ngủi lúc này còn hơn là một sự tổn thương kéo dài mãi mãi. Đây, phần của anh này.}''

Tôi ngẩn người nhìn hai hộp quà trên bàn. Kaigou nhận thấy vẻ ngạc nhiên của tôi, cô ấy hơi quay mặt đi, che giấu một chút bối rối hiếm thấy: ``\vie{Cậu biết đấy, ở thế giới cũ của tôi, nơi mà đàn ông chỉ là những kẻ hạ đẳng và cặn bã, cái ngày Valentine hay ngày Phụ nữ này chỉ là những thứ phù phiếm được tạo ra cho có lệ. Tôi chưa bao giờ nghĩ mình sẽ tặng quà cho bất cứ gã đàn ông nào.}''

Cô ấy dừng lại một chút, giọng nói trở nên trầm hơn và chân thành hơn: ``\vie{Nhưng cậu thì khác, Kuon-kun. Cậu là người đầu tiên khiến tôi cảm thấy rằng việc dành tặng một món quà cho ai đó không phải là điều gì quá tệ hại. Coi như đây là lời cảm ơn vì tất cả những gì cậu đã làm cho hai chị em chúng tôi.}''

Cầm lấy món quà từ hai chị em, tôi cảm thấy lòng mình ấm áp lạ thường. Giữa những bộn bề và tội lỗi của sự trốn chạy, sự thấu hiểu của gia đình chính là liều thuốc an thần quý giá nhất.

``\vie{Cảm ơn hai người\dots\ thực sự cảm ơn,}'' tôi khẽ nói, cảm thấy vị đắng chát trong lòng đã dịu đi đôi chút. Dù biết rằng năm nay Valentine có thể chỉ còn ba người chúng tôi - ba kẻ lữ khách lạc lõng giữa thế giới này - nhưng chỉ cần có họ bên cạnh, tôi đã cảm thấy đủ ấm. Ít nhất, trong ngày mà cả thế giới dường như đều có ai đó để nắm tay, tôi vẫn không hoàn toàn cô đơn. Và có lẽ, chỉ cần thế thôi cũng đủ để tôi tiếp tục bước tiếp.

Kaigou nhìn tôi một lúc rồi đứng dậy, ra hiệu cho tôi về phía phòng khách bên trong: ``\vie{Nhưng cậu biết đấy, dù cậu có trốn tránh thế nào, thì việc cậu nợ họ một lời xin lỗi chính thức là điều không thể tránh khỏi. Và có lẽ\dots\ cậu nên chuẩn bị tinh thần để đối mặt với nó sớm hơn cậu tưởng đấy.}''

Tôi nhìn Kaigou với vẻ khó hiểu: ``\vie{Ý cô là sao? Chẳng lẽ\dots?}''

Cô ấy không đáp, mà chỉ mỉm cười đầy ẩn ý. ``\vie{Thôi nào, Kuon-kun. Mau vào gặp họ thôi!}''

``\vie{Họ?}'' tôi bối rối hỏi.

Cánh cửa phòng khách vừa mở ra, tôi ngỡ như mình vừa bước xuyên qua một cánh cổng không gian. Nhưng không, đây là nhà của tôi, và thực tế trước mắt còn kinh khủng hơn bất cứ con quái vật nào tôi từng đối đầu.

``Cái\dots\ cái gì thế này?!''

Tôi đứng chôn chân tại chỗ, hàm muốn rơi xuống đất. Toàn bộ phòng khách - vốn dĩ chỉ vừa đủ cho ba người chúng tôi - giờ đây chật kín bóng người. Và không phải ai xa lạ, đó chính là toàn bộ 28 cô gái của lớp 1-C. Tất cả họ, từ Shino đang ngồi điềm tĩnh bên cửa sổ, Ayame đang nhai kẹo cao su trên bệ cửa, cho đến Nanase đang đứng khoanh tay với vẻ mặt 'tôi biết ngay mà'\dots\ tất cả đều đang nhìn chằm chằm vào tôi.

Sự bàng hoàng khiến đầu óc tôi quay cuồng. Làm thế nào? Bằng cách nào mà họ lại có thể xuất hiện ở đây cùng một lúc, ngay sau khi tôi vừa thực hiện một cuộc đào thoát quy mô lớn đến vậy?

``\eng{Hahaha! Incroyable ! Một bước đi mà chính kẻ tạo ra trò chơi cũng không ngờ tới, phải không Kuon-user?}'' Game cười sảng khoái trong đầu tôi. ``\eng{Tôi đã bảo là cậu sẽ không thể nào thoát khỏi họ mà, Kuon-user.}''

``Kuon-kun, cậu trốn kỹ thật đấy,'' Ayame là người lên tiếng đầu tiên, nụ cười tinh quái của cô nàng khiến tôi lạnh sống lưng. ``Cậu nghĩ là chỉ cần lẻn ra cửa sau và đi đường vòng là bọn này sẽ chịu thua sao? Cậu quá coi thường 'mạng lưới thông tin' của Hội học sinh bọn tôi rồi.''

\pov{Onidoro Ayame}

``Kuon-kun, cậu thực sự nghĩ rằng chỉ cần đóng tủ đồ lại và lẻn ra cửa sau là mọi chuyện sẽ kết thúc sao? Thật ngây thơ quá đi mất!''

Tôi đứng giữa phòng khách nhà Hikawa, khoanh tay trước ngực với một nụ cười đắc thắng. Nhìn vẻ mặt bàng hoàng của cậu ta lúc này, bao nhiêu công sức tôi bỏ ra trong một giờ qua đều hoàn toàn xứng đáng.

Ngay khi nhận được tin nhắn từ Shino về việc Kuon không xuất hiện tại các điểm hẹn, tôi đã biết ngay là tên ngốc này đang cố chơi trò trốn tìm. Nhưng cậu ta quên mất một điều: tôi là ai chứ? Với tư cách là Hội trưởng Hội học sinh, mạng lưới thông tin của tôi rải rác khắp cái thị trấn Aogami này. Chỉ cần một vài cuộc gọi cho Wata và các bạn học đang 'tản mát' trên đường, tôi đã nhanh chóng xác định được lộ trình 'đào thoát' của cậu ta.

``Nào mọi người, mục tiêu đang tiến về phía nhà Hikawa! Tập hợp quân số ngay lập tức!''

Tôi đã hô hào như thế trên nhóm chat của lớp. Và như đã thấy, các cô gái lớp 1-C khi đã đồng lòng thì còn đáng sợ hơn cả một đội quân tinh nhuệ. Chúng tôi đã có mặt ở đây từ mười lăm phút trước, được sự 'tiếp tay' nhiệt tình của Kaigou-chan và Suzuki-chan để dàn trận chờ sẵn. Một cái bẫy Valentine hoàn hảo mà cậu không tài nào thoát khỏi!

\pov{Hikawa Kuon}

Trước sự chủ động và quyết liệt đến mức 'đáng sợ' của Ayame, tôi chỉ biết đứng hình, cảm thấy mọi kỹ năng sinh tồn của mình đều trở nên vô dụng trước sức mạnh của sự đoàn kết phái đẹp.

Tôi nhìn sang Shino, người đang lặng lẽ quan sát tôi bằng đôi mắt sâu thẳm. ``Shino-san\dots\ cả cậu cũng ở đây sao?''

Cô nàng khẽ gật đầu: ``Mình biết cậu sẽ không chọn bất kỳ ai trong số những bức thư đó, Kuon-kun. Vì cậu lo sợ về sự chia ly, phải không?''

Câu hỏi của cô ấy như một đòn giáng mạnh vào tâm trí tôi. Hóa ra, họ đã thấu hiểu nỗi lòng của tôi từ lâu. Shino, với linh cảm nhạy bén và sự thấu hiểu về bản chất thực sự của chúng tôi, đã âm thầm kết nối với mọi người. Khi thấy tôi không xuất hiện tại các điểm hẹn, Ayame và Wata đã nhanh chóng tập hợp cả lớp lại. Với họ, việc tìm ra địa chỉ nhà Hikawa (dù gì họ cũng đã không còn lạ lẫm gì với chỗ này rồi) và 'đột kích' nơi này là chuyện trong lòng bàn tay.

``Cậu định cứ thế mà biến mất sao?'' Koishin bước lên phía trước, đôi mắt cô nàng đỏ hoe nhưng giọng nói vẫn cố tỏ ra đanh thép. ``Tặng sô-cô-la không phải là để trói buộc cậu, đồ ngốc! Bọn này chỉ muốn\dots\ chỉ muốn cậu biết là chúng tôi rất trân trọng cậu thôi!''

Phòng khách bỗng chốc trở nên náo nhiệt. Tiếng trách móc của Koishin, tiếng cười hì hì của Ayame, sự lo lắng của Wata, và cả những ánh mắt dịu dàng từ Mitsuki, Miku hay Touka\dots\ tất cả hòa quyện vào nhau tạo nên một bầu không khí hỗn loạn nhưng vô cùng ấm áp.

Tôi đứng đó, giữa tâm bão của những tình cảm chân thành, cảm thấy mọi rào cản và sự trốn chạy của mình trở nên thật nực cười. Hóa ra, họ không cần một 'lời chọn lựa' duy nhất. Họ chỉ cần tôi đối diện với họ.

Nhìn Suzuki và Kaigou đang đứng phía sau mỉm cười, tôi hiểu rằng họ cũng đã tham gia vào màn kịch này. Một màn kịch để kéo tôi ra khỏi cái kén của sự sợ hãi.

Tôi hít một hơi thật sâu, gỡ chiếc mũ áo xuống và nhìn thẳng vào hai mươi tám khuôn mặt đang chờ đợi. Có lẽ, thay vì trốn chạy khỏi nỗi đau của tương lai, tôi nên học cách trân trọng hiện tại này - ngay tại ngôi nhà này, cùng với những con người tuyệt vời này.

``Tôi\dots\ thực sự xin lỗi vì đã trốn chạy,'' tôi nói, giọng nói đã lấy lại sự bình tĩnh nhưng chứa chan cảm xúc. ``Và\dots\ cảm ơn mọi người đã đến đây.''

\ct

Bữa tiệc Valentine tại gia đình Hikawa chúng tôi bùng nổ theo một cách mà tôi chưa từng tưởng tượng nổi. Không còn sự trốn chạy, không còn những lời xin lỗi thầm lặng, chỉ còn tiếng cười nói vang vọng khắp ngôi nhà vốn thường ngày chỉ có ba người chúng tôi, và những loại sô-cô-la mà các cô gái đã cất công làm trong buổi sáng nay. Không khác nào một buổi tụ tập dã ngoại tại gia, với cái cớ là mừng ngày Lễ tình nhân.

``Nè Kuon-kun! Cầm lấy micro đi chứ!'' Yuki phấn khích vẫy tay, tay kia vẫn đang lăm lăm chiếc điện thoại để phát nhạc. ``Ở quán karaoke cậu  có thể trốn được, chứ ở đây thì đừng hòng! Chúng ta sẽ biến phòng khách này thành sân khấu 'Super Nova' thứ hai!''

Tôi chưa kịp phản ứng thì Saya đã lững thững tiến lại gần, đưa cho tôi một đĩa nhỏ đầy những viên sô-cô-la hình ngôi sao. ``Cậu nên ăn chút gì đi, trông cậu mệt mỏi như vừa chạy marathon xuyên thị trấn ấy\~{} Sô-cô-la này sẽ giúp cậu nạp lại năng lượng cho 'vận may' của mình đó\~{}''

Ở góc phòng, Honoka đang ồn ào khoe về hộp sô-cô-la siêu ngọt của mình với bất cứ ai đi ngang qua. ``Tôi đã bảo rồi mà! Sô-cô-la càng ngọt thì tình bạn càng bền lâu! Đúng không Kuon-kun?''

``Chậc, sai bét!'' Nanase cất tiếng phản bác, vừa bước tới vừa nhấp một ngụm cacao đắng. ``Sô-cô-la \emph{đắng} mới là thước đo chân thành - càng đắng, càng thể hiện nỗ lực tìm hiểu khẩu vị đối phương. Ngọt sẵn thì ai chẳng làm được?''

Honoka chu mỏ, ôm sát hộp kẹo: ``Ngọt thì dễ gần, dễ cưa! Đắng quá ai dám ăn?''

``Chính vì phải \emph{vượt qua} vị đắng, tình cảm mới sâu,'' Nanase đẩy kính, nhìn sang tôi. ``Kuon-san, cậu nói đi, cậu chọn phe nào?''

Bị kéo vào giữa, tôi giật mình, hai hộp quà trên tay suýt rơi. ``Ơ, tôi\dots\ thực ra\dots'' Tôi lùi một bước, đưa mắt tìm cứu viện. ``Shiina, cậu thích loại sô-cô-la nào?''

Shiina đang xếp ghế, nghe thấy tên mình thì quay lại, đáp gọn: ``Thế nào cũng được. Cái đó tùy từng khẩu vị thôi.''

Nanase và Honoka nhìn nhau, cùng ``Hử?'' một tiếng, còn tôi thì thầm thở phào, vừa thoát kiếp làm \emph{đồng xu tùy} mùi vị, vừa học được bài học: đừng bao giờ hỏi đội trưởng CLB thể thao về… chuyện nhạy cảm khẩu vị.

Cô nàng bờm Sư tử thì vẫn giữ phong cách khỏe khoắn, cô nàng đang cùng Kaoru sắp xếp lại bàn ghế để lấy chỗ cho mọi người ngồi. ``Kuon-kun, lần sau có tập luyện thì cứ nói một tiếng, đừng có lẻn ra cửa sau như thế. Nhìn cậu lúc nãy tôi cứ tưởng là kẻ trộm cơ đấy!'' Kaoru cười khúc khích phụ họa theo, ánh mắt không rời khỏi vẻ bối rối của tôi.

Cặp đôi Akane và Yuka thì đang cùng Suzuki bày biện thức ăn trong bếp. Tiếng cười của họ thỉnh thoảng lại vang lên khi Yuka trêu chọc Akane về việc cô nàng đã lo lắng thế nào khi Kuon không xuất hiện dưới tán cây lớn. Inari thì vẫn còn thở dốc, có lẽ vì vừa phải chạy đua với thời gian để đến đây kịp lúc: ``Hên là\dots\ hên là nhà cậu không quá xa trường\dots''

Miku đứng cạnh Mitsuki, cả hai trông như những tiểu thư đang tham dự một buổi dạ tiệc thượng lưu hơn là một buổi tụ tập học sinh. ``Tôi đã nói là cậu ấy sẽ về nhà mà,'' Miku khẽ mỉm cười, vẻ trưởng thành thường ngày pha chút dịu dàng. Mitsuki gật đầu tán thành, tay khẽ chỉnh lại nơ áo: ``Một kết thúc có hậu cho cuộc truy đuổi của chúng ta đấy na, Kuon-kun.''

Nhìn quanh phòng, tôi thấy Azuki đang say sưa thảo luận về một bản nhạc mới với Mari. Shion thì vẫn giữ vẻ 'ngầu' của mình, ngồi góc phòng đọc một cuốn sách về ma thuật nhưng thỉnh thoảng lại liếc nhìn về phía đám đông. Sakura và Shino đang cùng Kaigou trò chuyện, có vẻ như ba người họ đã tìm thấy tiếng nói chung trong việc ``giám sát'' tôi. Những người còn lại như Saki, Suzune, Touka\dots\ đều đang tận hưởng không khí náo nhiệt này theo cách riêng của mình.

Suzuki chạy đến, kéo tay tôi lắc lắc. ``\vie{Onii-sama, anh có thấy không? Mọi người thực sự rất có cảm tình với anh đó!}''

Kaigou bước đến bên tôi, đưa cho tôi một cốc nước cam. ``\vie{Thế nào? Một sự trừng phạt ngọt ngào chứ?}''

Tôi nhìn ly nước, rồi nhìn vào căn phòng đầy ắp tiếng cười. ``\vie{Có lẽ là vậy. Cảm ơn cô nhé, Kaigou-chan, Suzuki.}''

Một đêm Valentine không giống bất cứ điều gì tôi từng trải qua ở thế giới thực hay bất cứ vũ trụ nào khác. Một đêm mà ranh giới giữa người ngoại tộc và cư dân thị trấn Aogami dường như đã bị xóa nhòa bởi mùi hương nồng nàn của sô-cô-la và sự ấm áp của tình bạn.

\ct

\begin{table}[H]
  \centering
  \setlength{\arrayrulewidth}{1pt}
  \begin{tabular}{|>{\centering\arraybackslash}p{0.4\linewidth}|>{\centering\arraybackslash}p{0.4\linewidth}|}
    \hline
    \hyperref[choice]{Quay trở về lựa chọn} & \hyperref[ending]{Đi tới phần cuối} \\ \hline
  \end{tabular}
\end{table}


\newpage
\label{ending}

Lễ tình nhân ở thị trấn Aogami đã khép lại như thế đó.

Dù chúng tôi chọn bước đi cùng ai trên những con phố rực rỡ ánh đèn, hay có lúc lạc hướng giữa những cảm xúc chưa kịp gọi tên để rồi nhận ra mình vẫn luôn được bao bọc bởi tình yêu thương, thì vị đắng và ngọt của sô-cô-la ngày hôm nay vẫn sẽ trở thành một mảnh ký ức không thể phai mờ.

Đối với chúng tôi - Hikawa Kuon, Kaigou, Suzuki, cùng Misora Shino và Shinomiya Sakura - những ngày cuối xuân này không chỉ là về những hộp quà hay những lời tỏ tình e thẹn. Nó là lời nhắc nhở rằng, dù gánh nặng phía trước có lớn đến đâu, dù con đường chúng tôi đang đi vẫn còn phủ đầy những điều chưa biết, thì những rung động chân thành nhất vẫn luôn tồn tại. Chính những điều nhỏ bé đó đã giữ chúng tôi đứng vững, giữa một thị trấn mà quá khứ và hiện tại vẫn còn chồng lấp lên nhau.

Dư vị của buổi chiều tà ven sông, tiếng cười vang lên trong phòng karaoke, hay sự tĩnh lặng nơi đền thờ cổ\dots\ tất cả hòa quyện thành một giai điệu dịu dàng của tuổi trẻ: dù mong manh, nhưng đủ sức níu giữ một phần con người chúng tôi lại, không để bị cuốn trôi hoàn toàn bởi những gì đang chờ phía trước.

Valentine tại Aogami\dots\ có lẽ không chỉ là một lễ hội. Đó là một khoảng lặng. Một điểm dừng chân. Một lời nhắc rằng chúng tôi vẫn còn là những con người bình thường, trước khi phải quay lại đối mặt với thứ đã tồn tại suốt hơn một nghìn năm.

Nhưng chúng tôi đều hiểu rõ: khoảng lặng ấm áp này chỉ là tạm thời.

Ngoài kia, rừng Aogami vẫn âm thầm mở mắt. Những tán cây đen kịt như đang nhìn chúng tôi, chờ đợi một lời hứa còn dang dở. Mỗi lần gió thổi qua tán lá, tôi lại nghe thấy tiếng thì thầm của quá khứ, những ký ức chưa kịp khép lại, những linh hồn vẫn lẩn khuất trong bóng tối.

Lần tới quay lại đây, chúng tôi không còn là khách qua đường nữa. Không còn được phép lần mò, thắc mắc, ghi chép. Chúng tôi phải bước vào để kết thúc cho cái vòng lặp đã kéo dài hơn cả đời người. Một cái kết thật sự, không còn kéo dài thêm một khoảng nữa. Một nghìn ba trăm năm tìm kiếm sự công nhận của những hồn ma đó\dots\ cần phải kết thúc.

Sáng mai, mặt trời vẫn lên như thường lệ, nhuộm vàng những tán lá sồi già. Nhưng ánh nắng ấy sẽ không chỉ chiếu vào kẻ chứng kiến. Nó sẽ soi thẳng vào chúng tôi, những người cầm bút vẽ dấu chấm hết cho lời nguyền, cho những hồn ma bị bỏ lại, cho tất cả những gì đáng lẽ đã phải chìm vào quên lãng từ ngàn năm trước.

Đêm tháng Năm tràn xuống, trăng bạc vẽ những mảng sáng loang lổ trên mặt đất. Chúng tôi đứng thành hàng, im lặng, để hơi ấm từ bữa tiệc Valentine còn vương trong lồng ngực.

Chúng tôi biết: bước tiếp theo có thể khiến mọi thứ đổi màu mãi mãi. Nhưng đã đến lúc.

Và lần này, chúng tôi sẽ không chọn cách trốn chạy.

\end{document}